% concatenated from combined.tex
% GENERATED by make_combined.py. Do not edit; edit main.tex or
% supplementary.tex and regenerate.

\documentclass[11pt]{article}
\usepackage[margin=1in]{geometry}
\usepackage{amsmath,amssymb,amsthm}
\usepackage{graphicx}
\usepackage{booktabs}
\usepackage{array}
\usepackage{authblk}
\usepackage{lineno}
\usepackage{natbib}
\usepackage{hyperref}

% Figure files are placed in a figures/ subfolder
\graphicspath{{figures/}}

\newtheorem{theorem}{Theorem}
\newtheorem{lemma}[theorem]{Lemma}
\newtheorem{corollary}[theorem]{Corollary}
\newtheorem{definition}[theorem]{Definition}

\newcommand{\starG}{\star_G}
\newcommand{\starM}{\star_M}
\newcommand{\calA}{\mathcal{A}}
\newcommand{\calB}{\mathcal{B}}
\newcommand{\calC}{\mathcal{C}}
\newcommand{\calI}{\mathcal{I}}
\newcommand{\calS}{\mathcal{S}}
\newcommand{\calT}{\mathcal{T}}
\newcommand{\calU}{\mathcal{U}}
\newcommand{\calV}{\mathcal{V}}
\newcommand{\calX}{\mathcal{X}}
\newcommand{\bbR}{\mathbb{R}}
\newcommand{\bbZ}{\mathbb{Z}}
\newcommand{\Rsq}{R\textsuperscript{2}}

\hypersetup{colorlinks=true,linkcolor=blue,citecolor=blue,urlcolor=blue}

% ---- packages, macros and theorem environments only the SI uses ----
\usepackage{algorithm}
\usepackage{algpseudocode}
\newtheorem{proposition}[theorem]{Proposition}
\newtheorem{remark}[theorem]{Remark}
\newcommand{\starSig}{\star_{\Sigma}}
\newcommand{\calZ}{\mathcal{Z}}
\newcommand{\bbC}{\mathbb{C}}
\newcommand{\bbH}{\mathbb{H}}
\newcommand{\bbK}{\mathbb{K}}

\begin{document}


%\linenumbers   % required for peer-review submission

% Science Advances title candidates (capability-led; pick one before submission):
%   - "Exact symmetry as algebra: a machine-verified tensor calculus that enforces
%      physical selection rules"  [current working title]
%   - "Making symmetry exact: a machine-verified tensor calculus for equivariant
%      computation and physical selection rules"
%   - "Symmetry without approximation: an exact, machine-verified tensor calculus
%      of physical selection rules"
% Older framework-led candidates (kept for reference):
%   - "Making symmetry arithmetic: a group-algebraic tensor calculus"
%   - "A group-algebraic tensor calculus turns symmetry into computable structure"
%   - "Symmetry as algebra: equivariance and physical selection rules from a group-Fourier tensor calculus"
\title{\textbf{Exact Symmetry as Algebra: A Machine-Verified Tensor Calculus that Enforces Physical Selection Rules}}

\author[4]{Paulina Hoyos}
\author[1]{Shashanka Ubaru}
\author[5]{Dongsung Huh}
\author[1]{Vasileios Kalantzis}
\author[1]{{Kenneth L. Clarkson}}
\author[3]{Misha Kilmer}
\author[2]{Haim Avron}
\author[1]{Lior Horesh}

\affil[1]{IBM Research}
\affil[2]{Tel-Aviv University}
\affil[3]{Tufts University}
\affil[4]{UT Austin}
\affil[5]{Independent}

\date{}
\maketitle

% === Science Advances one-sentence summary (teaser; submit in the portal, ~125 chars): ===
% Approximate symmetry compounds away with depth; made the multiplication rule
% of an algebra, symmetry is exact and provable.
% (124 characters, within the ~125-char portal limit.)
% Alternate (119 chars): Make a symmetry the multiplication rule of a tensor
% algebra, and equivariance, compression, and selection rules follow.
%
% Science Advances uses NO Significance Statement; the block below is retained
% (disabled with \iffalse) only for a possible PNAS fallback.
\iffalse
\noindent\textbf{Significance Statement.} Machine learning usually captures a
physical symmetry by hand-building a specialized network for it. We show the symmetry
can instead become arithmetic: we define a tensor algebra whose multiplication rule is
a chosen group, so that equivariance, a machine-verified linear-algebra toolbox, and
interpretability follow from the algebra rather than from network design, and several
symmetries compose by simply multiplying their groups. Because the algebra resolves any
measurement into the channels of angular momentum, it acts as a spectroscope for
symmetry: with no quantum-mechanical input it surfaces angular-momentum signatures
consistent with the 1931 Wigner--Eckart theorem in molecular data, and on real
inorganic crystals it enforces the point-group selection rules of the elastic tensor
\emph{exactly}, where a trained graph network instead leaks energy into physically
forbidden channels. Made algebraic, symmetry becomes something a computer can both
exploit and read directly off measurement.
\fi

\begin{abstract}
Symmetry is central to the physical sciences, yet machine learning usually
captures it only approximately, leaving a residual per-step equivariance
error $\varepsilon$ that compounds with depth $M$ as $M\varepsilon$, whereas
exact equivariance holds at unbounded
depth; we demonstrate this divergence at fourteen orders of magnitude. We
show that a symmetry can be made exact by construction, as the multiplication
rule of a tensor algebra. In the resulting $\starG$ algebra, defined by any
finite group $G$, the group-Fourier transform block-diagonalizes every tensor
into irreducible-representation blocks, making equivariance intrinsic;
requiring equivariance conversely \emph{forces} this suitably normalized
transform, so the algebra is determined by $G$ rather than chosen. The
standard matrix toolbox, including a Frobenius-optimal low-rank
factorization, transfers blockwise, machine-checked in Lean~4 under an
explicit axiom budget, and extends unchanged to band-limited compact groups
and, under periodic boundary conditions, to all 230 crystallographic space
groups and the compact
little-group fibers of Euclidean and Poincar\'e symmetry. This exactness
is an applied capability: on inorganic-crystal elastic tensors the algebra
enforces point-group selection rules exactly on the output of \emph{any}
predictor, driving a trained graph network's forbidden-channel leakage from
$10^{-2}$ to machine zero, eliminating mechanically unstable predictions, and
recovering viable materials that an unconstrained screen discards;
on molecular data, with no quantum-mechanical input, it exposes octahedral
selection-rule signatures consistent with the Wigner--Eckart theorem. Matched
networks lead on pooled molecular accuracy, which we report plainly: the
contribution is a complementary algebraic calculus, structural and
diagnostic, exact at any depth.
\end{abstract}

\section{Introduction}

Symmetry enters computation in two ways that are usually treated as one:
exactly, as structure the computation cannot violate, or approximately, as
structure learned from data, imposed by a penalty, or left behind as a
residual error. This paper's starting point is that the difference is one of
kind, not degree. Symmetric computations are rarely applied once, and
iteration compounds any approximate symmetry: a single step whose
equivariance error is $\varepsilon$ can accumulate $O(1)$ symmetry violation
by depth $M \sim 1/\varepsilon$, whereas an exactly symmetric computation is
symmetric at every depth (Section~\ref{sec:compounding}). Our thesis is that
exactness is available by construction: a symmetry can be made the
multiplication rule of a tensor algebra, and every guarantee in this paper,
from machine-checked optimal factorizations to the exact enforcement of
physical selection rules on any predictor's output, flows from that single
move.

The obstacle this removes is familiar. Much of the data encountered in science and engineering is inherently
multidimensional: molecular configurations encode three-dimensional atomic
positions, quantum states exist in exponentially large Hilbert spaces, and sensor
arrays sample signals across spatial and temporal domains. Traditional machine
learning methods typically vectorize such data, collapsing its natural tensor
structure into flat feature vectors~\citep{kolda2009tensor,sidiropoulos2017tensor}.
This is akin to unfolding an origami crane into a flat sheet: the operation is
technically lossless, but the geometry that gives the object its meaning is
destroyed. All subsequent processing must then recover, implicitly and at great
computational cost, the structure that was discarded at the outset.

Symmetry sharpens this problem. A molecule exists in three-dimensional space:
it can be rotated without changing its properties (rotational symmetry), and its
identical atoms can be indexed in any order without changing the molecule itself
(permutation symmetry). These symmetries coexist and
interact~\citep{noether1918,bronstein2021geometric}, yet vectorization treats them
as incidental features of the data rather than as fundamental structural
constraints.

The dominant paradigm for incorporating symmetry is through Equivariant Neural
Networks
(ENNs)~\citep{cohen2016group,thomas2018tensor,fuchs2020se3,batzner2022e3}, which
have achieved remarkable success in molecular property
prediction~\citep{schutt2017schnet} and protein
structure~\citep{jumper2021alphafold}. ENNs handle symmetry through architecture:
to respect a rotation, one engineers rotation-equivariant layers; to respect a
permutation, one engineers permutation-equivariant layers. But when multiple
symmetries coexist, as they do in virtually all physical systems, the
architectural approach faces a combinatorial wall. The blueprint must be
redesigned from scratch for each new combination of symmetries, with no guarantee
that the resulting representation is optimal in any rigorous sense. The physics
is hard-coded into the network topology, and changing the physics means rebuilding
the network. And wherever the symmetry is achieved only approximately, the
compounding cost described above is incurred on top of the architectural one.

In this article, we propose a different philosophy. Instead of constraining the
architecture to fit the symmetry, we change the mathematics to fit the geometry
of the data. Building on the $\starM$ algebra of Kilmer and
collaborators~\citep{kilmer2021tensor,kernfeld2015tensor}, we construct a tensor algebra,
$\starG$, where any finite group $G$ defines the multiplication rule. The
resulting algebra inherits equivariance as an intrinsic property of
multiplication, not as an architectural constraint. Composing multiple symmetries
requires only specifying the direct product $G_1 \times G_2 \times \cdots \times
G_d$; no redesign is needed. Because the group-Fourier transform
block-diagonalizes every $\starG$ tensor, the algebra is fully
\emph{matrix-mimetic}: the standard toolbox of matrix linear algebra, the SVD,
QR factorization, symmetric eigendecomposition, least-squares, and functional
calculus, transfers to $\starG$ tensors blockwise, with every operation exactly
equivariant by construction. In particular the $\starG$-SVD carries a provable
Eckart--Young optimality (Theorem~\ref{thm:eckart_young}), and, by the
Peter--Weyl theorem~\citep{serre1977linear,peterweil1927}, the same
decomposition into irreducible-representation channels reveals the symmetry
content of physical observables (Section~\ref{sec:wigner_eckart}). These
operations are machine-checked in Lean~4, conditional on standard finite-group
Fourier and matrix-approximation facts declared as an explicit axiom budget rather
than reproved from foundations (SI~Section~\ref{sec:lean}). The Wigner--Eckart
theorem (1931) states that matrix elements of tensor operators between angular
momentum eigenstates factorize into a geometric part (Clebsch--Gordan
coefficient) and a reduced matrix element independent of magnetic quantum numbers,
implying selection rules: an operator of rank $l$ couples only states whose
angular momenta differ by at most $l$. We show empirically that signatures
consistent with these rules are recoverable from data alone via $\starG$
decomposition.

\emph{Inherited foundations.} The transform-domain reading of low-rank
optimality is inherited: for a fixed invertible transform $M$, the $\starM$/t-SVD
Eckart--Young theorem is due to Kilmer and
collaborators~\citep{kilmer2021tensor,kernfeld2015tensor}, and its tubal
(commutative) specialization has been studied by Mor and
Avron~\citep{mor2025quasitubaltensoralgebraseparable,mor2026sufficientnecessaryconditionseckartyoung}.
Three architecture-side ingredients are likewise inherited: the equivalence between
equivariance to a compact group and group convolution (established for general
groups by Kondor and Trivedi~\citep{kondor2018generalization} and by Cohen, Geiger,
and Weiler~\citep{cohen2019general}); computing directly in the non-abelian
group-Fourier domain, where activations are per-irrep matrices (Clebsch--Gordan
networks on $\mathrm{SO}(3)$~\citep{kondor2018clebsch}); and composing direct-product
symmetries without architectural redesign (steerable-CNN
programs~\citep{cesa2022program}). Each of these builds \emph{layers}.

\emph{What this paper adds.} We build the \emph{algebra} underneath those
layers: an associative
$\starG$ product under which those same transforms become a matrix-mimetic calculus
with its own factorizations, and composition acts on the algebra itself. Moreover
the transform is not a modeling choice: requiring the product to be equivariant
\emph{forces} it to be the group-Fourier transform $F_G$, unique up to a per-irrep
basis change and a relabeling of equal-dimension blocks
(SI~Theorem~\ref{thm:necessity}); the group determines the algebra, not the reverse.
Concretely (Figure~\ref{fig:framework}), what is new is (a)~the exactness
dividing line made quantitative: approximate equivariance
can compound to $O(1)$ symmetry violation under iteration, demonstrated at
every single-step error real approximation strategies achieve, while exact
equivariance holds at unbounded depth (Section~\ref{sec:compounding});
(b)~the specialization to the \emph{group}-Fourier
transform, which makes equivariance intrinsic and gives each block a physical
(angular-momentum) meaning; (c)~the extension to \emph{non-abelian} groups within
this tensor-algebra setting, where transform-domain blocks are matrices rather than
scalar tubes; (d)~the full
matrix-mimetic toolbox (SVD with Eckart--Young optimality, QR, symmetric
eigendecomposition, least-squares, functional calculus), with the algebra laws and
the SVD, QR, eigendecomposition, and least-squares reductions formalized in Lean~4;
(e)~product-group composability at the level of the algebra, by
$F_{G_1}\otimes F_{G_2}$; and (f)~the
per-irrep symmetry diagnostics on molecular and crystalline data, including an
exact selection-rule enforcement layer for any predictor's tensor output.

\begin{figure}[!t]
\centering
\includegraphics[width=\textwidth]{Figure1.pdf}
\caption{\textbf{The $\starG$ tensor algebra: symmetry as the multiplication
rule of the data.}
(\textbf{Top left, From molecules to algebra}) Molecular data measured under all
elements of a symmetry group $G$ form a structured tensor
$\calA \in \bbR^{n \times d \times |G|}$, preserving geometric information that
is destroyed by vectorization into $A \in \bbR^{n \cdot d \times |G|}$.
(\textbf{Top right, The $\starG$ product}) Two tensors are multiplied via group
convolution along the tube dimension, computed efficiently in the Fourier domain
via the Peter--Weyl theorem: $F_G$ transforms each tensor to its block-diagonal
spectral form, standard matrix products are applied per irreducible representation
block, and $F_G^{-1}$ returns the result. The group $G$ can be any finite group
(a single symmetry or a product $G_1 \times G_2 \times \cdots$) with no
architectural changes required.
(\textbf{Bottom left, The $\starG$-SVD}) Every $\starG$-tensor admits a
factorization $\calA = \calU \starG \calS \starG \calV^H$. The rank-$k$
truncation $\calA_k$ is provably optimal:
$\|\calA - \calA_k\|_F \leq \|\calA - \calB\|_F$ for any rank-$k$ equivariant
tensor $\calB$ (Eckart--Young theorem for $\starG$, Theorem~\ref{thm:eckart_young}).
(\textbf{Bottom right, From Eckart--Young to Wigner--Eckart}) The optimal
truncation and the physical selection rules are two readings of the same
irreducible-representation decomposition, so the algebra doubles as a
spectroscope for physical symmetry. Decomposing predictive power by irreducible representation
(irrep) over the octahedral group reads off octahedral selection-rule signatures
consistent
with the Wigner--Eckart theorem directly from molecular geometry data: scalar properties are dominated by the
$l\!=\!0$ (A$_1$) channel, the dipole magnitude draws on the $l\!=\!1$
(T$_1$) channel, and polarizability is uniquely insensitive to $l\!=\!1$.
Every operation in this pipeline is exactly equivariant, the property that
Sections~\ref{sec:compounding} and~\ref{sec:crystals} convert into
depth-stability and enforcement guarantees.}
\label{fig:framework}
\end{figure}

\section{Results}

\subsection{The \texorpdfstring{$\starG$}{starG} Algebra}
\label{sec:theory}

\subsubsection*{Theoretical Foundation}

Let $G$ be a finite group of order $n$ with elements $\{g_1, g_2, \ldots, g_n\}$. We define the \emph{convolution tensor} $\calT \in \bbR^{n \times n \times n}$ by
\begin{equation}
\calT(a, b, c) = \begin{cases} 1 & \text{if } ab = c \\ 0 & \text{otherwise}
\end{cases}
\end{equation}
for all $a, b, c \in G$. This tensor encodes the complete multiplication table of $G$: its frontal slices are permutation matrices, and it satisfies an associativity identity inherited directly from group associativity. By the Peter--Weyl theorem~\citep{serre1977linear,peterweil1927}, $\calT$ admits a spectral decomposition:
\begin{equation}\label{eq:spectral}
\calT(a, b, c) = \sum_{i,j,k} \calC(i, j, k) \, F_G(a, i) \, F_G(b, j) \, F_G^{-1}(c, k),
\end{equation}
where $F_G$ is a generalized Group Fourier transform matrix assembled from the
irreducible unitary representations (irreps) of~$G$, and $\calC$ is a sparse core
tensor encoding the block-diagonal matrix multiplication structure of the irreps
of~$G$. For abelian groups, $F_G$ is a generalized Fourier matrix (for cyclic groups, it reduces to the standard DFT matrix) and $\calC$ is
diagonal; for non-abelian groups, $F_G$ is an invertible matrix. The precise
definition of $F_G$ is provided in the Supplementary Information (SI Section~2).
Crucially, equation (\ref{eq:spectral}) means that group convolution in the
original domain corresponds to \emph{block-diagonal matrix multiplication} in the
Fourier domain: one independent matrix product per irrep block, enabling highly
efficient computation.

We view order-$(2+d)$ tensors $\calA \in \bbR^{\ell \times m \times n_1 \times
\cdots \times n_d}$ as $\ell \times m$ matrices whose entries (which are
$d$-order tensors) lie in the convolutional ring $\mathbb{K}_G$ for
$G = G_1 \times \cdots \times G_d$ and $n_i = |G_i|$ (SI Definition~\ref{def:conv_ring}).
The \emph{$\starG$ product}  of $\cal{A}$ with $\calB \in \bbR^{m \times p \times n_1 \times
\cdots \times n_d}$
is defined via group convolution along the group dimensions:
\begin{equation}
(\calA \starG \calB)_{ij}(c_1, \ldots, c_d) = \sum_{k} \sum_{(a_1, \ldots,
a_d) \in G} \calA_{ik}(a_1, \ldots, a_d) \, \calB_{kj}(a_1^{-1}c_1, \ldots,
a_d^{-1}c_d).
\end{equation}
This product defines a tensor algebra, specializing the $\starM$ construction of
Kernfeld, Kilmer, and Aeron~\citep{kernfeld2015tensor} to the group-Fourier
transform, that generalizes classical matrix algebra while embedding the symmetry
group $G$ directly into the multiplicative structure. The resulting algebraic system
supports a full suite of matrix-mimetic operations (inverses, transposes, norms,
and decompositions), all inheriting equivariance by construction. Equivariance
is thus a property of the algebra, not an imposed constraint. The $\starG$
product can be computed efficiently by (i)~applying $F_G$ to each tensor along
its group dimension, (ii)~performing standard matrix products at each of the
$|\hat{G}|$ Fourier irreps independently in parallel, and
(iii)~applying $F_G^{-1}$ to recover the result. For Abelian groups $G$, the total cost is
$O(n \ell m p + n \log n)$ including the Fourier transforms, matching the
complexity of a single matrix product up to logarithmic factors.

\paragraph{One group or many: the flattening convention.} The multi-group and
single-group views are interchangeable, and we use both. The bijection
$a \leftrightarrow (a_1, \ldots, a_d)$ identifies $G = G_1 \times \cdots
\times G_d$ with a single finite group of order $n = n_1 \cdots n_d$, whose
Fourier matrix is $F_G = F_{G_1} \otimes \cdots \otimes F_{G_d}$ and whose
irreducible representations are the tensor products of the factors' irreps
(Theorem~\ref{thm:product_ring} below). Under this identification an
order-$(2+d)$ tensor in $\bbR^{\ell \times m \times n_1 \times \cdots \times
n_d}$ \emph{is} an order-3 tensor in $\bbR^{\ell \times m \times n}$, so every
definition, factorization, and algorithm in this paper, when stated in the
order-3 single-group form for readability, applies verbatim to product groups
and higher-order tensors. We also retain the t-product term \emph{tube} for
the group-indexed entry $\calA_{ij} = \calA(i,j,:) \in \mathbb{K}_G$
throughout; for $d > 1$ a tube is a $d$-dimensional array rather than a
vector.

The \emph{$\starG$-Hermitian transpose} $\calA^H \in \bbR^{m \times \ell \times
n_1 \times \cdots \times n_d}$ is defined entry-wise by
\begin{equation}\label{eq:transpose}
(\calA^H)_{ij}(g_1,\ldots,g_d) = \overline{\calA_{ji}(g_1^{-1},\ldots,g_d^{-1})},
\end{equation}
where the overline denotes complex conjugation (trivial for real-valued tensors).
Equivalently, in the Fourier domain,
$\widehat{\calA^H}(:,:,\rho) = \hat{\calA}(:,:,\rho)^H$ for every irrep $\rho$
(see SI Section~3 for the precise definition of $\hat{\calA}(:,:,\rho)$),
so the $\starG$-transpose maps to the ordinary matrix Hermitian transpose at each
irrep block. This definition makes $\starG$-unitarity and the SVD factor
conditions below fully analogous to their matrix counterparts.

\subsubsection*{The $\starG$-SVD and Optimality}

Every tensor $\calA$ in the $\starG$ algebra admits a singular value
decomposition $\calA = \calU \starG \calS \starG \calV^H$, where $\calU$ and
$\calV$ are $\starG$-unitary (satisfying $\calU^H \starG \calU = \calI$) and
$\calS$ is f-diagonal (its frontal slices are diagonal matrices) with
non-negative real entries, the \emph{singular tubes} $\mathbf{s}_i$ (elements
of $\mathbb{K}_G$; $d$-dimensional arrays when $G$ is a product of $d$ groups,
per the flattening convention above). This
decomposition is computed exactly by applying the group Fourier transform along
the group dimension, performing standard matrix SVDs independently at each
Fourier irrep, and applying the inverse group Fourier transform, a procedure that
is both exact and computationally efficient.

Two rank notions are in play, and we state the optimality for both. Write the
group-Fourier blocks of $\calA$ as $\hat\calA(\rho)$ with singular values
$\sigma_1(\rho)\ge\sigma_2(\rho)\ge\cdots$ for $\rho\in\hat G$, and Plancherel weight
$d_\rho/|G|$.

\begin{theorem}[Eckart--Young for $\starG$]\label{thm:eckart_young}
\emph{(i) Per-block form (the machine-verified statement).} For any prescribed
per-block rank profile $(k_\rho)_{\rho\in\hat G}$, the blockwise SVD truncation that
keeps the top $k_\rho$ singular values of each $\hat\calA(\rho)$ minimizes
$\|\calA-\calB\|_F^2$ over all tensors $\calB$ whose block at $\rho$ has matrix rank
at most $k_\rho$, with error $\sum_{\rho}\frac{d_\rho}{|G|}\sum_{j>k_\rho}\sigma_j(\rho)^2$.
\emph{(ii) Tubal form.} Ordering the singular tubes $\mathbf{s}_i$ by Frobenius norm,
the truncation $\calA_k$ keeping the leading $k$ tubes minimizes $\|\calA-\calB\|_F^2$
over all $\calB$ of $\starG$-rank at most $k$, with
$\|\calA-\calA_k\|_F^2=\sum_{i>k}\|\mathbf{s}_i\|_F^2$; this is the instance of~(i) at
the coupled profile $k_\rho=k\,d_\rho$.
\end{theorem}

\noindent The $\starG$-rank of $\calA$ is the number of non-zero singular tubes in
the $\starG$-SVD $\calA=\calU\starG\calS\starG\calV^H$. The two competitor classes,
and hence the two rank notions, coincide when every irrep is one-dimensional (abelian
$G$) and differ otherwise. Full proofs of both forms are in the Supplementary
Information (SI~Section~5), and the per-block form~(i), together with the underlying
matrix Eckart--Young theorem, is machine-checked in Lean~4 (see Code Availability).

This result is a direct analogue of the classical matrix Eckart--Young
theorem~\citep{eckart1936approximation}: for the rank profile specified in the
group-Fourier domain, the blockwise $\starG$ truncation is Frobenius-optimal
among all tensors obeying the same block-rank constraints, and the tubal
statement is the coupled-profile specialization. This instantiates the Eckart--Young optimality of the
transform-domain $\starM$/t-SVD~\citep{kilmer2021tensor,kilmer2011factorization}
at the group-Fourier transform $M = F_G$. This optimality is not special to $F_G$;
it holds for a broad transform family~\citep{mor2025eckart}, while $F_G$ itself is
fixed by equivariance rather than chosen (SI~Section~\ref{sec:necessity}). Two
concurrent results frame that necessity: within the commutative tubal family,
the $\starM$ axioms force the product's \emph{form} without determining the
transform~\citep{avron2025demystifying}, and group-compatibility admits an
if-and-only-if characterization in the abelian scalar case together with a
proof that \emph{no} tubal $\starM$ serves a non-abelian
group~\citep{dunbar2025tensor}; the $\starG$ necessity theorem is the
non-abelian, matrix-block resolution of exactly that impossibility. What is
specific here is the
group-algebraic reading of that transform and its three consequences: the
extension from the scalar-tube (abelian) case to non-abelian groups, where the
irrep blocks have dimension greater than one and the per-block operation is a
matrix rather than a scalar product; the further extension from finite groups to
band-limited compact groups (Theorem~\ref{thm:bandlimited}); and the result that,
on the regular and band-limited representation spaces treated here, every
group-equivariant linear map is itself a $\starG$ operation
(Corollary~\ref{cor:unification}), of which the rank-$k$ $\starG$-SVD truncation is the Frobenius-optimal member at each rank. Among symmetry-preserving approximation schemes, each under its native rank notion, this is the strongest guarantee available: Tucker/HOSVD
gives only quasi-optimal bounds with a $\sqrt{d}$ factor~\citep{desilva2008tensor},
CP decomposition is NP-hard to compute optimally, and tensor-train has no global
optimality guarantee. The error has a closed-form expression in terms of the
discarded singular tube norms, enabling principled rank selection with full
control over approximation quality. In practice, the leading singular tubes give
the smallest Frobenius reconstruction error attainable at the same
$\starG$-rank, an error available in closed form.

\paragraph{Relationship to the tubal Eckart--Young literature.} The transform-domain
matrix-mimetic program originates with Kilmer and
collaborators~\citep{kilmer2021tensor,kilmer2011factorization}, and the precise
conditions under which a tube product admits an Eckart--Young theorem, with an
extension to separable Hilbert-space tubes, have been developed for \emph{tubal}
tensors by Mor and
Avron~\citep{mor2025quasitubaltensoralgebraseparable,mor2026sufficientnecessaryconditionseckartyoung},
whose transform-domain fibers form a \emph{commutative} algebra; that
characterization is established specifically for the commutative case. The present
group-algebraic construction is broader on four axes. (i)~The transform is the
group-Fourier transform of a possibly \emph{non-abelian} group, so a
transform-domain component is a $d_\rho\!\times\!d_\rho$ \emph{matrix} block rather
than a scalar tube; the fiber algebra is non-commutative and lies outside the tubal
(commutative) setting, and the two coincide only when every irreducible
representation is one-dimensional (the abelian case).
(ii)~We establish the full matrix-mimetic toolbox (QR, symmetric
eigendecomposition, least-squares, functional calculus), not the SVD alone.
(iii)~Reading the transform as a group makes equivariance an intrinsic property and
the irrep blocks physically meaningful, which is what produces the Wigner--Eckart
diagnostics (Section~\ref{sec:wigner_eckart}) and the equivariant-network
unification (Corollary~\ref{cor:unification}). (iv)~Several symmetries compose by
$F_{G_1}\!\otimes\!F_{G_2}$ with no redesign (Theorem~\ref{thm:product_ring}). The
Eckart--Young optimality is one shared waypoint; the group-algebraic construction,
its non-abelian reach, and its physical consequences are what is specific here.
The $\starM$ product of Kilmer and collaborators (an arbitrary invertible
transform) and $\starG$ (the group-Fourier transform) are, in this sense,
distinct generalizations of the t-product that agree on the abelian case,
with neither containing the other.
A parallel development in statistical representation learning, arrived at
independently for equivariant conditional-expectation operators, block-diagonalizes
by isotypic component and applies the classical Eckart--Young theorem per block to
reach the same equivariant low-rank optimality (ICML
2026)~\citep{ordonez2025equivariant}; it
targets operator learning with non-asymptotic guarantees rather than a tensor
algebra, and develops neither a $\starG$ product with its $\starG$-SVD algorithm
nor the matrix-mimetic toolbox and formalization given here. Any priority we
claim for Eckart--Young optimality is therefore scoped: first for a
symmetry-encoding \emph{tensor algebra}, and first machine-checked, not first
per-isotypic-block.

\subsubsection*{Composing Multiple Symmetries}

\begin{theorem}[Product Groups]\label{thm:product_ring}
For $G = G_1 \times \cdots \times G_d$, the convolution tensor factorizes as
$\calT_G = \calT_{G_1} \otimes \cdots \otimes \calT_{G_d}$, and the generalized
Fourier matrix is $F_G = F_{G_1} \otimes \cdots \otimes F_{G_d}$.
\end{theorem}

\noindent Full proof is provided in the Supplementary Information (SI Section~6).

This Kronecker structure is the algebraic reason why multiple symmetries compose
without architectural redesign. The factorization of $\calT_G$ implies that the
irreps of a product group are exactly the tensor products of the factors' irreps,
so the Fourier-domain block-diagonal structure is the Kronecker product of the
individual block-diagonal structures. Concretely, for
$G = \bbZ_{n_1} \times \bbZ_{n_2}$, the Group Fourier transform
$F_G = \mathrm{DFT}_{n_1} \otimes \mathrm{DFT}_{n_2}$ computes a 2D DFT,
resolving coupled frequencies $(f_1, f_2)$ that are entirely invisible to either
factor group alone. Adding a new symmetry $G_{d+1}$ to an existing $\starG$
model requires only replacing $F_G$ with $F_G \otimes F_{G_{d+1}}$; no layers
are redesigned and no weights are reinitialized. This theorem is also what
makes the flattening convention above operational: because the irreps of the
product are the Kronecker products $\rho_1 \otimes \cdots \otimes \rho_d$ of
the factors' irreps, running any algorithm of this paper on the flattened
group, with $F_G = F_{G_1} \otimes \cdots \otimes F_{G_d}$, is
block-for-block the same computation as working multi-index, so the
$\starG$-SVD, QR, eigendecomposition, and least-squares apply unchanged to
order-$(2+d)$ tensors over product groups.

\subsection{Exact equivariance does not compound; approximate equivariance does}
\label{sec:compounding}

The distinction between exact and approximate equivariance is often treated as a
matter of degree: a method whose single-step symmetry error is $10^{-2}$ or
$10^{-5}$ is called equivariant ``up to a small error.'' The distinction is in
fact qualitative, and the reason is composition. Symmetric computations are
rarely applied once: deep networks stack layers, dynamical simulations and
rollouts iterate a step map, and self-consistent solvers loop until convergence.
The accumulated symmetry defect then grows without bound. Iterating a
\emph{fixed} step of equivariance error $\varepsilon$, the regime of rollouts
and self-consistent loops and the regime we demonstrate below, accumulates the
defect coherently, at order $M\varepsilon$ in stable norm-preserving
computation (where the growth cannot be attributed to any blow-up of the
dynamics itself), reaching $O(1)$ at depth $M \sim 1/\varepsilon$ before
saturating; a composition of $M$ \emph{distinct}, independently erring layers
accumulates more slowly, at least of order $\sqrt{M}\,\varepsilon$ in
expectation, with $M\varepsilon$ the known worst-case
bound~\citep{assaad2023vntransformer}. Every rate in this range diverges with
depth; only $\varepsilon = 0$ is depth-stable. (Cancellation between steps can
slow the accumulation, and contrived non-generic steps, an orthogonal
involution for instance, return exactly to equivariance at special depths; but
no non-equivariant invertible step is exactly equivariant at all sufficiently
large depths, and generic steps produce no exact zeros.) An operator
whose equivariance is exact, by contrast, commutes with the group at every step,
so its iterates are exactly equivariant at \emph{unbounded} depth. Exactness is
therefore a capability rather than a percentage, and the $\starG$ algebra, in
which every product, factorization, and projection is equivariant identically
in exact arithmetic, sits on the exact side of that line by construction. The
qualitative dichotomy, that exact rollouts hold their symmetry while
unconstrained ones drift away from it, has been observed concurrently by
Wang~\citep{wang2026exactequivariance}; what this section establishes is the
\emph{quantitative} form of the law: the measured single-step error
$\varepsilon$, the $M\varepsilon$ rate and its $1/\varepsilon$ depth horizon,
the fate of a physically forbidden channel, and the placement of real methods'
measured error levels on the resulting curves.

We make the law concrete on the octahedral group $O$
(\texttt{symmetry\_compounding.py}; offline and deterministic, all invariants
verified). A single step is an orthogonal, norm-preserving linear map on the
regular representation, so iteration is a stable flow rather than a collapse to
a dominant eigenmode. The exact step is $T = \exp(S)$ with $S$ an equivariant
antisymmetric generator obtained by Reynolds symmetrization; its measured
single-step equivariance error is at floating-point zero
($2.5\times10^{-16}$), and because
it is block-diagonal across irreps it can never move energy into a
symmetry-forbidden channel. The approximate step is $T_\varepsilon = \exp(S +
\varepsilon B)$ with $B$ a generic non-equivariant antisymmetric generator: an
orthogonal map whose single-step equivariance error is \emph{measured}, not
assumed, and swept over $\varepsilon \approx 7\times10^{-2}$, $7\times10^{-3}$,
$7\times10^{-4}$, the range that finite data, finite augmentation, relaxed
equivariance, or soft symmetry penalties realistically leave behind. This range
is documented, not stipulated: measured equivariance errors of trained models
span it, from augmentation-trained networks at the top of the range down to
architecturally constrained ones~\citep{gruver2023lie,wang2022approximately},
and our own augmented baselines sit near $10^{-2}$
(Section~\ref{sec:enn_comparison}). The sweep is anchored on both sides of the
line by two operators that are \emph{trained} rather than constructed, fitted
to the same exact step from the same $256$ example pairs. An orbit-augmented
multilayer perceptron, the standard augmentation route to approximate
invariance, trains to a relative fit error of $8\times10^{-3}$ with each pair
presented in five group frames, yet its measured single-step equivariance
error is $\varepsilon \approx 1.7\times10^{-2}$, squarely inside the swept
range. An architecturally equivariant linear layer, the least-squares fit over
the commutant of the regular representation and hence, by
Corollary~\ref{cor:unification}, precisely the linear layer of an equivariant
network, recovers the step with measured error $9\times10^{-16}$. The same
data, fitted through two standard routes, lands on opposite sides of the
exact/approximate line.

Figure~\ref{fig:compounding} shows the result of iterating each operator
$M = 1$ to $400$ times. The equivariance drift of the exact iterate stays at
floating-point noise for all depths, $2\times10^{-16}$ at $M{=}1$ and
$1.8\times10^{-14}$ at $M{=}400$, while every approximate operator compounds as
the law predicts: the $\varepsilon \approx 7\times10^{-2}$ operator reaches
$O(1)$ symmetry violation by $M \approx 50$, the $7\times10^{-3}$ one by
$M \approx 400$, and the $7\times10^{-4}$ one is at $10\%$ violation by
$M{=}400$ and still rising. The trained operators obey the same law at
inference: the augmented perceptron, unrolled exactly as a rollout would
unroll it, drifts as $\sim\!M\varepsilon$ from its measured
$\varepsilon = 1.7\times10^{-2}$ to total symmetry loss by $M{=}400$, with the
forbidden-channel fraction of its state growing to $17\%$, while the trained
equivariant layer holds its drift at floating-point noise ($10^{-13}$ at
$M{=}400$). The forbidden-channel energy tells the same story
physically. Here the forbidden channel is the $A_2$ isotypic component, the
one-dimensional pseudoscalar irreducible representation of $O$ (the sign of
the permutation action on the body diagonals; see the irrep table in
Section~\ref{sec:wigner_eckart}), absent from the initial state by
construction: a state initialized in the symmetry-allowed subspace acquires
\emph{zero} $A_2$ energy under the exact iterate at every
depth ($\le 4\times10^{-15}$), whereas each approximate iterate leaks energy
into the forbidden channel and the leak grows with depth. The gap between the
exact and approximate classes is fourteen orders of magnitude at $M{=}400$ and
\emph{widens} with exactly the depth and composition that modern pipelines
scale.

Three readings of this experiment organize the rest of the paper. First, it is
not a strawman: $\varepsilon$ is measured, the sweep covers the error levels
measured for real trained models~\citep{gruver2023lie,wang2022approximately},
two of the curves are themselves trained models rather than constructions,
and the compounding is a mathematical
law, not a weakness of a particular baseline. An architecturally exact
equivariant network is on the exact side of the line; the dividing line is
exact versus approximate equivariance, and $\starG$ is the closed-form,
training-free member of the exact class whose equivariance is additionally
machine-certified (Lean~4, Discussion); Table~\ref{tab:comparison} summarizes
the structural contrast. Second, the forbidden-channel leak is
precisely the failure mode that the crystal physicality layer of
Section~\ref{sec:crystals} repairs on real learned predictors, whose trained
outputs leak $10^{-2}$ of their energy into symmetry-forbidden channels; the
projection restores the exact class post hoc, on any predictor's output.
Third, the capability costs nothing: exactness is a property of the algebra's
multiplication rule, not a constraint traded against expressivity or accuracy.

\begin{figure}[!ht]
\centering
\includegraphics[width=\textwidth]{fig_compounding.pdf}
\caption{\textbf{Exact equivariance does not compound; approximate equivariance
does.} One-step operators on the regular representation of the
octahedral group $O$, iterated $M$ times; the exact operator commutes with the
group to floating-point precision, the approximate operators have measured
single-step equivariance error $\varepsilon$, swept over the range measured
for trained models in
practice~\citep{gruver2023lie,wang2022approximately}. Two operators are
trained on the same step from identical data: an orbit-augmented perceptron
unrolled at inference (approximate class, measured
$\varepsilon \approx 1.7\times10^{-2}$) and an equivariant linear layer
(exact class, $\varepsilon \approx 10^{-15}$). (\textbf{Left}) Relative equivariance violation of the $M$-fold
iterate: the exact operators stay at floating-point noise
($1.8\times10^{-14}$ at $M{=}400$ for the $\starG$ step) while every
$\varepsilon > 0$, trained or synthetic, compounds as
$\sim\!M\varepsilon$ to $O(1)$. (\textbf{Right}) Energy in the
symmetry-forbidden $A_2$ channel of a state initialized in the allowed
subspace: exactly zero under the exact iterates at every depth, growing leakage
under every approximate one. The gap is fourteen orders of magnitude at
$M{=}400$ and widens with depth.}
\label{fig:compounding}
\end{figure}

\begin{table}[h]
\centering
\caption{Structural comparison: ENNs vs.\ $\starG$ algebra.}
\label{tab:comparison}
\small
\begin{tabular}{@{}p{2.4cm}p{5.8cm}p{7.0cm}@{}}
\toprule
\textbf{Feature} & \textbf{ENNs} & \textbf{$\starG$ Algebra} \\
\midrule
Symmetry     & Architectural constraint             & Algebraic property \\
Group domain & Continuous, truncated to finite $l_{\max}$ & Finite, or compact via band-limit (Thm.~\ref{thm:bandlimited}) \\
Composition  & Fixed at architecture design time    & $F_{G_1} \otimes F_{G_2}$, no redesign (Thm.~\ref{thm:product_ring}) \\
Optimality   & None established for learned layers  & Eckart--Young (Thm.~\ref{thm:eckart_young}) \\
Equivariance & Exact, up to float precision         & Exact; machine-verified (Lean~4) \\
Training     & Iterative gradient descent           & Closed form (SVD + ridge) \\
\bottomrule
\end{tabular}
\end{table}

\subsection{Symmetry diagnostics: Wigner--Eckart signatures from data}
\label{sec:wigner_eckart}

A first, distinctive use of the algebra is not prediction but diagnosis: the
$\starG$ decomposition can \emph{surface} physical symmetry structure that was
not provided as input. The Wigner--Eckart
theorem states that scalar observables ($l\!=\!0$) couple only to the trivial
representation, vector observables ($l\!=\!1$) require the fundamental
representation, and rank-2 tensor observables ($l\!=\!2$) require the $l\!=\!0$
and $l\!=\!2$ channels. A generalized Wigner--Eckart theorem has already been
brought into equivariant learning by Lang and Weiler~\citep{lang2021wigner} to
\emph{parameterize} steerable convolution kernels; here we use the same
selection-rule structure in the opposite direction, as a diagnostic read out of
data rather than as a design constraint. We test whether the $\starG$ framework
surfaces signatures consistent with these selection rules from molecular geometry
data alone.

\subsubsection*{Setup}

We work with the chiral octahedral group $O$
(order 24, a subgroup of SO(3)) whose five irreducible representations correspond
directly to angular momentum channels:

\begin{center}
\begin{tabular}{@{}llcl@{}}
\toprule
\textbf{Irrep} & \textbf{Dim} & \textbf{$l$-channel} & \textbf{Physical content} \\
\midrule
A$_1$ & 1 & $l=0$ & Scalar (spherically symmetric) \\
A$_2$ & 1 & $l=0$ & Pseudoscalar \\
E     & 2 & $l=2$ & Quadrupole ($d$-wave) \\
T$_1$ & 3 & $l=1$ & Vector ($p$-wave, dipolar) \\
T$_2$ & 3 & $l=2$ & Quadrupole ($d$-wave) \\
\bottomrule
\end{tabular}
\end{center}

\noindent We study two real datasets. On the full QM9 set (130{,}831 molecules,
3 seeds) we featurise each molecule under all 24 octahedral rotations and take
the rotation-invariant per-irrep Fourier power, reporting the test \Rsq{} of each
irrep's features alone for the scalar properties HOMO--LUMO gap, isotropic
polarizability $\alpha$, and ZPVE, and for the dipole magnitude
$|\boldsymbol{\mu}|$. To test the rank-2 selection rule \emph{directly}, rather
than through a charge-derived proxy, we use the QM7-X
dataset~\citep{hoja2021qm7x}, which provides the
full molecular polarizability \emph{tensor} (a real symmetric $3\times3$ matrix),
and decompose it into octahedral irreps.

\subsubsection*{Results}

Table~\ref{tab:wigner}, Figure~\ref{fig:wigner}, and
Table~\ref{tab:qm7x_polarizability} show five findings.

\textbf{(i) Scalars live in A$_1$.} The gap, ZPVE, and isotropic polarizability
are dominated by the A$_1$ ($l\!=\!0$) channel (\Rsq{} $= 0.32, 0.85, 0.61$),
while the A$_2$ (pseudoscalar) channel contributes essentially nothing
($|\Rsq| < 0.001$). Realised correctly as the sign of the $S_4$ action on the
body diagonals, A$_2$ vanishes for these true-scalar properties, as a genuine
pseudoscalar must.

\textbf{(ii) Directional observables shift to T$_1$.} The dipole magnitude
$|\boldsymbol{\mu}|$, though a scalar, draws most of its predictive signal from
the T$_1$ ($l\!=\!1$) channel (\Rsq{}$_{T_1} = 0.28$ vs.\ A$_1 = 0.13$): it is the
norm of a vector, and the directional content lives in $l\!=\!1$.

\textbf{(iii) Controls confirm the structure is real and group-dependent.}
Permuting the labels collapses every channel to \Rsq{} $\approx 0$ (no
over-fitting). Replacing $F_G$ with a random orthonormal basis of the same block
sizes spreads the signal uniformly across all blocks (each block reaches the
full-model \Rsq{} $\approx 0.34$ for the gap), erasing the selectivity entirely.
It is therefore the octahedral group structure, not generic block features or a
coordinate convention, that produces the A$_1$/T$_1$ pattern.

\textbf{(iv) The exact rank-2 selection rule.} On the real QM7-X polarizability
tensor, the octahedral decomposition is A$_1$ (isotropic) $0.93$, E $0.03$,
T$_2$ $0.04$, and T$_1$ \emph{exactly} $0$ (maximum over 3{,}000 molecules: $0$ to
machine precision; Figure~\ref{fig:wigner}b). A symmetric rank-2 operator has no
$l\!=\!1$ component, and $\starG$ recovers this selection rule \emph{exactly} from
the data, not merely statistically.

\textbf{(v) The irrep index is an explicit, restrictable feature axis.} Because
the decomposition is available before any training, a model can be
\emph{restricted} to a single irrep channel by construction. On the QM7-X
polarizability components, a two-parameter ridge model built from the $E$ irrep
alone predicts the $E$-channel magnitude and is structurally unable to predict
the $T_2$ magnitude, and vice versa ($96$--$97\%$ cross-selectivity;
Table~\ref{tab:qm7x_polarizability}). This is an operational affordance of the
representation rather than a representational edge over trained networks:
given the full input, $\starG$ with all irreps, symmetry-blind baselines, and
three ENN architectures all sit near zero cross-selectivity, and a trained ENN's
internal irrep structure can be recovered with a bespoke post-hoc probe
(SI~Section~\ref{sec:enn_probe}). What the algebra adds is the channel
restriction as an explicit input-side control, with no probe and no training.

\begin{table}[h]
\centering
\caption{\textbf{QM7-X polarizability components: cross-selectivity table.}
For each method we train two single-target models, one predicting the
$E_g$ irrep magnitude $\|\boldsymbol{\alpha}_E\|$ of the molecular
polarizability tensor, the other predicting the $T_{2g}$ magnitude
$\|\boldsymbol{\alpha}_{T_2}\|$, on identical 60/20/20 per-molecule
splits across three seeds (42, 43, 44). Cross-selectivity is
$\mathrm{cs} = (R^2_{\mathrm{on}} - R^2_{\mathrm{off}}) / (R^2_{\mathrm{on}}
+ R^2_{\mathrm{off}} + 10^{-12})$, where the ``on'' model targets the
indicated component and the ``off'' model is the architecturally-
identical model trained on the other component. A method whose two
single-target architectures achieve similar $R^2$ values has
cross-selectivity near zero, indicating that the architecture treats
the two targets symmetrically. A method with very asymmetric $R^2$ has
cross-selectivity near one, indicating a structural prior that aligns
with one component and not the other.}
\label{tab:qm7x_polarizability}
\begin{tabular}{@{}lrccc@{}}
\toprule
\textbf{Method} & \textbf{Params} & $R^2(\|\boldsymbol{\alpha}_E\|)$ & $R^2(\|\boldsymbol{\alpha}_{T_2}\|)$ & \textbf{Cross-selectivity} \\
\midrule
$\starG$ E-only + Ridge          & $2$         & $0.591$ & $0.012$ & $\mathbf{96.2\%}$ \\
$\starG$ T$_2$-only + Ridge      & $2$         & $0.008$ & $0.593$ & $\mathbf{97.3\%}$ \\
$\starG$ all irreps + Ridge      & $6$         & $0.622$ & $0.609$ & $1.1\%$ \\
\midrule
Raw features + Ridge             & $37$        & $0.481$ & $0.512$ & $-3.2\%$ \\
Coulomb matrix + Ridge           & $277$       & $0.442$ & $0.446$ & $-0.5\%$ \\
Raw + MLP$(128, 64, 32)$         & $15{,}105$  & $0.586$ & $0.591$ & $-0.5\%$ \\
Coulomb + MLP$(128, 64, 32)$     & $45{,}825$  & $0.495$ & $0.456$ & $4.2\%$ \\
\midrule
SchNet (PyG)                     & $455{,}809$ & $0.664 \!\pm\! 0.026$ & $0.662 \!\pm\! 0.018$ & $0.15\%$ \\
e3nn (SE(3)-equivariant)         & $794{,}032$ & $0.708 \!\pm\! 0.027$ & $0.698 \!\pm\! 0.006$ & $0.7\%$ \\
MACE                             & $1{,}041{,}808$ & $0.687 \!\pm\! 0.032$ & $0.672 \!\pm\! 0.022$ & $1.1\%$ \\
\bottomrule
\end{tabular}
\medskip

\footnotesize
\noindent\emph{The three ENNs achieve modest predictive $R^2$ in the
$0.66$--$0.71$ range on each polarizability component individually,
matching the $\starG$ all-irreps + Ridge baseline at $6$ parameters.
\textbf{The $\starG$ $E$-only and $T_2$-only models, with $2$ trainable
parameters each, retain $96.2\%$ and $97.3\%$ cross-selectivity} as a
direct consequence of restricting the algebra's per-irrep features to a
single irrep. The three ENN architectures, given the full input rather than a
restricted one, stay near $1\%$, as does $\starG$ itself with all irreps present;
the effect measures input restriction, not a representational separation the
networks cannot make. The operative point of finding (v) is
that $\starG$ exposes the irreducible-representation decomposition as a
controllable input axis that the standard ENN workflow does not.}
\end{table}

By decomposing data over a physically meaningful group, $\starG$ thus acts as a
\emph{spectroscope for symmetry}: it reads off which angular-momentum channel
carries which observable, and it certifies the T$_1$-forbiddenness of the
polarizability exactly. The exact zero is structural rather than statistical:
any symmetric rank-2 input has no $l\!=\!1$ component, so this row certifies
that the pipeline preserves the algebraic identity to machine precision (a
learned model has no such guarantee; see the crystal experiment below), rather
than discovering a contingent property of the molecular data. We frame this as
a diagnostic that surfaces known
selection rules from data, not as a re-derivation of the Wigner--Eckart theorem.

\begin{table}[h]
\centering
\caption{Per-irrep test \Rsq{} on QM9 (full set, 130{,}831 molecules, 3 seeds;
chiral octahedral $O$). Scalars are A$_1$-dominated; the dipole magnitude
$|\boldsymbol{\mu}|$ shifts to T$_1$. The corrected A$_2$ (sign) representation is
$\approx 0$ throughout. Control rows: shuffling labels collapses every channel; a
random orthonormal basis of the same block sizes erases the selectivity (every
block reaches the full-model \Rsq).}
\label{tab:wigner}
\begin{tabular}{@{}lcccccr@{}}
\toprule
\textbf{Property} & \textbf{A$_1$} & \textbf{A$_2$} & \textbf{E} &
\textbf{T$_1$} & \textbf{T$_2$} & \textbf{all} \\
\midrule
HOMO--LUMO gap                   & \textbf{0.319} & 0.000 & 0.009 & 0.058 & 0.002 & 0.344 \\
$\alpha$ (isotropic)             & \textbf{0.614} & 0.000 & 0.052 & 0.005 & 0.066 & 0.632 \\
ZPVE                             & \textbf{0.854} & 0.000 & 0.000 & 0.031 & 0.032 & 0.856 \\
$|\boldsymbol{\mu}|$ (magnitude) & 0.130 & 0.000 & 0.007 & \textbf{0.278} & 0.000 & 0.359 \\
\midrule
\;\;gap, shuffled labels         & 0.000 & 0.000 & 0.000 & 0.000 & 0.000 & 0.000 \\
\;\;gap, random basis            & 0.343 & 0.340 & 0.345 & 0.346 & 0.346 & 0.348 \\
\bottomrule
\end{tabular}
\end{table}

\begin{figure}[!ht]
\centering
\includegraphics[width=\textwidth]{irrep_bars.pdf}
\caption{\textbf{Symmetry spectroscopy and the exact polarizability selection
rule.} (\textbf{a})~Per-irrep test \Rsq{} on full QM9: the scalar gap, ZPVE, and
isotropic polarizability are A$_1$ ($l\!=\!0$) dominated, while the dipole
magnitude $|\boldsymbol{\mu}|$ shifts to T$_1$ ($l\!=\!1$); the corrected A$_2$
(sign) representation is $\approx 0$. (\textbf{b})~Octahedral irrep energy of the
real QM7-X polarizability tensor: A$_1 \oplus$ E $\oplus$ T$_2$, with T$_1$
\emph{exactly} zero over all 3{,}000 molecules, the exact Wigner--Eckart selection
rule for a symmetric rank-2 operator.}
\label{fig:wigner}
\end{figure}

\subsection{Exact tensorial selection rules in crystals}
\label{sec:crystals}

The experiments so far apply $\starG$ to molecular data whose symmetry is only
\emph{approximate}. Crystalline solids are the setting where the symmetry is exact
and discriminative: the space group is an exact symmetry of the structure, and
tensorial properties decompose by irreducible representation with symmetry-forced
selection rules that are physically central. On the $1{,}181$ inorganic crystals of
the de Jong elastic-tensor dataset~\citep{dejong2015charting}, the exact SO(3)
decomposition of the elastic tensor into its $l=0,2,4$ channels reproduces these
rules from the raw tensors with no symmetry input: the $452$ cubic crystals carry a
median $l=2$ energy fraction of $1.4\times10^{-10}$, seven orders of magnitude
below the mean of $6.6\times10^{-5}$ that imperfect DFT relaxations leave
($99\%$ below $10^{-3}$), their
anisotropy residing entirely in $l=4$, whereas every uniaxial or lower-symmetry
class (hexagonal, trigonal, tetragonal, orthorhombic, monoclinic) carries a nonzero
$l=2$ channel. This is the crystalline counterpart of the Wigner--Eckart signatures
above (Section~\ref{sec:wigner_eckart}), now exact rather than a
discretised diagnostic.

Two prediction experiments locate $\starG$'s advantage precisely where symmetry,
not chemistry, carries the signal. First,
predicting properties from structure with a symmetry-blind gradient-boosted
regressor on a $\starG$ per-site descriptor, the descriptor's advantage over a
composition-only baseline concentrates exactly where symmetry, not chemistry,
carries the signal: on the composition-dominated moduli it is marginal
($\log K_{\mathrm{VRH}}$ $R^2$ $0.84\!\to\!0.86$, $\log G_{\mathrm{VRH}}$
$0.67$, unchanged within run-to-run variation), whereas on the anisotropy channels, which are pure symmetry
content, it reaches $R^2 = 0.27$ ($l=2$, uniaxial) and $0.24$ ($l=4$, cubic) and
lifts the anisotropy \emph{type} $c_4=l_4/(l_2+l_4)$ from $R^2 = 0.10$ to $0.43$.
Second, the $\starG$ decomposition supplies a closed-form \emph{physicality layer}
that enforces the crystal's selection rules exactly on the output of \emph{any}
predictor. Symmetry-aware tensor prediction is an active line, and group
averaging itself is not the contribution: equivariant architectures build
crystal-system constraints into the model (MatTen for the elasticity tensors
of all seven crystal systems~\citep{wen2024matten}; GMTNet, whose space-group
projection is likewise a Reynolds average, applied inside an equivariant
network at numerical tolerance~\citep{yan2024gmtnet}), canonicalization
wrappers orient the crystal before prediction but by their own account do not
enforce selection rules~\citep{hua2024goectp}, inference-time corrections
target the zero elements~\citep{jin2026ceitnet}, and post-hoc symmetrization
of tensors by group averaging is standard practice in materials
software~\citep{ong2013pymatgen}. What the algebra adds is the projector's
provenance and the guarantees that provenance buys: it is the orthogonal
idempotent of the verified $\starG$ irrep calculus, the point-group Reynolds
operator and the $\Gamma$-point instance of the Brillouin-zone
selection-rule projector of SI~Section~\ref{sec:spacegroups}; it acts at
machine precision rather than at solver tolerance; it applies post hoc to an
already-trained model, with a proof (below) that it cannot increase error,
where output averaging imposed during training has been reported to degrade
accuracy~\citep{yan2024gmtnet}; and its effect on downstream screening is
quantified. Generalizing the cubic
case above to all seven crystal systems and applied to the $1{,}181$ real
tensors (which span the six crystal systems represented in the dataset), it
drives the symmetry-forbidden channel energy to machine precision for every system
($\le 10^{-30}$; the cubic $l=2$ fraction from $6.6\times10^{-5}$ to
$1.6\times10^{-16}$) while leaving the allowed channels unchanged. The consequence
for prediction is the point. An unconstrained tensor regressor violates these rules
universally: across three seeds, every prediction of a tuned gradient-boosted model
leaks forbidden energy (mean fraction $5\times10^{-3}$), which the projection removes
to machine zero. The same holds for a trained crystal graph network: across a sweep
of training-set sizes its predicted tensors leak $0.9$--$2.6\times10^{-2}$ of their
energy into forbidden channels, which the projection drives to $1.5\times10^{-32}$ at
every size (\texttt{crystal\_gnn\_data\_efficiency.py}), so the guarantee is genuinely
architecture-agnostic, holding on a deep learned model's output as on a boosted one.
The demonstration deliberately targets predictors that do not build the
constraint in, the class most widely deployed; for architecture-constrained
models such as MatTen~\citep{wen2024matten}, which satisfy the rules to float
tolerance on exactly symmetric inputs, the layer is complementary rather than
corrective, supplying the exact-zero certificate and the error guarantee below.
This is the single-step instance of the compounding law of
Section~\ref{sec:compounding} caught in the wild, and the projection restores the
exact class post hoc: left unprojected, a predictor whose tensors feed further
symmetric computation would compound this leakage with depth.
Composing it with a standard positive-semidefinite projection
(nearest-PSD by eigenvalue clipping) additionally eliminates every violation of Born
mechanical stability the model emits, $1.7$ of $1{,}181$ predictions for the tuned
model and $80$ of $1{,}181$ for a weaker linear one, both driven to zero, whereas
the symmetry projection alone leaves the weak model's stability violations at $63$.
The layer never alters symmetry-invariant screening (bulk-modulus ranking, Spearman
$0.92$, unchanged by construction) and modestly improves tensor reconstruction
($R^2$ by $+0.009$ for the tuned model, $+0.047$ for the weak one). This is a
provable, architecture-agnostic guarantee of physical consistency rather than an
average-accuracy gain: it removes physically impossible predictions,
symmetry-forbidden and mechanically unstable, with a certificate that no
unconstrained model provides. The downstream effect is concentrated where symmetry,
not magnitude, carries the signal: on ranking crystals by anisotropy \emph{type}
($c_4$), the layer lifts the Spearman correlation with the true type from $0.32$ to
$0.75$ (and from $0.16$ to $0.73$ for a weaker model), because the raw predictor's
forbidden $l=2$ leakage corrupts the cubic crystals' type (true $c_4\approx1$) and
the projection restores it exactly, whereas on anisotropy \emph{magnitude}, which a
tuned model already captures, the effect is small.

\begin{figure}[!ht]
\centering
\includegraphics[width=\textwidth]{fig_enforcement.pdf}
\caption{\textbf{The physicality layer in action
(\texttt{make\_enforcement\_fig.py}; artifacts committed).}
(\textbf{Left}) Forbidden-channel energy of a trained crystal graph network's
predicted elastic tensors across training sizes: the raw model leaks
$0.9$--$2.6\times10^{-2}$, above even the imperfect-DFT data floor, and the
$\starG$ projection drives every prediction to machine zero
($\sim\!1.5\times10^{-32}$).
(\textbf{Middle}) Discovery consequence: viable materials wrongly rejected by
a raw Born-stability screen at each training size, all recovered by the
composed (symmetry + nearest-PSD) layer; at the smallest size the fully
mean-reverted model rejects nothing to begin with.
(\textbf{Right}) Ranking crystals by anisotropy \emph{type} ($c_4$): the
projection lifts the Spearman correlation from $0.32$ to $0.75$ (3 seeds,
error bars $\pm$1 s.d.) by exactly
restoring the cubic crystals' forbidden-channel zeros
(\texttt{anisotropy\_screening.py}).}
\label{fig:enforcement}
\end{figure}

The guarantee has a simple reason and a sharp regime (Figure~\ref{fig:enforcement}). Because the projector is
orthogonal and the true tensor lies in its range, the squared error splits exactly
into an approximation term and the forbidden-channel energy, so projection can only
reduce the error, by precisely the forbidden mass it removes
(SI~Proposition~\ref{prop:enforce_guarantee}). The benefit is therefore largest in
the data-scarce discovery regime, where a model trained on the few measured tensors
must extrapolate to many candidates. Sweeping the training set down to a few tens of
crystals, the projection reduces the reconstruction error of $98$--$100\%$ of
individual predictions (by $5$--$6\%$ once out of the trivial mean-reverting limit),
and, at training sets of $71$ crystals and larger (below which the fully
mean-reverted model rejects nothing), a raw Born-stability screen wrongly
rejects $10$--$20$ otherwise-viable materials
per held-out set, all of which the composed layer recovers; this false-rejection
count grows linearly with the size of a screen, so at discovery scale it is the
difference between discarding and retaining thousands of viable candidates
(\texttt{enforcement\_screening.py}).

\paragraph{Equal-input data efficiency: the symmetry-structured representation as a
sample-complexity prior.} The value of the symmetry structure can be isolated on
identical information. Both representations of a crystal's local environment are built
from the same chemistry-weighted radial, dipole, and quadrupole moments of each site's
neighbours; the $\starG$ descriptor is their octahedral group-Fourier power spectrum,
exactly invariant under the octahedral group $O$, whereas a symmetry-blind descriptor
reads the same moments in one chosen frame (keeping, in fact, \emph{more} information,
since the power spectrum discards phase). Because the targets are octahedral scalar
invariants, a held-out crystal presented in any octahedral orientation carries the
same label, so testing on a randomly reoriented held-out set isolates the
sample-complexity cost of learning the invariance that $\starG$ has by construction
(\texttt{crystal\_data\_efficiency.py}, Table~\ref{tab:data_efficiency}). On the
symmetry-carried anisotropy channel $c_4$, $\starG$ ranks crystals two to three times
better than the symmetry-blind representation of the same information at every
training size, and better than that representation even when it is given explicit
orbit augmentation over all $24$ frames. On the composition-carried bulk modulus
$\starG$ is the most sample-efficient in the low-data regime (Spearman $0.84$ at $40$
crystals), the symmetry-blind representation catching up only at the largest size
swept, and then only with explicit orbit augmentation ($0.910$ versus $0.894$ at
$640$ crystals), the expected featurized-versus-learned crossover; on the
symmetry-carried channel $c_4$ no crossover occurs anywhere in the sweep. The $\starG$ score is
identical on canonically- and rotated-oriented test sets (an empirical confirmation of
its exact invariance), whereas the symmetry-blind representation degrades under
reorientation. The advantage is therefore the sample-complexity of built-in
invariance, apples-to-apples on identical input, not a difference in information; a
full equivariant-network head-to-head fed the same structures is the natural extension
(\texttt{crystal\_gnn\_headtohead.py}).

\begin{table}[h]
\centering
\caption{Equal-input crystal data efficiency: rotated-frame test ranking (Spearman,
mean over 3 seeds) versus training-set size, for the octahedral-invariant $\starG$
descriptor and a symmetry-blind representation of the \emph{identical} local-environment
moments (read in one frame, and with orbit augmentation over the $24$ octahedral
frames). Anisotropy $c_4$ is symmetry-carried; $\log K$ is composition-carried.
$\starG$'s rotated- and canonical-frame scores coincide (exact invariance).}
\label{tab:data_efficiency}
\begin{tabular}{@{}llccc@{}}
\toprule
\textbf{Target} & \textbf{$N_{\mathrm{train}}$} & \textbf{$\starG$} & \textbf{sym-blind} & \textbf{sym-blind\,+\,aug} \\
\midrule
anisotropy $c_4$ & $40$  & \textbf{0.358} & 0.114 & 0.179 \\
                 & $160$ & \textbf{0.483} & 0.140 & 0.237 \\
                 & $640$ & \textbf{0.614} & 0.217 & 0.306 \\
\midrule
$\log K$         & $40$  & \textbf{0.840} & 0.704 & 0.763 \\
                 & $160$ & 0.873 & 0.825 & 0.876 \\
                 & $640$ & 0.894 & 0.870 & 0.910 \\
\bottomrule
\end{tabular}
\end{table}

\subsection{Optimal symmetry-preserving compression}
\label{sec:compression}

Theorem~\ref{thm:eckart_young} is, at heart, a compression statement: the
$\starG$-SVD gives the best $\starG$-rank-$k$ approximation of group-structured
data. We test it directly on real molecular tensors, against the matrix SVD and
classical tensor decompositions, following the comparison convention of the
transform-domain t-SVD~\citep{kilmer2021tensor}: at matched rank $k$, the matrix
baseline is the SVD of the lateral-slice matricization (each lateral slice
vectorised to a column), and the tensor methods are the $\starG$-SVD and
Tucker/HOSVD. Two real datasets, each arranged as an order-3 tensor over the
chiral octahedral group: (i)~QM9 molecular features sampled over the 24
octahedral frames, and (ii)~QM7-X molecular polarizability tensors $\alpha$,
which transform as $\alpha \mapsto R_g\, \alpha\, R_g^{\!\top}$ and therefore
decompose as $A_1 \oplus E \oplus T_2$.

The advantage is largest at low rank (Figure~\ref{fig:compression}). On both datasets the
$\starG$-SVD reconstructs the data \emph{exactly at rank~3} (relative error
$<10^{-6}$), whereas at the same rank the matrix SVD still has $27\%$ (QM9) and
$21\%$ (QM7-X) error and needs rank~11 and rank~6 respectively to reach exact
recovery; Tucker is uniformly worse. This is not a claim that $\starG$-rank and
matrix rank are the same notion; they count different algebraic objects. The
comparison asks how many degrees of freedom each representation needs, under its
native rank, to reconstruct the same structured tensor. Empirically, at \emph{every} rank $k$ below full the
$\starG$-SVD error is at most the matrix SVD error, and the two nearly coincide
only on i.i.d.\ random tensors (a few percent apart), the worst case predicted by
the transform-domain theory~\citep{kilmer2021tensor}. The reason is structural:
molecular tensors over a symmetry group are intrinsically low $\starG$-rank (the
polarizability is literally rank~3, one tube per irrep $A_1, E, T_2$), and only
the group-Fourier transform exposes this; a matricization that ignores the group
cannot. This is the low-$\starG$-rank structure that Theorem~\ref{thm:eckart_young}
compresses optimally within its rank class, made concrete on real data, and it is a
property no matrix or generic-tensor method shares.

\begin{figure}[!ht]
\centering
\includegraphics[width=\textwidth]{compression.pdf}
% Generated by python/large_scale/make_compression_fig.py
\caption{\textbf{Optimal symmetry-preserving compression.} Relative
reconstruction error versus rank for the $\starG$-SVD, the matrix SVD (of the
lateral-slice matricization, matched rank), and Tucker/HOSVD, on real molecular
tensors. (\textbf{a})~QM9 features over the 24 octahedral frames;
(\textbf{b})~QM7-X polarizability tensors. The $\starG$-SVD reaches exact
recovery at rank~3 on both, where the matrix SVD retains $21$--$27\%$ error,
because the group structure makes the data intrinsically low $\starG$-rank.
Empirically the $\starG$-SVD error is at most the matrix SVD error at every rank;
Theorem~\ref{thm:eckart_young} establishes the $\starG$-SVD's optimality within its
own $\starG$-rank class.}
\label{fig:compression}
\end{figure}

\subsection{Synthetic validation}
\label{sec:synthetic_exp}

A controlled synthetic setting isolates, statically, the same exactness
dichotomy that Section~\ref{sec:compounding} demonstrated dynamically: when
the data's symmetry is exact, algebraic invariance sits at machine precision
while learned invariance is approximate at best. We generated
1{,}000 synthetic molecules with exact $\bbZ_{12}$ rotational symmetry and
compared $\starG$-SVD with ridge regression against four baselines (Augmented
MLP, Neural $\starG$, Standard MLP, and Invariant MLP) across three metrics:
predictive accuracy (\Rsq), rotational invariance (variance of predictions under
unseen rotations), and parameter efficiency. Cyclic structure was verified at
machine precision ($4.2 \times 10^{-16}$). The target property $\mathbf{y}$ combines mean interatomic distance, distance variance, and atomic number contributions, representing a rotationally invariant scalar analogous to size-dependent molecular properties such as polarizability.

Results are summarized in Table~\ref{tab:synthetic} and Figure~\ref{fig:synthetic}.
The $\starG$-SVD achieves perfect prediction (\Rsq{} $= 1.000 \pm 0.000$) and
exact invariance (rotation variance $5.8 \times 10^{-31}$) using only 101
parameters, compared to 5{,}249--14{,}465 for all neural baselines. The Standard
MLP achieves \Rsq{} $= 0.377$ with rotation variance $0.14$, confirming that
without explicit symmetry handling the model neither learns well nor respects the
symmetry. The Augmented MLP (\Rsq{} $= 0.998$) achieves near-perfect accuracy
but retains residual rotation variance ($3.7 \times 10^{-5}$) five orders of
magnitude larger than $\starG$-SVD, illustrating that augmentation approximates
but does not algebraically guarantee invariance. Figure~\ref{fig:synthetic}b
shows a $26$ to $30$ order-of-magnitude gap in rotation variance between
$\starG$-SVD and the approximately invariant neural baselines (augmented and
standard MLP).

\begin{table}[h]
\centering
\caption{Synthetic validation ($\bbZ_{12}$, 1{,}000 molecules, 3 seeds).}
\label{tab:synthetic}
\begin{tabular}{@{}lccc@{}}
\toprule
\textbf{Method} & \textbf{Test \Rsq} & \textbf{Rot.\ Variance} & \textbf{Params} \\
\midrule
$\starG$-SVD + Ridge & $\mathbf{1.000 \pm 0.000}$ & $5.8 \times 10^{-31}$ & \textbf{101} \\
Augmented MLP       & $0.998 \pm 0.000$           & $3.7 \times 10^{-5}$  & 5{,}249 \\
Neural $\starG$     & $0.697 \pm 0.063$           & $1.0 \times 10^{-30}$ & 8{,}641 \\
Standard MLP        & $0.377 \pm 0.054$           & $1.4 \times 10^{-1}$  & 5{,}249 \\
Invariant MLP       & $0.327 \pm 0.147$           & $\sim 0$              & 14{,}465 \\
\bottomrule
\end{tabular}
\end{table}

\begin{figure}[!ht]
\centering
\includegraphics[width=\textwidth]{fig2_synthetic.pdf}
\caption{\textbf{Synthetic validation ($\bbZ_{12}$).}
(\textbf{a})~Test \Rsq{}.
(\textbf{b})~Rotation variance (log scale): a 26-to-30-orders-of-magnitude gap
between
$\starG$-SVD and the approximately invariant baselines, the static counterpart
of the compounding dichotomy of Section~\ref{sec:compounding}.
(\textbf{c})~Parameter efficiency (Pareto frontier): on this exactly symmetric
synthetic set, $\starG$-SVD dominates all
baselines on both axes simultaneously.
(\textbf{d})~Multi-metric normalized summary scores.}
\label{fig:synthetic}
\end{figure}


\subsection{Composing symmetries: product groups}
\label{sec:product_exp}

To test whether the algebraic composition theorem
(Theorem~\ref{thm:product_ring}) translates into empirical performance, we
constructed a task with two independent, commuting symmetries:
$G_1 = \bbZ_6$ (discrete rotations in the $xy$-plane) and
$G_2 = \bbZ_4$ (periodic translations along $z$). The target quantity was
designed to be dominated by coupled 2D Fourier frequencies $(f_1, f_2)$ that
require both symmetries simultaneously; neither factor alone can resolve these
modes. This setting models the physical situation in which a molecular property
depends on two structural degrees of freedom simultaneously, as is common in
materials with layered or helical symmetry.

Results are shown in Table~\ref{tab:product} and Figure~\ref{fig:product}. The
$\starG$ model over $G_1 \times G_2$ achieves perfect prediction
(\Rsq{} $= 1.000 \pm 0.000$) with just 186 parameters. The single-factor models
capture at most 23\%: $G_2$ alone (\Rsq{} $= 0.229$) and $G_1$ alone
(\Rsq{} $= 0.155$). The 2D frequency map (Figure~\ref{fig:product}b) visualizes
why: the coupled frequency cells (highlighted in red) carry 87\% of the target
energy and are resolved only by $F_{\bbZ_6} \otimes F_{\bbZ_4}$; the individual
transforms see only the axis-aligned marginals. The cyclic approximation
$\bbZ_{24}$ (treating the same 24-dimensional group dimension as a single cyclic
group) reaches \Rsq{} $= 0.986$, close but not exact, because it cannot
distinguish the tensor-product structure of the irreps. The ablation cascade
(Figure~\ref{fig:ablation}) confirms a strict performance hierarchy: product group
$>$ wrong cyclic $>$ single factor $>$ no symmetry, each step removing algebraic
information and reducing performance monotonically. On this task the implementation realizes
Theorem~\ref{thm:product_ring} exactly: only the exact product-group specification
achieves exact recovery, and every coarsening of it degrades performance. Because
the target is constructed to contain coupled product-group frequencies, the
experiment validates that the algebra recovers the intended structure when it is
present, not that such coupled structure is ubiquitous in real materials.

\begin{table}[h]
\centering
\caption{Product group $\bbZ_6 \times \bbZ_4$ (1{,}000 molecules, 3 seeds).}
\label{tab:product}
\begin{tabular}{@{}lcc@{}}
\toprule
\textbf{Method} & \textbf{Test \Rsq} & \textbf{Params} \\
\midrule
$\starG$ over $G_1 \times G_2$ + Ridge & $\mathbf{1.000 \pm 0.000}$ & \textbf{186} \\
$\bbZ_{24}$ cyclic + Ridge             & $0.986 \pm 0.002$           & 157 \\
$\starG$ over $G_1 \times G_2$ + MLP   & $0.826 \pm 0.099$           & 9{,}473 \\
Standard MLP                           & $0.488 \pm 0.116$           & 4{,}481 \\
Invariant MLP                          & $0.324 \pm 0.267$           & 6{,}785 \\
Augmented MLP                          & $0.250 \pm 0.239$           & 4{,}481 \\
$G_2$ only ($\bbZ_4$) + Ridge          & $0.229 \pm 0.191$           & 42 \\
$G_1$ only ($\bbZ_6$) + Ridge          & $0.155 \pm 0.244$           & 55 \\
\bottomrule
\end{tabular}
\end{table}

\begin{figure}[!ht]
\centering
\includegraphics[width=\textwidth]{fig4_product.pdf}
\caption{\textbf{Product group $\bbZ_6 \times \bbZ_4$: compositional advantage.}
(\textbf{a})~Eight-method comparison. The product group achieves $R^2 = 1.000$;
each factor alone captures $\leq 23$\%.
(\textbf{b})~2D frequency map: coupled cells (red borders) carry 87\% of target
energy and are resolved only by $F_{\bbZ_6} \otimes F_{\bbZ_4}$.}
\label{fig:product}
\end{figure}

\begin{figure}[!ht]
\centering
\includegraphics[width=0.65\textwidth]{fig_ablation.pdf}
\caption{\textbf{Ablation cascade.} Progressively removing symmetry components
reveals a strict performance hierarchy: product group $\to$ wrong cyclic
approximation $\to$ single factor $\to$ no symmetry.}
\label{fig:ablation}
\end{figure}

\paragraph{Non-abelian composition and data efficiency against a rotation-only
network.} The $\bbZ_6 \times \bbZ_4$ task above uses two abelian factors.
Composition matters most when one factor is non-abelian, so we built a second task
whose target is a joint invariant under the chiral octahedral group $O$
($|O|\!=\!24$, acting on 3D point clouds by rotation) and a global cyclic shift
$\bbZ_C$ ($C\!=\!4$) of a discrete per-point label. The target is a fixed
$\star_{O \times \bbZ_C}$ power-spectrum readout on the joint-nontrivial product
irreps, exactly invariant under both symmetries and dependent on neither alone
(Methods~\S\ref{sec:composed_methods}). Since it is generated by a product-group
readout, the experiment measures data efficiency and inductive bias rather than
expressive capacity. We compare four models over a sample-size sweep (3 seeds;
Table~\ref{tab:composed}), the three symmetry-aware models matched near
$\sim\!50$k parameters and the unconstrained MLP baseline deliberately given
$1.8$M: the product-group
$\star_{O \times \bbZ_C}$ readout, which carries both symmetries by construction; a
single-symmetry $\star_{O}$ readout built from the identical spatial machinery with
the label folded into ordinary channels, exactly rotation-invariant but required to
learn the color-shift invariance from data; a rotation-equivariant SchNet given the
label as an input feature; and a plain MLP. Validation confirms the intended
structure to floating-point precision: the product readout is invariant to both
symmetries ($\leq 5\times 10^{-5}$), whereas $\star_{O}$ is rotation-invariant but
not color-shift-invariant.

The controlled comparison is between the two $\star_G$ readouts, which share a
hypothesis class and differ only in whether the second symmetry is built in. Their
gap is a clean data-efficiency signature: negligible once both saturate ($+0.001$
at $3{,}000$ samples) and growing as data becomes scarce ($+0.043$ at $300$,
$+0.225$ at $100$), where the built-in color-shift invariance spares the product
model from spending data to discover it. The rotation-only baselines sit far below
across the whole sweep: SchNet rises only from $0.03$ to $0.31$ and the MLP barely
leaves zero, so a network that must absorb the second symmetry by rote does not
catch up even at $3{,}000$ samples. We read the SchNet gap cautiously, since it
also reflects SchNet approximating the $\star_G$ environment kernel through
distance messages and carries large seed variance; the matched $\star_{O}$
comparison is the controlled one. The experiment validates
Theorem~\ref{thm:product_ring} for a non-abelian factor and locates the
compositional advantage in the data-scarce regime.

\begin{table}[h]
\centering
\caption{Composed-symmetry task, chiral octahedral $O$ $\times$ color shift
$\bbZ_4$: test \Rsq{} (mean $\pm$ std, 3 seeds) over a sample-size sweep; the
three symmetry-aware models are matched near $\sim\!50$k parameters, while the
unconstrained MLP baseline has $1.8$M. The product-group readout carries both
symmetries by construction; $\star_{O}$-only is exactly rotation-invariant but must
learn color-shift invariance from data; SchNet is rotation-equivariant with the
label as an input feature. The target is a fixed product-group readout, so the
comparison measures data efficiency, not expressive capacity.}
\label{tab:composed}
\begin{tabular}{@{}lcccc@{}}
\toprule
\textbf{$n_{\mathrm{train}}$} & \textbf{$\star_{O \times \bbZ_4}$} &
\textbf{$\star_{O}$ only} & \textbf{SchNet} & \textbf{MLP} \\
 & \textbf{(59k)} & \textbf{(43k)} & \textbf{(63k)} & \textbf{(1.8M)} \\
\midrule
$100$     & $\mathbf{0.840 \pm 0.043}$ & $0.615 \pm 0.099$ & $0.033 \pm 0.267$ & $-0.003 \pm 0.029$ \\
$300$     & $\mathbf{0.993 \pm 0.000}$ & $0.950 \pm 0.018$ & $0.212 \pm 0.232$ & $0.012 \pm 0.044$ \\
$1{,}000$ & $\mathbf{0.998 \pm 0.001}$ & $0.989 \pm 0.013$ & $0.257 \pm 0.251$ & $0.031 \pm 0.083$ \\
$3{,}000$ & $\mathbf{0.999 \pm 0.000}$ & $0.998 \pm 0.001$ & $0.312 \pm 0.169$ & $0.136 \pm 0.153$ \\
\bottomrule
\end{tabular}
\end{table}

% Experiment 4 (symmetry and factorization selection) relocated to the
% Supplementary Information (Section: Group and Factorization Selection).

\subsection{Permutation symmetry and simultaneous composition}
\label{sec:permutation}

A recurring claim for $\starG$ is that a new symmetry costs only a change of the
group argument (the cost is computational, set by the group order, rather than an
architectural redesign). We make this concrete for the
symmetric group $S_n$, the permutation symmetry central to molecular learning,
and for the simultaneous handling of two symmetries at once. The $\starG[S_n]$
algebra is built with exactly the machinery used for the cyclic and octahedral
cases: the irreducible representations of $S_n$ (Young's orthogonal form over
standard Young tableaux) assemble the unitary group-Fourier matrix $F_{S_n}$, and
nothing else changes. Three facts, all verified to machine precision for $S_4$:
(i)~$\starG[S_n]$ convolution block-diagonalises under $F_{S_n}$, so the product
is one matrix multiplication per Specht-module irrep (residual
$1.8\times10^{-14}$); (ii)~the decomposition is permutation-equivariant: under
left translation by any $\tau \in S_n$ the per-irrep singular values are
invariant ($2.7\times10^{-15}$); and (iii)~by the product-group theorem
(Theorem~\ref{thm:product_ring}), $F_{\bbZ_k \times S_n} = F_{\bbZ_k} \otimes
F_{S_n}$ handles a combined cyclic-and-permutation symmetry with one transform.
On a signal whose structure is genuinely coupled across both factors, the
product transform compacts it completely (top-$10\%$ energy fraction $1.000$)
while the cyclic factor alone ($0.565$) and the permutation factor alone
($0.445$) each leave most of the energy spread out. A single algebra thus
captures a joint symmetry that neither factor recovers, and adding the
permutation factor required no new architecture, only the Kronecker product of
Theorem~\ref{thm:product_ring}. These are algebraic and numerical demonstrations on a
finite permutation group: they establish that permutation-equivariance and
composition are intrinsic to the algebra, not that $\starG$ matches graph equivariant
networks for full three-dimensional, permutation-invariant molecular prediction,
which is an architectural question outside this algebraic claim.

\subsection{Scope: pooled predictive accuracy}
\label{sec:enn_comparison}

The results above are capabilities, and each is exact: enforcement, diagnosis,
composition, optimality. Pooled predictive accuracy on standard molecular
benchmarks is deliberately not among them, and we state the scope plainly. On
the full QM9 benchmark (130{,}831 molecules; a deliberate stress test whose
symmetry is only approximate and whose discriminating signal is largely
bond-topological), graph-based equivariant networks lead on pooled $R^2$
(SchNet $0.996$, MACE $0.985$ on the HOMO--LUMO gap, at $10^5$--$10^6$
parameters), and on the identical $48$-row molecule-level feature tensor a
well-trained standard MLP ($0.687$) exceeds both $\starG$ variants
($\starG$-SVD~+~Ridge $0.481$ at $144$ parameters). A within-isomer audit shows
that most of the ENN margin is an input-representation gap rather than an
algorithmic one (every molecule-summary method caps near within-isomer
$R^2 \approx 0.35$ versus $\approx 0.96$--$0.99$ for the graph networks, which
consume the bond topology that distinguishes isomers), and a data-scarce sweep
locates the regime boundary explicitly: the closed-form $\starG$-SVD leads
every trained baseline through $316$ training molecules, and the unconstrained
MLPs overtake between $316$ and $1{,}000$, beyond which the sample-complexity
prior should not be extrapolated; both effects are regime- and
input-specific rather than a ranking of methods. The complete accuracy tables,
learning curves, and audit protocol are in SI~Section~\ref{sec:qm9_accuracy};
the contribution of this paper is structural and diagnostic, not an
improvement in predictive accuracy, and where chemistry-aware pooled accuracy
on a single target with generous data and compute is the goal, graph ENNs
remain the right tool.

One accuracy result does belong to the exactness axis of
Section~\ref{sec:compounding}. The augmented MLP, which attempts to
\emph{learn} the invariance from $|G|$-fold orbit-augmented training data and
is the closest non-equivariant analogue of $\starG$, collapses at full QM9
scale ($R^2 \approx 0.002$ on gap, $\alpha$, and $\mu$), and a capacity sweep
to $575{,}489$ parameters lifts it only to $R^2 \approx 0.03$
(SI~\S\ref{sec:aug_capacity}): two orders of magnitude below a plain MLP on
identical features. Orbit augmentation sits on the approximate side of the
exact/approximate line ($\varepsilon \sim 10^{-2}$), and at scale it fails to
deliver even approximate invariance; the algebra builds the invariance in
exactly, at zero training cost.

\subsection{Symmetry in the representation, not the architecture}
\label{sec:representation}

The accuracy comparisons of Section~\ref{sec:enn_comparison} place $\starG$ at
an input disadvantage: it consumes a molecule-level summary while the ENNs
consume the full atomic graph, and the within-isomer audit
(SI~Section~\ref{sec:qm9_accuracy}) showed that every tested model on that
summary plateaus near $R^2_{\mathrm{within}} \approx 0.35$. That plateau is a
property of the featurization, not of the algebra. To test the algebra on even ground we build
a \emph{per-atom} $\starG$ descriptor: for each atom the local environment
(neighbours within a cutoff, resolved by element) is sampled over the $24$
octahedral frames and reduced to its per-irrep power spectrum, a
$\starG$-native readout that is exactly rotation-invariant (to $10^{-16}$), and
the per-atom invariants are pooled over atoms. The symmetry lives entirely in
this representation; the predictor is a symmetry-\emph{blind} gradient-boosted
regressor with no equivariant architecture.

On full QM9 this reaches within-isomer $R^2 = 0.785 \pm 0.003$ and pooled
$R^2 = 0.895 \pm 0.001$ (mean over 3 seeds) on the HOMO--LUMO gap, versus the
$\approx 0.35$ plateau of all tested molecule-level-summary models and $\approx 0.99$
for the graph-based ENNs; even a linear ridge regressor on the same descriptor
clears it ($0.46$). The equivariance benefit therefore does not require an equivariant
network: it can be carried by the data representation and handed to an
arbitrary learner. In this role $\starG$ is a composable equivariant
featuriser, one that decouples symmetry from architecture and composes several
symmetries by the Kronecker rule $F_{G_1}\otimes F_{G_2}$ with no redesign.
The remaining gap to the graph ENNs is honest: message passing still resolves
bond topology that a fixed local descriptor does not fully capture, and closing
it is a matter of richer descriptors rather than of the algebra. No algebraic
rigidity prevents combining the two: the descriptor is already computed per
atom, so its per-irrep spectra can serve directly as node features inside any
message-passing network, and message functions could themselves be built from
$\starG$ products, making each layer's equivariance exact by construction. We
do not pursue that hybrid here because it would re-entangle the two axes this
section deliberately separates (representation versus architecture), but it is
an engineering question, not a mathematical obstruction.
Figure~\ref{fig:paradigm} summarizes the two routes to equivariance this
result separates: the symmetry can live in the layers of the model, or in the
algebra of the data.

\begin{figure}[!ht]
\centering
\includegraphics[width=0.85\textwidth]{fig11_paradigm.png}
\caption{\textbf{Two routes to equivariance.}
\textit{Left (architectural equivariance):} the symmetry constraint is built
into the layers, so the set of symmetries is fixed when the architecture is
designed.
\textit{Right (algebraic equivariance):} the same algebra handles any finite
or band-limited compact group; composing symmetries requires only specifying
$G_1 \times G_2$ in the Fourier transform.}
\label{fig:paradigm}
\end{figure}

\subsection{Extension to Compact Groups: the Band-Limited \texorpdfstring{$\starG$}{starG} Algebra}
\label{sec:compact}

The construction above used a finite group only to make the convolution tensor
$\calT$ finite-dimensional; the algebraic content is supplied entirely by the
Peter--Weyl theorem, which holds for \emph{every compact group}. For compact $G$
with normalized Haar measure, $L^2(G)$ decomposes as a Hilbert direct sum over
the countable set $\widehat{G}$ of irreducible unitary representations, the group
Fourier transform sends $f \mapsto (\widehat{f}(\rho))_{\rho \in \widehat G}$ with
$\widehat{f}(\rho) = \int_G f(g)\,\rho(g)\,dg \in \mathbb{C}^{d_\rho \times d_\rho}$,
convolution becomes block-diagonal matrix multiplication exactly as in
Eq.~(\ref{eq:spectral}), and the Plancherel identity reads $\|f\|^2 = \sum_{\rho}
d_\rho \,\|\widehat{f}(\rho)\|_{\mathrm{HS}}^2$. For $G=\mathrm{SO}(3)$ the blocks
are the Wigner $D^{(l)}$ matrices with $d_\rho = 2l+1$, $l=0,1,2,\dots$: the
angular-momentum channels of physics are \emph{literally} the irrep blocks of the
$\starG$ algebra. The $\starG$ product, transpose, norm, and equivariance
therefore carry over verbatim; the only change is that the irrep sum is now
countably infinite. The $\starG$-SVD likewise generalizes to the simultaneous SVD
of the irrep blocks, and optimal low-rank approximation reduces block by block to
the matrix Eckart--Young theorem, with the Plancherel weights $d_\rho$ selecting
which singular values to retain across blocks; for Hilbert--Schmidt operators the
weighted singular-value sequence is summable and the optimum is attained.

The case that matters in practice is finite-dimensional and needs no such
hypothesis. Call a signal \emph{band-limited} to a finite set $S \subset
\widehat{G}$ if $\widehat{f}(\rho)=0$ for $\rho \notin S$ (for $\mathrm{SO}(3)$,
$S = \{l \le l_{\max}\}$).

\begin{theorem}[Band-limited reduction]\label{thm:bandlimited}
Let $G$ be a compact group and $S \subset \widehat{G}$ finite. The space of
signals band-limited to $S$ is closed under the $\starG$ product and carries a
$\starG$ algebra isomorphic to the finite-dimensional $\starG$ algebra over the
irrep set $S$. Consequently every finite-group result of this paper, including
the $\starG$-SVD Eckart--Young optimality (Theorem~\ref{thm:eckart_young}),
holds verbatim for band-limited signals on any compact group.
\end{theorem}

The proof is immediate: band-limiting truncates the Peter--Weyl sum to the
finitely many blocks in $S$, and block-diagonal multiplication preserves the
truncation, so the band-limited algebra \emph{is} the finite $\starG$ algebra
over $S$. Its reduction target is therefore the same finite-dimensional algebra
the Lean~4 development formalizes over the irrep set $S$; the Peter--Weyl
truncation that effects the reduction is a pen-and-paper argument and is not
itself separately machine-checked.

We validate this reduction numerically for $G=\mathrm{SO}(3)$ band-limited to
$l \le 3$, applying the \emph{same} $\starG$-product and $\starG$-SVD code used
for the finite groups directly to the Wigner-$D$ Fourier blocks
(\texttt{bandlimited\_so3\_experiment.py}). Every finite-group guarantee holds
to machine precision: per-block SVD agreement $2\times10^{-14}$, Eckart--Young
closed-form tail-sum match $2\times10^{-12}$, exact
$\mathrm{SO}(3)$-equivariance of the per-irrep power and singular values to
$2\times10^{-15}$, and a compression advantage on genuinely
$\mathrm{SO}(3)$-structured operators that vanishes on unstructured controls.
Band-limiting does not reintroduce the approximate side of
Section~\ref{sec:compounding}'s dichotomy: truncation is an equivariant
projection onto a subalgebra closed under the product, so iterating a
band-limited operator remains exactly equivariant at every depth, and the
$O(L^{-s/2})$ tail incurred on merely smooth signals
(SI~Proposition~\ref{prop:bandlimit_error}) is a resolution error confined to
the symmetry-allowed channels, chosen in practice from the observable spectral
tail of the data (SI~Remark~\ref{rem:truncation_exactness}).

The same mechanism reaches the crystallographic space groups, the actual symmetry
of matter. A space group is an extension of a point group by the lattice
translations, and because the translations are dual to the compact Brillouin zone,
imposing periodic boundary conditions replaces the space group by a finite group;
the calculus, its Eckart--Young optimality, and its selection-rule enforcement then
apply verbatim, resolved per Brillouin-zone point through the little-group
irreducible representations (SI~Section~\ref{sec:spacegroups}). This is a reduction
to the finite theory rather than new analysis: it extends the point-group treatment
of crystals used above to all $230$ space groups (and, with a projective twist, the
non-symmorphic ones), and its $\Gamma$-point instance is exactly the point-group
selection-rule projection applied to cubic crystals in
Section~\ref{sec:crystals}. Antiunitary time reversal (magnetic and grey groups)
and band topology are out of scope, as detailed there. Nor does the reach stop
at crystals: the Mackey construction extends the same machinery through
compact little-group fibers to the noncompact Euclidean and Poincar\'e
groups, validated numerically on the Poincar\'e mass shell and on an
$\mathrm{SU}(2)$ lattice-gauge plaquette
(SI~Section~\ref{sec:wignermackey}; Discussion).

This has a direct bearing on equivariant neural networks. Every
$\mathrm{SO}(3)$-equivariant architecture used in practice (e3nn, NequIP,
MACE) represents features in a finite set of angular-momentum channels $l \le
l_{\max}$ (typically $l_{\max}\in\{1,2,3\}$); their feature spaces are precisely
band-limited $\star_{\mathrm{SO}(3)}$ modules. By Theorem~\ref{thm:bandlimited}
the $\starG$-SVD furnishes the \emph{provably optimal} equivariant low-rank
operation on those modules in closed form (Corollary~\ref{cor:unification}), whereas the network fits a map within that same equivariant class by gradient descent, with no optimality characterization.
Architectural and algebraic equivariance thus act on the same band-limited
representation space; what the $\starG$ algebra adds is the Eckart--Young
guarantee and the closed-form per-irrep decomposition
(Table~\ref{tab:wigner}) that standard ENN training does not expose. The finite octahedral
computation of our Wigner--Eckart experiment is one instance of this band-limited
algebra; a continuous $\star_{\mathrm{SO}(3)}$ computation at $l_{\max}=2$ is
another.

This can be stated precisely. Let the feature space be band-limited to a finite
$S \subset \widehat{\mathrm{SO}(3)}$.

\begin{corollary}[Equivariant linear layers are $\starG$ operations]\label{cor:unification}
By Schur's lemma, every $\mathrm{SO}(3)$-equivariant linear map on this space is
block-diagonal in the irrep basis, i.e.\ a $\star_{\mathrm{SO}(3)}$ operation; and
among such maps of $\starG$-rank at most $k$, the $\starG$-SVD truncation
is Frobenius-optimal (Theorem~\ref{thm:eckart_young}).
\end{corollary}

\noindent The \emph{linear} equivariant layers of $\mathrm{SO}(3)$-equivariant
networks (e3nn, NequIP, MACE), their self-interaction and channel-mixing maps, are
therefore particular $\star_{\mathrm{SO}(3)}$ operations, fit by gradient descent;
the Clebsch--Gordan tensor-product layers are \emph{bilinear} couplings between
irreps and lie outside this operation class. The block-diagonality itself is
classical (Schur), and the correspondence between equivariant linear maps and
convolutions is due to Cohen, Geiger, and
Weiler~\citep{cohen2019general}; what the algebra contributes is the
\emph{optimal} member of that linear class in closed form. This is a structural
relationship, not a claim that a trained network reduces to a single $\starG$
layer: such networks compose many layers with nonlinear gates. The optimum is
simply present in the operation class and left unused.

\section{Discussion}

The through-line of this work is exactness. Once symmetry is made an algebraic
property of the data, equivariance holds identically rather than approximately,
and Section~\ref{sec:compounding} showed why that distinction is qualitative
rather than quantitative: approximate equivariance compounds under
iteration and composition to $O(1)$ symmetry violation, in the worst case and
generically in stable norm-preserving computation, while exact
equivariance holds at unbounded depth. On that exact footing, one construction
serves two ends usually pursued apart: optimal
equivariant compression and the read-out of physical symmetry structure. A single
consequence runs through every experiment: because the symmetry lives in the
algebra of the data rather than in the layers of a model, the equivariant structure
can be handed to an arbitrary, symmetry-blind learner (a gradient-boosted regressor
breaks a featurization ceiling that every molecule-summary model hits; a projection
repairs a trained graph network's forbidden-channel leakage), and the guarantees it
carries, optimal low-rank compression, exact selection rules, and per-irrep
separation, are closed-form rather than learned. The most
vivid instance of the second is that the $\starG$ tensor algebra, applied
over a physically meaningful group (the octahedral subgroup of SO(3)), surfaces,
directly from molecular geometry data, octahedral finite-subgroup signatures
consistent with the angular
momentum selection rules of the Wigner--Eckart theorem. The T$_1$/A$_1$ predictive power ratio separates the
vector observable (dipole magnitude, $\approx 2.1$) from the scalar observables
($\leq 0.2$; Table~\ref{tab:wigner}) by an order of magnitude,
and the isotropic polarizability's near-zero T$_1$ dependence confirms the
representation-theoretic absence of $l\!=\!1$ content in symmetric rank-2 tensors.
These patterns emerge without any quantum-mechanical input, demonstrating that the
$\starG$ framework functions as a \emph{spectroscope for physical symmetry}: it
decomposes empirical predictions into irreducible representation channels that
expose the group-representation structure of the observables.

This spectroscopic capability rests on the theoretical results established above: (i)~the Peter--Weyl spectral decomposition of the convolution tensor,
which expresses group structure through the sparse core tensor $\calC$ (equation
(\ref{eq:spectral})); (ii)~the Eckart--Young optimality guarantee for the
$\starG$-SVD (Theorem~\ref{thm:eckart_young}), instantiating the transform-domain
t-SVD optimality~\citep{kilmer2021tensor} at the group-Fourier transform and
extending it to non-abelian groups; (iii)~the product group ring
isomorphism (Theorem~\ref{thm:product_ring}), which composes multiple symmetries
seamlessly (\Rsq{} $= 1.000$ for $\bbZ_6 \times \bbZ_4$ vs.\ $\leq 0.23$ for
single factors). Because the group is a parameter, it can in principle be selected
by explicit model comparison over a candidate library; we do not develop this as a
contribution. Automatic symmetry discovery is a separate, active line of work, for
continuous groups~\citep{dehmamy2021lieconv,yang2023gan}, for symmetry breaking from
an assumed group~\citep{wang2023relaxed}, and for finite
groups~\citep{karjol2023discrete}.

Underlying these is a canonicity result (SI~Theorem~\ref{thm:necessity}): the
$\starG$ algebra is not one arbitrary member of the transform-based
($\starM$/t-product) family but the unique one compatible with the group action.
Requiring equivariance forces the transform, and with it the algebra's entire
matrix-mimetic toolbox (the product, and the QR, symmetric-eigendecomposition,
least-squares, and SVD factorizations built on it), to be determined by $G$ rather
than chosen. The low-rank optimum inherits this canonicity: it is intrinsic to the
data and its symmetry (SI~Corollary~\ref{cor:intrinsic}). In this sense the
construction is a statement about $G$ rather than a design for it.

The practical implications follow the same honest line as the experiments. Where
data is scarcest, the closed-form $\starG$-SVD ridge regressor is the most
sample-efficient method on QM9: at $100$ molecules it leads the unconstrained
networks with two orders of magnitude fewer parameters and no training
(SI~Section~\ref{sec:qm9_accuracy}), a lead that narrows and reverses as data
grows. A compact,
exactly equivariant, closed-form representation is thus most valuable in exactly
the resource-constrained scientific settings where interpretability and parameter
efficiency matter, even if scale eventually wins on aggregate accuracy.

\subsection*{Scope and limitations}

\paragraph{Realizability of the low-rank operator.} It is worth stating precisely
what the $\starG$-SVD certificate does and does not guarantee. The discarded
singular-tube mass is the exact Frobenius error of the rank-$k$ truncation, so it
certifies \emph{approximation quality}; it does not certify that the retained
equivariant subspace admits a \emph{localized} symmetric basis. For finite groups
this distinction is immaterial, since the transform domain is a finite direct sum
of irrep blocks. It becomes substantive in the band-limited continuous-group
setting (Section~\ref{sec:compact}), where the group-Fourier fibers vary over a
continuous parameter, for a crystalline system the Brillouin zone. There the
singular values are gauge-invariant scalars, whereas any topological content
resides in the winding of the singular vectors; consequently two operators with
identical singular-value fields, and hence an identical certificate at every
rank, can carry different Chern number and therefore differ in whether their
retained subspace admits an exponentially-localized symmetric (Wannier) basis.
The certificate is, by construction, blind to this obstruction. More structurally,
wherever $\starG$ is exact its group algebra is semisimple, so the obstruction
vanishes identically in the finite-group regime and can arise only in the
continuous-group extension. We regard this as
a scope statement rather than a defect: $\starG$ delivers provably optimal,
symmetry-preserving approximation together with its closed-form error, but whether
the compressed operator admits a physically-localizable realization is a separate
representation-theoretic question, to be settled with a symmetry-indicator or
curvature diagnostic rather than the approximation certificate alone.

\paragraph{The per-irrep readout is a second-order invariant.} The per-irrep power
spectrum $\|\hat f(\rho)\|^2$ underlying our interpretability analysis, and the
singular-tube energies defining the $\starG$-SVD certificate, are \emph{second-order}
(power-spectral) invariants. They are therefore invariant under a strictly larger
group than $G$: any per-irrep unitary $U$ with $|\hat U(\rho)|=1$ acts as an isometry
of the $\starG$ Frobenius norm, so the readout cannot in general separate distinct
$G$-orbits (it is orbit-\emph{incomplete}). Restoring orbit separation requires a
third-order (bispectral) invariant, as established for compact and finite groups by
the bispectrum-completeness literature~\citep{sanborn2023bispectral,mataigne2024selective}.
Our contribution is not orbit completeness but the closed-form, per-irrep
\emph{attribution} of predictive power to angular-momentum channels; where orbit
discrimination is required, the leading singular vectors (which carry the discarded
phase) or an explicit bispectral invariant should be used.

\paragraph{The compression advantage is a bounded, group-determined constant.} The
gap between $\starG$-SVD rank and matrix-SVD rank on octahedral and QM7-X data (exact
recovery at rank~3 versus rank~6--11) should not be read as an unbounded algebraic
separation. It is controlled by the block structure of the group-Fourier transform:
the achievable rank reduction is bounded by the largest irreducible-representation
dimension $d_{\max}$, is exactly~$1$ for abelian (e.g.\ cyclic) groups where every
irrep is one-dimensional, and can be~$1$ even for a non-abelian group when the data
occupies a single high-dimensional irrep. The effect is the classical
Peter--Weyl/Schur block-diagonalization~\citep{serre1977linear} made quantitative for
low-rank approximation; its value is the closed-form, per-block optimality and error,
not a large or unbounded advantage over matrix methods.

\paragraph{The compact setting is the natural maximal domain.} The restriction
to finite and compact groups is not a modeling convenience; it marks, in a
precise sense, the boundary of the theory. The $\starG$-SVD needs two
ingredients: a decomposition of the group algebra into irreducible blocks, and
a Plancherel weight with which to compare singular values across blocks. Both
survive exactly one step further. For separable unimodular type~I locally
compact groups (the Euclidean motion groups, for instance) the Plancherel
theorem replaces the sum over $\widehat{G}$ by a direct integral, and the
per-irrep Eckart--Young statement becomes a measurable field of blockwise
SVDs~\citep{folland2016course}; the measurable-selection analysis this
requires is available in the literature~\citep{franco2024measurability}, so
the extension is corollary-level rather than new mathematics. For the
semidirect-product groups that carry the physics, the Euclidean motion groups
and the Poincar\'e group, this direct integral is moreover rarely the object one
must compute. The Mackey construction realizes each irreducible as a
\emph{compact} little-group representation over a continuous orbit of momenta, so
the non-compactness lives in the base while the fiber stays finite-dimensional;
linearizing a covariant system about a background hands one exactly that compact
little group, and the boost covariance is then discharged by a symbolic
conjugation along the orbit, under which the singular values are invariant,
rather than by any numerical integration. Carried out on the universal cover
$\bbR^{1,3} \rtimes \mathrm{SL}(2,\mathbb{C})$, whose compact little groups are
$\mathrm{SU}(2)$ and the double cover of $\mathrm{SO}(2)$, the finite
$\starG$-SVD thus reaches
the entire physical one-particle spectrum, massive and massless-helicity,
integer and half-integer spin alike, with continuous-spin representations
(never observed in nature) and genuinely boost-covariant or bootstrap
observables reserved for
the direct integral; we develop this Wigner--Mackey compact-fiber reduction, and delimit
its residual, in Supplementary Information Section~\ref{sec:wignermackey}. Beyond
type~I the construction fails in principle rather than in practice: by Glimm's
theorem~\citep{glimm1961type}, a non-type-I group, a class that includes free
groups and most infinite discrete groups, admits no standard Borel dual, and
direct-integral decompositions into irreducibles, though they exist, are no
longer essentially unique, so ``per irrep'' ceases
to be well-defined; what remains available there, low-rank approximation in the
group von Neumann algebra through generalized $s$-numbers, is classical
functional analysis~\citep{fack1986generalized} and carries no irrep
structure. Genuinely infinite-dimensional groups such as loop groups and
$\mathrm{Diff}(S^1)$ fail earlier still: they are not locally compact and have
no Haar measure, and the von Neumann algebras attached to their
positive-energy representations are type~III, with no trace and hence no
Frobenius norm to optimize. In the other direction, ``beyond semisimple'' is
vacuous where data lives: over $\mathbb{R}$ or $\mathbb{C}$, Maschke's theorem
makes every finite (and compact) group algebra semisimple, and non-semisimple
group algebras arise only in positive characteristic, where no inner product,
and hence no SVD, exists. The band-limited compact setting of
Theorem~\ref{thm:bandlimited} is therefore essentially the largest arena in
which a per-irrep Eckart--Young theorem with a canonical block decomposition
can be stated.

\paragraph{Choice of benchmark.} QM9 is a demanding setting for a symmetry
algebra, and we used it deliberately as a stress test rather than a showcase:
molecular conformers are only \emph{approximately} symmetric, and the
property-discriminating signal is largely bond-topological, which favours graph
message passing over any symmetry decomposition. The framework's home turf is data
whose symmetry is exact and discriminative, which is why the crystalline
selection-rule results (Section~\ref{sec:crystals}), together with
the ENN-class within-isomer discrimination reached even on the unfavourable QM9
turf (Section~\ref{sec:representation}), are where $\starG$ shows its clearest
advantage.

\paragraph{Angular resolution of finite subgroups.} The $\starG$ framework
extends from finite groups to the band-limited setting of any
compact group through
the band-limited reduction of Section~\ref{sec:compact}, so continuous
symmetries such as $\mathrm{SO}(3)$ are covered once a maximum angular-momentum
channel $l_{\max}$ is fixed. The chiral octahedral group $O$ ($|G|=24$) used in
our Wigner--Eckart experiment is instead a discrete \emph{subgroup} of
$\mathrm{SO}(3)$; the qualitative selection-
rule signatures (scalars $A_1$-dominated, dipoles $T_1$-dominated,
polarisability $T_1$-insensitive) survive this discretisation cleanly
because they reflect the irrep content of each tensor rank rather than
the angular resolution. The quantitative accuracy ceiling, by
contrast, is set by the finite-group resolution: a molecule that is
well-resolved by $\mathrm{SO}(3)$ but not by $O$ has its angular
information aliased onto the 24-element orbit, and per-irrep $R^2$
saturates accordingly. Larger polyhedral groups lift this ceiling, and we
confirm it directly: with a machine-validated icosahedral $\starG$ algebra
($I$, $|G|=60$; five irreps $A,T_1,T_2,G,H$ with $\sum_\rho d_\rho^2 = 60$),
random signals band-limited to angular degree $\ell=4$ are aliased by the
octahedral 24-frame orbit (test $R^2 = 0.91$) but resolved exactly by the
icosahedral 60-frame orbit (test $R^2 = 1.00$), consistent with $I$ carrying a
dimension-5 irrep where $O$ tops out at dimension~3
(SI~Section~\ref{sec:ico_resolution}). A natural next step is predicting full
tensor-valued properties (not just scalar invariants), which would enable
complete Wigner--Eckart decomposition including $l\!=\!2$ selection rules for
the quadrupole moment.

\subsection*{Broader Impact}

The framework opens a path from Eckart--Young (optimal low-rank tensor
approximation preserving group structure) to Wigner--Eckart (angular momentum
selection rules for physical observables) through a single algebraic construction,
closing a circle between two theorems that, fittingly, share a common author in
Carl Eckart. By changing only the group $G$, the same machinery that achieves
parameter-efficient molecular property prediction in the data-scarce regime, without
matching graph networks on pooled accuracy, can decompose physical observables
by angular momentum channel and compose multiple symmetries. We anticipate applications in
materials science (crystallographic point groups), particle physics (gauge
symmetries), and drug discovery (molecular chirality), where the ability to
simultaneously predict properties and reveal their symmetry content could
accelerate scientific understanding. In lattice gauge theory in particular the
compact gauge group is truncated per link to a finite band of irreducibles,
precisely the band-limited compact setting in which the $\starG$-SVD and its
selection rules apply verbatim.

\subsection*{Formal Verification}

The $\starG$-specific reductions and algebraic identities have been machine-checked in the
Lean~4 proof assistant~\citep{demoura2021lean4} with
Mathlib~\citep{mathlib2020}, with zero unresolved proof obligations
(\texttt{sorry}), conditional on a short list of standard linear-algebra and
Fourier-analysis facts declared as axioms (enumerated below and in
SI~Section~\ref{sec:lean}). The core
algebraic theorems are proved outright, depending only on Lean's
foundational axioms: $\starG$ associativity, identity, distributivity,
transpose reversal, left equivariance (right equivariance for abelian $G$), the
product-group factorization of the convolution tensor and the Kronecker
construction of product-group irreps (Theorem~\ref{thm:product_ring}), and
invariance of the per-irrep Fourier power under the group action. Beyond the SVD, the matrix-mimetic character of the
algebra is itself machine-checked: the $\starG$-QR factorization (with
$\starG$-orthonormality $Q^{\!H}\!\star_G Q=\mathcal I$), the symmetric
$\starG$ eigendecomposition, and $\starG$ least-squares are each verified
by the same block reduction, resting on the axiom-free bridge lemma that
the group-Fourier transform sends the $\starG$ transpose to the transpose
of each irrep block. The Eckart--Young optimality result
(Theorem~\ref{thm:eckart_young}) rests on
classical facts the formalisation declares explicitly rather than
reproving. The real development carries exactly \emph{six} declared axioms:
the Plancherel identity for finite groups, the injectivity and surjectivity
of the group-Fourier transform (Peter--Weyl theory), two elementary lemmas
behind the matrix Eckart--Young theorem (capture of a low-rank matrix by an
orthogonal projection, and Ky~Fan's 1949 max-trace inequality), from which
the matrix Eckart--Young theorem is itself derived rather than assumed, and
the existence of the matrix QR factorization, pending its addition to
Mathlib. The symmetric
eigendecomposition and least-squares theorems carry no matrix axiom, being
proved from
Mathlib's spectral theorem and the normal equations respectively (their
$\starG$-orthogonality reductions rest on Fourier injectivity), and a
concrete axiom-free irreducible-representation system is constructed at which
the whole optimality suite is instantiated, so the theorems are demonstrably
non-vacuous. A parallel complex-typed module adds three complex-local
analogues of the Fourier axioms, used only to build cyclic-group (DFT)
witnesses that the real development cannot express; no real-development
theorem depends on them, bringing the repository's total to nine
\texttt{axiom} declarations. The octahedral Wigner--Eckart selection rules are not
asserted but \emph{derived} from the character table of the octahedral
group $O \cong S_4$, computed over the group itself: the irreducible
characters are proved orthonormal, and the polarisability rule
$\mathrm{Sym}^2(T_1) = A_1 \oplus E \oplus T_2$ (which carries no $T_1$)
follows by character arithmetic, relying additionally on Lean's compiled
evaluator (\texttt{native\_decide}). Full details, including the per-theorem
axiom certificate, are in
SI~Section~\ref{sec:lean}.

\subsection*{Relation to adjacent literatures}

The relationship to the transform-domain tubal literature, and to the
tubal Eckart--Young results of Mor and Avron (established for commutative
tubes), was set out with the $\starG$-SVD above
(Theorem~\ref{thm:eckart_young}). A second, physics-side lineage deserves the same care. In quantum many-body
computation, truncating in a symmetry-adapted block basis has been standard
practice for over two decades: non-abelian density-matrix renormalization
represents block states by irreducible representations of the symmetry group
and truncates irrep block by irrep block~\citep{mcculloch2002nonabelian}, and
symmetric tensor-network formalisms factor every tensor into degeneracy and
structural (Clebsch--Gordan) parts, with Wigner--Eckart machinery built into
standard
implementations~\citep{singh2010tensor,singh2012tensor,weichselbaum2012nonabelian}.
Per irrep block, the $\starG$-SVD performs the same symmetry-respecting
truncation, and we read that literature as two decades of independent evidence
for the utility of the operation. What it does not contain, and what this
paper contributes, is the operator-level theory of the truncation itself: the
group algebra as the ambient tensor algebra, the Eckart--Young theorem stating
that blockwise truncation under the Plancherel weights is the provably optimal
equivariant low-rank approximation together with its closed-form error, a
machine-checked proof of that optimality, the product-group composition
theorem, and the use of the per-irrep spectrum as a data-analysis and
symmetry-diagnostic instrument rather than as a compression step inside a
variational algorithm.

\subsection*{Outlook}

More broadly, the $\starG$ framework inverts the conventional relationship between
data and mathematics: rather than forcing tensorial data into the algebra of
vectors, we adapt the algebra to the symmetry of the data. The value of doing so
is not higher pooled accuracy. A well-trained network can match or exceed
$\starG$ on aggregate $R^2$, and we do not claim otherwise. The value is
exactness: a closed-form, exactly equivariant, provably optimal
representation (Theorem~\ref{thm:eckart_young}) obtained at roughly a hundred
parameters and no training, from which the per-irrep structure of the physics can
be read directly (Section~\ref{sec:wigner_eckart}, Table~\ref{tab:wigner}),
whose symmetry survives unbounded composition
(Section~\ref{sec:compounding}), and whose guarantees are machine-certified
rather than empirically estimated. The lesson is not that one never needs a
larger model; it is that when the data carries symmetry, the right algebra
supplies guarantees, exact selection rules, exact invariance, certified
optimality, that no amount of training and scale can, because approximation
error, however small, is on the wrong side of a qualitative line.

\section{Materials and Methods}

\subsection{Algorithmic Overview}
\label{sec:methods_overview}

The molecular prediction and diagnostics experiments share a common
$\starG$ pipeline of
four stages: (i)~\emph{group selection}, in which a finite group $G$ is fixed
and its convolution tensor $\calT_G$ together with the generalized Fourier
matrix $F_G$ are precomputed once and cached; (ii)~\emph{tensorial
featurization}, in which each input molecule is mapped to a tensor
$\calX \in \bbR^{n_f \times |G|}$ (or $\bbR^{n_f \times |G_1| \times \cdots
\times |G_d|}$ for product groups) by sampling a measurement basis at every
group element; (iii)~\emph{algebraic decomposition}, in which group-invariant
features are extracted from $\calX$ via the generalized Fourier transform and
the $\starG$-SVD; and (iv)~\emph{prediction}, in which a downstream regressor
(ridge regression for the linear pipeline; an MLP or a Neural-$\starG$
network for the neural pipelines) maps these features to the target property.
For the Wigner--Eckart experiment a fifth stage replaces the global
generalized Fourier power with a per-irrep decomposition
(Section~\ref{sec:wigner_eckart}). The
end-to-end procedure, the $\starG$ product computation, the $\starG$-SVD, the
invariant feature extractor, and the per-irrep decomposition are written as explicit
algorithm blocks in the Supplementary Information. Reference implementations are provided in
MATLAB (\texttt{core/StarGAlgebra.m}, \texttt{core/extractStarGFeatures.m},
\texttt{core/NeuralStarGFramework.m}) and in Python
(\texttt{python/StarGAlgebra.py}), the two passing a common regression-test
suite. The synthetic, data-scarce, and product-group results were generated
by the MATLAB scripts under \texttt{experiments/}; the flagship experiments
of Sections~\ref{sec:compounding}, \ref{sec:wigner_eckart},
\ref{sec:crystals}, \ref{sec:enn_comparison}, and \ref{sec:representation},
together with the band-limited $\mathrm{SO}(3)$ validation, were generated by
the deterministic Python scripts named in those sections (under
\texttt{python/}), each of which writes its result file directly.

\subsection{Symmetry-compounding protocol}
\label{sec:methods_compounding}

The compounding experiment (Section~\ref{sec:compounding};
\texttt{symmetry\_compounding.py}) is deterministic, offline, and seedless in
substance: one fixed draw (seed $0$) generates the generators, and every
reported quantity is a measured property of the resulting operators, not a
statistic over runs. The exact step is $T = \exp(S)$ with
$S = \tfrac{0.6}{\|A - A^{\!\top}\|_2}(A - A^{\!\top})$ and $A$ the Reynolds
average $\frac{1}{|G|}\sum_g P_g A_0 P_g^{\!\top}$ of a Gaussian matrix, so
$S$ commutes with the regular representation by construction; approximate
steps are $T_\varepsilon = \exp(S + \varepsilon B)$ with $B$ a fixed generic
antisymmetric generator normalized in spectral norm. The two trained
operators are fitted to the same exact step from the same $256$ Gaussian
input-output pairs. The approximate-class model is a two-hidden-layer
$\tanh$ perceptron ($24 \to 256 \to 256 \to 24$) trained by full-batch Adam
($5{,}000$ epochs, learning rate $10^{-3}$) with orbit augmentation, each
pair presented in five frames (the identity and four random group
translates); the exact-class model is the least-squares fit over the
$24$-dimensional commutant of the regular representation, an equivariant
linear layer in the sense of Corollary~\ref{cor:unification}, which recovers
the step to $10^{-15}$. Both are unrolled at inference exactly as the
synthetic operators are, and the perceptron's iterates remain on the training
scale throughout (state norm within $8\%$ of $\sqrt{24}$ at every
checkpoint), so its drift is not a norm artifact. The single-step
equivariance error is measured as the mean relative commutator residual over
$64$ Gaussian probe signals and all $24$ group elements; drift of the
$M$-fold iterate and the $A_2$ isotypic fraction are recorded at
logarithmically spaced checkpoints $M \le 400$. All invariants
(orthogonality, exact-step commutation, projector idempotence) are verified
at runtime and the result JSON is regenerable by a single script run.

\subsection{Crystal enforcement layer}
\label{sec:methods_enforcement}

The physicality layer of Section~\ref{sec:crystals} is the point-group
Reynolds projector assembled per crystal from its spglib-determined point
group, applied in the harmonic ($l = 0, 2, 4$) decomposition of the elastic
tensor; its composition with the nearest-positive-semidefinite projection
(eigenvalue clipping in Voigt form) yields the Born-stability repair. Models:
a tuned gradient-boosted regressor and a Ridge baseline
(\texttt{symmetry\_consistency.py}, \texttt{enforcement\_screening.py},
\texttt{anisotropy\_screening.py}; 3 seeds where train/test splits enter),
and a trained crystal graph network over a training-size sweep
(\texttt{crystal\_gnn\_data\_efficiency.py}; GPU job). Data: the $1{,}181$
crystals of the de Jong elastic-tensor set loaded through matminer.

\subsection{Feature Construction}
\label{sec:methods_featurizer}

For the single-group and product-group experiments, molecular features are inner
products with a rotating measurement basis at angles $\theta_g = 2\pi g/n$. For
the product group, axial features use periodic $z$-embeddings and coupled features
are angular$\times$axial products; the two actions commute because $z$-rotation
modifies $(x,y)$ not $z$, while $z$-translation modifies $z$ not $(x,y)$.

For the Wigner--Eckart experiment, features are computed under the 24 rotations
of the chiral octahedral group $O$ (6 face rotations at
$90^\circ/180^\circ/270^\circ$ about coordinate axes, 8 vertex rotations at
$120^\circ/240^\circ$ about body diagonals, 6 edge rotations at $180^\circ$
about edge midpoints, plus identity). This group is a subgroup of SO(3) whose
irreps (A$_1$, A$_2$, E, T$_1$, T$_2$) correspond to angular momentum channels
$l = 0, 0, 2, 1, 2$.

\subsection{Dipole Vector Computation}

The QM9 .xyz files include Mulliken partial charges $q_i$ as a fifth column.
The dipole vector is computed as $\boldsymbol{\mu} = \sum_i q_i \mathbf{r}_i$,
yielding components $\mu_x, \mu_y, \mu_z$ that transform as a rank-1 tensor
(vector) under rotation. This provides ground-truth targets with known
transformation properties for the Wigner--Eckart test.

\subsection{Per-Irrep Fourier Decomposition}

For each irrep $\rho$ of dimension $d_\rho$ with representation matrices
$\{\rho(g)\}_{g \in G}$, the Fourier transform of feature row $j$ is
$\hat{X}_j^\rho = \sqrt{d_\rho/|G|} \sum_{g} X(j,g) \, \rho(g)$, a
$d_\rho \times d_\rho$ matrix. The per-irrep power is $\|\hat{X}_j^\rho\|_F^2$.
Per-irrep features (one power value per feature row per irrep) are used
independently as predictors for each quantum property via ridge regression,
yielding per-irrep \Rsq{} values. Pseudocode is given in
Algorithm~\ref{alg:per_irrep}.

\subsection{Invariant Feature Extraction}
\label{sec:invariant_features}

Given a sample tensor $\calX \in \bbR^{n_f \times |G|}$, the invariant feature
vector concatenates seven complementary descriptors, all of which are exactly
invariant under the left action of $G$ (proofs in the SI Equivariance Proofs section):
(a)~the DC component $\bar{\calX}_j = \tfrac{1}{|G|} \sum_g \calX(j,g)$ for
each feature row $j$;
(b)~the AC energy $\sigma_j = \mathrm{std}_g\,\calX(j,g)$;
(c)~the total per-frequency power $\sum_j |\hat{\calX}(j,k)|^2$ for each
generalized Fourier bin $k$, where $\hat{\calX} = \calX F_G$ is the
generalized Fourier transform of $\calX$ along the group dimension;
(d)~per-row generalized Fourier power $|\hat{\calX}(j,k)|^2$ for the first
$K = 14$ equivariant rows;
(e)~the singular tube norms $\|\mathbf{s}_i\|_F$ for $i = 1, \ldots,
\min(p, q)$ obtained from the $\starG$-SVD of a reshaped $p \times q \times
|G|$ tensor (the reshape $(p, q)$ is chosen to maximize $\min(p, q)$);
(f)~the rows of $\calX$ identified as invariant by row-variance (constant
under the group action); and
(g)~four spectral statistics of the unfolded matrix (nuclear norm, spectral
norm, condition number, and entropy of the singular-value distribution).
Features are $z$-score normalized using statistics computed from training
data, and an unregularized intercept column is appended. Pseudocode is given
in Algorithm~\ref{alg:features}.

\subsection{Ridge Regression and Rank Selection}

The downstream linear model is ridge regression with hyperparameter
$\lambda$ selected from the geometric grid
$\{10^{-3}, 10^{-2}, \ldots, 10^{2}, 10^{3}\}$ by validation MSE. The
intercept is unregularized. The number of singular tubes retained in the
$\starG$-SVD feature block is set to $\min(p, q)$ where $(p, q)$ is the
optimal rectangular reshape of $n_f$ (Section~\ref{sec:invariant_features});
no further truncation is applied because the Eckart--Young theorem guarantees
a closed-form bound on the truncation error
(Theorem~\ref{thm:eckart_young}). Total parameter count for $\starG$-SVD +
Ridge is $1 + d_{\mathrm{feat}}$, where $d_{\mathrm{feat}}$ is the number of
non-degenerate feature columns retained after $z$-score normalization (for
example $143$ on full QM9 with the $48$-row featurizer and $63$ in the
data-scarce $14$-row setting, and $185$ on the product-group task).

\subsection{Baseline Architectures and Training}
\label{sec:baselines}

All four neural baselines share the same default hidden width
$[64, 32]$, ReLU activations on hidden layers, a linear output, He
initialization, full-batch validation, and the Adam optimizer
($\beta_1 = 0.9, \beta_2 = 0.999, \varepsilon = 10^{-8}$) with early stopping
on validation MSE (patience 20). They differ in three respects: (i)~the
input representation $X_{\mathrm{in}}$, (ii)~the training-set construction,
and (iii)~the parameter count. Table~\ref{tab:baselines_main} summarizes
these differences; the complete per-method specifications and pseudocode are
in SI~Section~\ref{sec:baseline_alg}, and the per-experiment hyperparameter
sheet is SI Table~1.

\textbf{Training recipe (all MLP baselines).} To ensure the baselines are
trained to a modern standard and the comparison is not confounded by an
unfavorable recipe, every MLP uses: (i)~target standardization on the
training split, with predictions de-standardized before scoring, matching the
implicit scaling of the closed-form $\starG$ ridge regressor; (ii)~L2 weight
decay selected on the validation split over the grid
$\{0, 10^{-4}, 10^{-3}, 10^{-2}\}$; and (iii)~a \texttt{ReduceLROnPlateau}
learning-rate schedule on validation loss. On the full-QM9 matched-input
comparison (SI~Section~\ref{sec:qm9_accuracy}) all molecule-level methods, both
$\starG$ variants and the three MLPs, consume the identical $48$-row angular
feature tensor of Section~\ref{sec:methods_featurizer}. The default hidden
width is $[64,32]$; the augmented-MLP capacity sweep
(SI~\S\ref{sec:aug_capacity}) varies it up to $[1024,512]$.

\begin{table}[h]
\centering
\caption{Baseline summary. Hidden $= [64, 32]$, ReLU, Adam, He init,
patience 20. Parameter counts are for QM9.}
\label{tab:baselines_main}
\begin{tabular}{@{}p{3.2cm}p{5.0cm}p{3.2cm}p{1.6cm}p{1.5cm}@{}}
\toprule
\textbf{Method} & \textbf{Input $X_{\mathrm{in}}$} & \textbf{Train data} &
\textbf{Equiv.} & \textbf{Params} \\
\midrule
Standard MLP   & raw slice $\calX(:, e)$, normalized        & native ($n$)        & approx & 5{,}249 \\
Invariant MLP  & $[\mathrm{mean}, \mathrm{std}, \min, \max]_g \calX$ & native ($n$) & exact  & 14{,}465 \\
Augmented MLP  & raw slice $\calX(:, e)$, normalized        & augmented ($|G|\,n$) & approx & 5{,}249 \\
Neural $\starG$ & $\starG$-features (Sec.~\ref{sec:invariant_features}) & native ($n$) & exact  & 62{,}625 \\
$\starG$-SVD + Ridge & $\starG$-features (Sec.~\ref{sec:invariant_features}) & native ($n$) & exact & \textbf{144} \\
\bottomrule
\end{tabular}
\end{table}

\subsection{Composed-symmetry task}
\label{sec:composed_methods}

Each sample is a point cloud of $8$ to $12$ points in $\bbR^3$, each carrying a
discrete label $c_i \in \{0, \ldots, C-1\}$ ($C=4$). For every point we build a
local rotation-covariant descriptor sampled over the $|O|=24$ octahedral frames and
tensor it with the label one-hot $\delta_{c_i}$, giving a per-point $\starG$ feature
on the product group $G = O \times \bbZ_C$; these features are mean-pooled over
points. The target applies a fixed random $\starG$ operator to the pooled feature
and reads out the power spectrum on the joint-nontrivial product irreps (spatial
irrep nontrivial \emph{and} color character nontrivial), weighted by fixed random
coefficients. Each joint power is a rotation-invariant pairwise coupling of point
environments modulated by $\cos\!\big(2\pi k (c_i - c_j)/C\big)$, so the target is
exactly invariant under any octahedral rotation and any global label shift and
vanishes if either the spatial or the label variation is removed. The
$\star_{O}$-only baseline uses the identical spatial descriptor with the label
folded into ordinary feature channels, which a global shift permutes and the
rotation-invariant readout does not undo. All models train for $300$ epochs (Adam,
learning-rate warmup, early stopping on validation loss); the SchNet baseline is
the PyTorch Geometric reference implementation with the label supplied as an
additional node feature. The construction, the four models, the validation suite,
and the sweep driver are in the code repository
(\texttt{composed\_symmetry\_knockout.py}).

\subsection{Reproducibility}

All \emph{learning} experiments use at least 3 random seeds (5 for the
full-QM9 $\starG$, MLP,
and SchNet rows and the data-scarce sweep) and a 70/15/15 train/validation/test
split unless a section states its own protocol; the deterministic
demonstrations (the compounding experiment, the enforcement projections, and
the band-limited $\mathrm{SO}(3)$ validation) involve no training and are
exactly reproducible by a single script run. The same random seeds (\texttt{seed}, $111 \cdot \texttt{seed}$ and
$31 \cdot \texttt{seed}$) drive train/val/test partitioning, weight
initialization, and mini-batch shuffling so that runs are reproducible up to
floating-point non-determinism. End-to-end runtimes (single seed,
1{,}000 molecules) on a 2024 desktop CPU are reported in SI
Table~2 and range from 0.7~s ($\starG$-SVD~+~Ridge) to 4.0~s (Augmented
MLP).

\section*{Data Availability}

All data needed to evaluate the conclusions of this study are publicly
available. The QM9 dataset~\citep{ramakrishnan2014qm9} main structure file
(\texttt{dsgdb9nsd.xyz.tar.bz2}) can be downloaded directly from
\url{https://figshare.com/ndownloader/files/3195389}; the full collection,
including property labels and auxiliary files, is archived at
\url{https://doi.org/10.6084/m9.figshare.c.978904.v5}. The QM7-X
dataset~\citep{hoja2021qm7x}, which supplies the molecular polarizability
tensors, is archived at \url{https://doi.org/10.5281/zenodo.4288677}. The
elastic-tensor dataset of de Jong \emph{et al.}~\citep{dejong2015charting}
(1,181 inorganic crystals) is obtained through \texttt{matminer} via
\texttt{load\_dataset("elastic\_tensor\_2015")}.

\section*{Code Availability}

% Zenodo deposit published 2026-07-14 (curated archive: code, Lean, committed
% result data, project READMEs; no manuscript sources). Concept DOI below
% resolves to the latest version; version 1.0.0 DOI: 10.5281/zenodo.21359576.
Open-source implementations of the $\starG$ algebra in MATLAB and Python,
together with scripts to reproduce every experiment, table, and figure (see
\texttt{python/large\_scale/REPRODUCE.md}), are archived on Zenodo at
\url{https://doi.org/10.5281/zenodo.21359575}.  The Lean~4 formalization of the
algebraic proofs (zero \texttt{sorry}) is included under \texttt{lean/}, with
its per-theorem axiom certificate in \texttt{lean/StarG/Audit.lean}.

\section*{Acknowledgements}

We wish to acknowledge Tammy Kolda for early feedback and her foundational
contributions to the tensor algebra field, and Tess Smidt and Joe Kileel
for engaging discussions. P.H. gratefully acknowledges the IBM internship program, which enabled key components of this study.

\section*{Funding}

This work received no external funding. P.H. was supported by the IBM
internship program.

%\section*{Author Contributions}
%
%\textbf{Conceptualization:} L.H., S.U.
%\textbf{Methodology:} H.A., L.H., P.H., M.K.
%\textbf{Formal analysis:} H.A., L.H., P.H., M.K.
%\textbf{Software:} L.H., P.H., S.U.
%\textbf{Formal verification (Lean~4):} L.H., P.H.
%\textbf{Investigation:} K.L.C., L.H., P.H., D.H., V.K., S.U.
%\textbf{Visualization:} L.H., P.H., S.U.
%\textbf{Writing (original draft):} L.H., P.H., S.U.
%\textbf{Writing (review and editing):} all authors.
%\textbf{Supervision:} L.H.
%\textbf{Project administration:} L.H.
%\textbf{Funding acquisition:} L.H.
%
%\section*{Competing Interests}
%
%The authors declare no competing interests.

% Science Advances requires a single reference list, in the main text, covering
% citations in BOTH the paper and the supplementary materials. These nine are
% cited only in supplementary.tex; \nocite pulls them into this list so that
% main.pdf and combined.pdf carry identical, identically numbered references.
\nocite{batatia2022mace,bekka2025crystals,bradley1972,geiger2022e3nn,
jaming2009characterization,mackey1958unitary,schutt2018schnet,thoma1968,
zhang2025singlearchitecture}

%\bibliographystyle{unsrtnat}
%\bibliography{references}
\input{combined.bbl}

\clearpage
\setcounter{section}{0}
\setcounter{figure}{0}
\setcounter{table}{0}
\setcounter{theorem}{0}
\renewcommand{\thesection}{S\arabic{section}}
\renewcommand{\thesubsection}{S\arabic{section}.\arabic{subsection}}
\renewcommand{\thefigure}{S\arabic{figure}}
\renewcommand{\thetable}{S\arabic{table}}
\renewcommand{\thetheorem}{S\arabic{theorem}}


\begin{center}
 {\LARGE \textbf{Supplementary Information}}\\[10pt]
% {\Large Exact Symmetry as Algebra: A Machine-Verified Tensor Calculus that Enforces Physical Selection Rules}\\[10pt]
% {\large Hoyos, Ubaru, Huh, Kalantzis, Clarkson, Kilmer, Avron, Horesh}
\end{center}


%\newpage

%% ========================================================================
\section{Mathematical Foundations}
%% ========================================================================

\subsection{Notation}

\begin{itemize}
\item $G$: finite group of order $n = |G|$ with identity $e$.
\item $\hat{G}$: set of equivalence classes of irreducible unitary representations.
\item $d_\rho = \dim(\rho)$ for $\rho \in \hat{G}$.
\item $\calA(:,:,g)$: frontal slice at group element $g$. $\calA_{ij} = \calA(i,j,:)$: tube at indices $i,j$.
\end{itemize}

\subsection{The Group Algebra}

\begin{definition}[Group Algebra]
$\bbR[G]$ is the vector space of formal sums $\sum_{g \in G} a_g g$ with convolution product:
\begin{equation}
\left(\sum_g a_g g\right) \cdot \left(\sum_h b_h h\right) = \sum_{c \in G} \left(\sum_{g \in G} a_g b_{g^{-1}c}\right) c,
\end{equation}
 where $a_g$ and $b_h$ are scalar coefficients in $\bbR$.
\end{definition}

\subsection{The Convolutional Ring $\mathbb{K}_G$}

\begin{definition}[Convolutional Ring $\mathbb{K}_G$]\label{def:conv_ring}
The convolutional ring $\mathbb{K}_G$, for $G = G_1 \times \cdots \times G_d$ and $n_i = |G_i|$, is the ring of tensors in $\mathbb{R}^{n_1 \times n_2 \times \cdots \times n_d}$ with pointwise addition and convolution product:
\begin{align}\label{eq: group convolution of tensors}
    (\mathcal{A} \cdot \mathcal{B})(c_1, \ldots, c_d) = \sum_{(a_1 \ldots, a_d) \in G}
\mathcal{A}(a_1, \ldots, a_d) \mathcal{B}(a_1^{-1}c_1, \ldots, a_d^{-1}c_d),
\end{align}
for all $(c_1, \dots, c_d) \in G$. As a ring, $\mathbb{K}_G$ is isomorphic to
the group algebra $\bbR[G]$ above under $a \leftrightarrow (a_1, \ldots, a_d)$,
and it factors as
$\mathbb{K}_G \cong \mathbb{K}_{G_1} \otimes \cdots \otimes \mathbb{K}_{G_d}$.
The entries (tubes) of a $\starG$ tensor are elements of $\mathbb{K}_G$.
\end{definition}


%% ========================================================================
\section{The Convolution Tensor, Spectral Decomposition, and Generalized Fourier Matrix}
%% ========================================================================

\begin{definition}[Convolution Tensor]
$\calT \in \bbR^{n \times n \times n}$ is defined by $\calT(a,b,c) = \delta(g_a g_b = g_c)$ (also known as the structure constants of the group algebra).
\end{definition}

\begin{proposition}[Properties of $\calT$]\label{prop:T_props}
\begin{enumerate}
\item[(i)] Associativity:
\[
\sum_d \calT(a,b,d)\,\calT(d,c,e)
        \;=\; \sum_d \calT(a,d,e)\,\calT(b,c,d).
\]
\item[(ii)] Identity: $\calT(e,b,c) = \delta_{bc}$.
\item[(iii)] Each slice $\calT(a,:,:)$ is a permutation matrix.
\end{enumerate}
\end{proposition}

\begin{proof}
(i) follows from associativity of group multiplication; both sides equal $\delta(g_a g_b g_c = g_e)$. (ii)--(iii) follow from the group axioms.
\end{proof}

\begin{definition}[Generalized Fourier Transform Matrix]
$F_G \in \bbC^{n\times n}$ is defined row-wise: the row $F_G(g, :)$ is given by the concatenation
$[\mathrm{rvec}(\rho_1(g)), \dots, \mathrm{rvec}(\rho_\ell(g))]$,
where $\mathrm{rvec}(\rho_i(g))$ denotes the row-vectorization of the matrix $\rho_i(g)$ for each $\rho_i \in \hat{G}$. For abelian groups, $F_G$ is a generalized Fourier matrix and for cyclic groups, it reduces to the standard DFT matrix. For non-abelian groups $F_G$ is invertible but in general neither unitary nor block-unitary.
\end{definition}

\begin{theorem}[Peter--Weyl Spectral Decomposition]\label{thm:pw_spectral}
\begin{equation}
\calT(a,b,c) = \sum_{i,j,k} \calC(i,j,k) \, F_G(a,i) \, F_G(b,j) \, F_G^{-1}(c,k)
\end{equation}
where $\calC$ is a core tensor  that is typically sparse. For abelian groups, $\calC$ is diagonal.
\end{theorem}

\begin{proof}
By the Peter--Weyl theorem the matrix elements
$\{\sqrt{d_\rho}\,\rho_{st}(g) : \rho \in \hat{G},\ 1 \le s,t \le d_\rho\}$ form an
orthonormal basis of $L^2(G)$, and this is exactly the basis assembled column by
column into $F_G$: column $i \leftrightarrow (\rho, s, t)$ carries $\rho_{st}$. In
this basis the group-algebra multiplication is the convolution theorem
$\widehat{(f*h)}(\rho) = \hat{f}(\rho)\,\hat{h}(\rho)$, i.e.\ block-diagonal matrix
multiplication: the product of the basis elements $(\rho, s, u)$ and $(\rho', u', t)$
is nonzero only when $\rho = \rho'$ and $u = u'$, and then equals the basis element
$(\rho, s, t)$ (the matrix-unit rule $E^{\rho}_{su} E^{\rho}_{u't} = \delta_{uu'}
E^{\rho}_{st}$). Transporting the convolution tensor $\calT$, with
$\calT(a,b,c) = \mathbf{1}[ab=c]$, to this basis through $F_G$ (contracting mode~3
with $F_G^{-1}$, the exact inverse) therefore gives the stated decomposition with
core $\calC$ supported on the index triples
\begin{equation}
\big((\rho,s,u),\ (\rho,u,t),\ (\rho,s,t)\big),
\qquad \rho \in \hat{G},\ \ 1 \le s,u,t \le d_\rho,
\end{equation}
on which $\calC$ carries the multiplication structure constants of the block
$M_{d_\rho}(\bbC)$ (the Peter--Weyl normalization $d_\rho/n$ being absorbed into
$F_G$ and $F_G^{-1}$). This support is the union over $\rho$ of the $M_{d_\rho}$
multiplication tables, so $\calC$ is sparse; when every irrep is one-dimensional
(abelian $G$) each index collapses to $(\rho)$, the support reduces to
$\{(\rho,\rho,\rho)\}$, and $\calC$ is diagonal.
\end{proof}

\begin{corollary}
For $G = \bbZ_n$: $F_G = \mathrm{DFT}_n$ and $\calC$ is diagonal, recovering the circular convolution theorem  (i.e., the $t-$product).
\end{corollary}

%% ========================================================================
\section{The Group Fourier Transform of a Tensor}
%% ========================================================================

\begin{definition}[Group Fourier Transform of a Tensor]\label{def:tensor_ft}
For $\calA \in \bbR^{\ell \times m \times n}$, the Group Fourier transform $\mathcal{F}_G$ assigns to each irrep $\rho \in \hat{G}$ the $\ell d_\rho \times m d_\rho$ block matrix whose $(i,j)$ block is
\begin{equation}
    \hat{\calA}(i,j,\rho) = \sum_{g \in G} \calA(i,j,g)\,\rho(g) \in \bbC^{d_\rho \times d_\rho}, \qquad i \le \ell,\ j \le m.
\end{equation}
The full Fourier representation is the block-diagonal matrix $\oplus_{\rho \in \hat{G}} \hat{\calA}(:,:,\rho)$.
\end{definition}

\begin{proposition}[Group Fourier Inversion Theorem]
Given $\hat{\calA}(:,:,\rho)$ for each $\rho \in \hat{G}$, the inverse Group Fourier transform recovers
\begin{equation}
    \calA(i,j,g) = \sum_{\rho \in \hat{G}} \frac{d_\rho}{n}\,\mathrm{Tr}\!\left[\hat{\calA}(i,j,\rho)\,\rho(g)^H \right].
\end{equation}
\end{proposition}
\begin{proof}
Follows by applying the standard Peter--Weyl inversion theorem for each fixed $i$ and $j$.
\end{proof}

\begin{remark}
The two notions of ``Fourier transform'' used in the paper are related as follows. The contraction of $\calA$ with $F_G$ along its group dimension (as in SI Section~2) and the block-diagonal form $\oplus_\rho \hat{\calA}(:,:,\rho)$ (Definition~\ref{def:tensor_ft}) are equivalent via the row-vectorization reshaping $\mathrm{rvec}$ used in the construction of $F_G$. The block-diagonal form is used for the SVD computation (one standard matrix SVD per irrep block); the $F_G$-contraction form is used for the product computation and for establishing the spectral decomposition of $\calT$.
\end{remark}

\begin{algorithm}[h]
\caption{$\starG$ product (used by every $\starG$-based method)}
\label{alg:starg_product}
\begin{algorithmic}[1]
\Require $\calA \in \bbR^{\ell \times m \times n}$, $\calB \in \bbR^{m \times p \times n}$, generalized Fourier matrix $F_G \in \bbC^{n \times n}$, irrep block sizes $\{d_\rho\}_{\rho \in \hat{G}}$ \Comment{product groups and order-$(2+d)$ tensors via Remark~\ref{rem:flatten}}
\Ensure $\calC = \calA \starG \calB \in \bbR^{\ell \times p \times n}$
\State $\hat{\calA} \gets \calA \times_3 F_G$ \Comment{Generalized Fourier transform along group dimension}
\State $\hat{\calB} \gets \calB \times_3 F_G$
\For{$\rho \in \hat{G}$ \textbf{in parallel}}
    \State extract $\hat{\calA}_\rho \in \bbC^{\ell d_\rho \times m d_\rho}$, $\hat{\calB}_\rho \in \bbC^{m d_\rho \times p d_\rho}$ from block-diagonal form
    \State $\hat{\calC}_\rho \gets \hat{\calA}_\rho \cdot \hat{\calB}_\rho$ \Comment{ordinary matrix product}
\EndFor
\State assemble $\hat{\calC}$ from $\{\hat{\calC}_\rho\}_\rho$ in block-diagonal form
\State $\calC \gets \mathrm{Re}(\hat{\calC} \times_3 F_G^{-1})$ \Comment{inverse Generalized Fourier transform; imaginary part is zero up to roundoff for real inputs}
\State \Return $\calC$
\end{algorithmic}
\end{algorithm}

%% ========================================================================
\section{The $\starG$ Algebra: Properties and Proofs}
%% ========================================================================

\begin{definition}[$\starG$ Product]\label{def:starg_product}
For $\calA \in \bbR^{\ell \times m \times n}$, $\calB \in \bbR^{m \times p \times n}$:
\begin{equation}
(\calA \starG \calB)_{ij}(c) = \sum_k \sum_{a \in G} \calA_{ik}(a) \calB_{kj}(a^{-1}c).
\end{equation}
Equivalently: $\widehat{(\calA \starG \calB)}(:,:,\rho) = \hat{\calA}(:,:,\rho) \cdot \hat{\calB}(:,:,\rho)$ for each irrep $\rho \in \hat{G}$.
\end{definition}

\begin{remark}[One group or many: the flattening convention]\label{rem:flatten}
Definition~\ref{def:starg_product} is stated for a single group $G$ and
order-3 tensors, whereas the main text (Results, Eq.~(3)) states the product
for $G = G_1 \times \cdots \times G_d$ and order-$(2+d)$ tensors, with entries
in the convolutional ring $\mathbb{K}_G$ of
Definition~\ref{def:conv_ring}. The two are the same definition. The bijection
$a \leftrightarrow (a_1, \ldots, a_d)$ identifies the product group with a
single group of order $n = n_1 \cdots n_d$; under it, tubes become
$d$-dimensional arrays (elements of $\mathbb{K}_G$), the single-group
convolution of Definition~\ref{def:starg_product} becomes the multi-index
convolution of Eq.~(3) of the main text, the Fourier matrix factors as
$F_G = F_{G_1} \otimes \cdots \otimes F_{G_d}$, and the irreducible
representations are the Kronecker products
$\rho = \rho_1 \otimes \cdots \otimes \rho_d$ of the factors' irreps
(Section~6). Consequently every algorithm in this document, stated over
$\bbR^{\ell \times m \times n}$ for readability, applies verbatim to
order-$(2+d)$ tensors over product groups: one runs it on the flattened
group, and each Fourier block is indexed by a tuple of irreps. The order-3
form is a notational flattening, not a restriction.
\end{remark}

\begin{proposition}[Algebraic Properties]
(i)~Associativity. (ii)~Distributivity. (iii)~Identity: $\calI(:,:,e)=I$, $\calI(:,:,g\neq e)=0$.
(iv)~$(\calA \starG \calB)^H = \calB^H \starG \calA^H$.
\end{proposition}

\begin{proof}
All follow from the corresponding matrix properties applied per-irrep in the Fourier domain (Definition~\ref{def:tensor_ft}), plus linearity and invertibility of the Fourier transform. In the Lean~4 development these identities are instead proved directly, by combinatorial reindexing of the group-convolution sum, so they depend on none of the Fourier-analytic axioms.
\end{proof}

%% ========================================================================
\section{The $\starG$-SVD: Full Proof of Optimality}
%% ========================================================================

\begin{theorem}[$\starG$-SVD Existence]\label{thm:svd_exist}
Every $\calA \in \bbR^{\ell \times m \times n}$ admits $\calA = \calU \starG \calS \starG \calV^H$ where $\calU,\calV$ are $\starG$-unitary and $\calS$ is f-diagonal.
\end{theorem}

\begin{proof}
For each irrep $\rho \in \hat{G}$, $\hat{\calA}(:,:,\rho)$ is a standard matrix admitting SVD: $\hat{\calA}(:,:,\rho) = U_\rho \Sigma_\rho V_\rho^H$. Setting $\hat{\calU}(:,:,\rho) = U_\rho$, $\hat{\calS}(:,:,\rho) = \Sigma_\rho$, $\hat{\calV}(:,:,\rho) = V_\rho$ and applying the inverse Fourier transform yields the $\starG$-SVD. Unitarity: $\widehat{\calU^H \starG \calU}(:,:,\rho) = U_\rho^H U_\rho = I$ for all $\rho$, so $\calU^H \starG \calU = \calI$. When the tensor is real and its irreps occur in complex-conjugate pairs, choosing the per-block SVDs conjugate-consistently on paired irreps ($U_{\bar\rho} = \overline{U_\rho}$, and likewise for $\calS$ and $\calV$) makes the inverse transform real, the same convention that yields a real t-SVD; over the real-irrep formalization used in Lean this holds automatically.
\end{proof}

\begin{algorithm}[h]
\caption{$\starG$-SVD}
\label{alg:svd}
\begin{algorithmic}[1]
\Require $\calA \in \bbR^{\ell \times m \times n}$ \Comment{product groups and order-$(2+d)$ tensors via Remark~\ref{rem:flatten}}
\Ensure $\calU, \calS, \calV$
\State Compute $\hat{\calA}(:,:,\rho)$ for all $\rho \in \hat{G}$ using $\mathcal{F}_G$ (Definition~\ref{def:tensor_ft})
\For{$\rho \in \hat{G}$ \textbf{in parallel}}
    \State $[U_\rho, \Sigma_\rho, V_\rho] \gets \mathrm{SVD}(\hat{\calA}(:,:,\rho))$
\EndFor
\State Apply $\mathcal{F}_G^{-1}$ to $\{U_\rho\}$, $\{\Sigma_\rho\}$, $\{V_\rho\}$ to obtain $\calU, \calS, \calV$
\end{algorithmic}
\end{algorithm}

\begin{theorem}[Eckart--Young for $\starG$]\label{thm:ey_full}
Let $\calA = \calU \starG \calS \starG \calV^H$ be the $\starG$-SVD returned by
Algorithm~\ref{alg:svd}, whose per-irrep singular values $\sigma_j(\rho)$ are in
non-increasing order.
\emph{(i) Per-block form.} For any per-block rank profile $(k_\rho)_{\rho\in\hat G}$,
the blockwise truncation keeping the top $k_\rho$ singular values of each
$\hat\calA(:,:,\rho)$ satisfies
\begin{equation}
\Big\|\calA-\calA_{(k_\rho)}\Big\|_F^2 = \sum_{\rho\in\hat G}\frac{d_\rho}{n}
\sum_{j>k_\rho}\sigma_j(\rho)^2 \;\le\; \|\calA-\calB\|_F^2
\end{equation}
for any $\calB$ with $\mathrm{rank}\,\hat\calB(:,:,\rho)\le k_\rho$ for every $\rho$.
\emph{(ii) Tubal form.} Let $\mathbf{s}_1,\ldots,\mathbf{s}_r$ be the singular tubes
(with norms automatically ordered, $\|\mathbf{s}_1\|_F\ge\cdots\ge\|\mathbf{s}_r\|_F$),
and $\calA_k$ the truncation retaining the leading $k$. Then
\begin{equation}
\|\calA - \calA_k\|_F^2 = \sum_{i=k+1}^r \|\mathbf{s}_i\|_F^2 \leq \|\calA - \calB\|_F^2
\end{equation}
for any $\calB$ with $\starG$-rank $\leq k$; this is~(i) at the coupled profile
$k_\rho=k\,d_\rho$.
\end{theorem}

\begin{proof}
Part~(i) is the classical matrix Eckart--Young theorem applied independently to each
block $\hat\calA(:,:,\rho)$ at rank $k_\rho$ (Step~2), weighted by the Parseval
isometry (Step~1); part~(ii) follows by specializing to $k_\rho=k\,d_\rho$ and
regrouping the discarded block singular values by tube (Step~3).

\textbf{Step 1 (Parseval).} By the generalized Fourier transform's isometry (Peter--Weyl):
\begin{equation}
\|\calA - \calB\|_F^2 = \sum_{\rho \in \hat{G}} \frac{d_\rho}{n}\|\hat{\calA}(:,:,\rho) - \hat{\calB}(:,:,\rho)\|_F^2.
\end{equation}

\textbf{Step 2 (Per-irrep Eckart--Young).} Because $\calS$ is f-diagonal, each
singular tube contributes one $d_\rho \times d_\rho$ diagonal block to
$\hat{\calS}(:,:,\rho)$; a tensor of $\starG$-rank $\leq k$ therefore has at most
$k$ nonzero such blocks, so $\mathrm{rank}(\hat{\calB}(:,:,\rho)) \leq k\,d_\rho$
for each $\rho$. Since Algorithm~\ref{alg:svd} sorts the singular values of every
block in non-increasing order, the leading-$k$-tube truncation
$\hat{\calA}_k(:,:,\rho)$ retains exactly the top $k\,d_\rho$ singular values of
$\hat{\calA}(:,:,\rho)$, which by the classical Eckart--Young theorem at rank
$k\,d_\rho$ is the best rank-$(k\,d_\rho)$ approximation of that block:
\begin{equation}
\|\hat{\calA}(:,:,\rho) - \hat{\calA}_k(:,:,\rho)\|_F^2 \leq \|\hat{\calA}(:,:,\rho) - \hat{\calB}(:,:,\rho)\|_F^2.
\end{equation}

\textbf{Step 3 (Summation).} Tube $i$ collects the block singular values
$\sigma_{(i-1)d_\rho + 1}(\rho), \ldots, \sigma_{i d_\rho}(\rho)$, so
$\|\mathbf{s}_i\|_F^2 = \sum_\rho \frac{d_\rho}{n}\sum_{t=1}^{d_\rho}\sigma_{(i-1)d_\rho + t}(\rho)^2$
and the singular values discarded from block $\rho$ (those indexed $j > k\,d_\rho$)
are exactly those carried by the tubes $\mathbf{s}_{k+1}, \ldots, \mathbf{s}_r$.
Summing Step 2 over $\rho$ and regrouping by tube,
\begin{equation}
\|\calA - \calA_k\|_F^2 = \sum_{\rho \in \hat{G}} \frac{d_\rho}{n}\sum_{j > k d_\rho} \sigma_j(\rho)^2
= \sum_{i=k+1}^r \|\mathbf{s}_i\|_F^2 \leq \|\calA - \calB\|_F^2.
\end{equation}
\end{proof}

\begin{remark}
This is \emph{exact} optimality. By contrast: CP decomposition is NP-hard to compute optimally; Tucker/HOSVD provides only quasi-optimal guarantees with factor $\sqrt{d}$ \cite{desilva2008tensor}; tensor-train has no global optimality. The $\starG$-SVD achieves both polynomial-time computation and exact optimality by leveraging group structure.
\end{remark}

\begin{remark}[Semisimple generality]
The proof of the per-block form uses only that the transform is a
$*$-isomorphism onto a direct sum of matrix algebras carrying a weighted
Frobenius norm. The per-block Eckart--Young statement therefore holds verbatim
for any finite-dimensional semisimple $*$-algebra with a fixed
Artin--Wedderburn decomposition; the group algebras of this paper are the
equivariant instance, where the blocks additionally carry the physical
(irrep) labels that the diagnostics and selection rules exploit.
\end{remark}

%% ========================================================================
\section{Product Groups: Full Proof}
%% ========================================================================

\begin{theorem}[Product Group Ring Isomorphism]\label{thm:prod_full}
For $G = G_1 \times \cdots \times G_d$:
(i)~$\bbK_G \cong \bbK_{G_1} \otimes \cdots \otimes \bbK_{G_d}$.
(ii)~$\calT_G = \calT_{G_1} \otimes \cdots \otimes \calT_{G_d}$.
(iii)~$F_G = F_{G_1} \otimes \cdots \otimes F_{G_d}$.
\end{theorem}

\begin{proof}
\textbf{(i)} Standard: $\bbR[G_1 \times \cdots \times G_d] \cong \bbR[G_1] \otimes \cdots \otimes \bbR[G_d]$ \cite{serre1977linear}.

\textbf{(ii)} The product group multiplication $(a_1,\ldots,a_d)(b_1,\ldots,b_d) = (a_1 b_1,\ldots,a_d b_d)$ gives
\begin{equation}
\calT_G(\mathbf{a},\mathbf{b},\mathbf{c}) = \prod_{i=1}^d \calT_{G_i}(a_i,b_i,c_i) = (\calT_{G_1} \otimes \cdots \otimes \calT_{G_d})(\mathbf{a},\mathbf{b},\mathbf{c}).
\end{equation}

\textbf{(iii)} Over $\bbC$ (the field of $F_G$ used here), the irreps of $G_1 \times \cdots \times G_d$ are exactly the tensor products $\rho_1 \otimes \cdots \otimes \rho_d$. Matrix elements factor: $(\rho_1 \otimes \cdots \otimes \rho_d)(g_1,\ldots,g_d) = \rho_1(g_1) \otimes \cdots \otimes \rho_d(g_d)$. By the $\mathrm{rvec}$ construction of $F_G$ and the mixed-product property of Kronecker products, $F_G = F_{G_1} \otimes \cdots \otimes F_{G_d}$.
\end{proof}

\begin{corollary}[2D Frequency Resolution]
For $G = \bbZ_{n_1} \times \bbZ_{n_2}$, the Fourier transform $F_G = \mathrm{DFT}_{n_1} \otimes \mathrm{DFT}_{n_2}$ computes a 2D DFT, resolving coupled frequencies $(f_1, f_2)$ that are invisible to either factor alone.
\end{corollary}

%% ========================================================================
\section{Invariant Feature Extraction Algorithm}
\label{sec:feat_alg}
%% ========================================================================

The $\starG$-feature extractor used in every linear and Neural-$\starG$
experiment of this paper is given in Algorithm~\ref{alg:features}. The seven
feature blocks (a)--(g) below correspond to the seven concatenated columns
of the output. Block (a) is the DC component; (b) the AC standard
deviation; (c) the global per-frequency power; (d) the per-row Fourier power
restricted to the first $K$ equivariant rows; (e) the singular tube norms
from the $\starG$-SVD; (f) the rows of $\calX$ that are constant under the
group action (recovered by row-variance thresholding); and (g) four spectral
statistics of the unfolded matrix $\calX_{(1)}$. Every feature is exactly
$G$-invariant for every finite group $G$; proofs are in SI Section~8. For the
Fourier-power features (c) and (d) the invariant quantity is the per-irrep
Frobenius power $\|\hat{\calX}(i,:,\rho)\|_F^2$, summed over each $d_\rho \times
d_\rho$ irrep block, which the extractor computes directly (invariance of the
per-irrep Fourier power, proved in SI Section~8): the group action left-multiplies
each block by $\rho(g)$, which is unitary and preserves the block Frobenius norm.
For abelian $G$, where every irrep is one-dimensional, this reduces to the
per-frequency power.

\begin{algorithm}[!tb]
\caption{$\starG$ invariant feature extraction (\texttt{extractStarGFeatures.m})}
\label{alg:features}
\begin{algorithmic}[1]
\Require batch $\calX \in \bbR^{N \times n_f \times n}$, group algebra $G$, optional normalization parameters $\Theta$ from training
\Ensure feature matrix $\Phi \in \bbR^{N \times d_{\mathrm{feat}} + 1}$ ($+1$ for the unregularized intercept) and updated $\Theta$
\If{$\Theta$ is not provided}
    \State $(p, q) \gets \arg\max_{p \cdot q \leq n_f} \min(p, q)$ \Comment{best near-square reshape for $\starG$-SVD}
    \State $n_{\mathrm{svd}} \gets \min(p, q)$
    \State \(\sigma_{\mathrm{row}} \gets \mathrm{var}_g(\calX(1, :, :))\); $\mathrm{inv\_mask} \gets \sigma_{\mathrm{row}} < 10^{-8} \cdot \max(\sigma_{\mathrm{row}})$ \Comment{rows constant under $G$}
    \State $\mathrm{eq\_idx} \gets \{j : \mathrm{inv\_mask}_j = \text{false}\}$;\ $K \gets \min(14, |\mathrm{eq\_idx}|)$
\EndIf
\State $\bar{\calX}_{ij} \gets \tfrac{1}{n} \sum_g \calX(i, j, g)$ \Comment{(a) DC component, $N \times n_f$}
\State $\sigma_{ij} \gets \mathrm{std}_g \calX(i, j, g)$ \Comment{(b) AC energy, $N \times n_f$}
\State $\hat{\calX} \gets \calX \times_3 F_G$ \Comment{Generalized Fourier transform along group dim}
\State $P^{\mathrm{col}}_{i\rho} \gets \sum_j \|\hat{\calX}(i, j, \rho)\|_F^2$ \Comment{(c) per-irrep power, $N \times |\hat{G}|$}
\State $P^{\mathrm{row}}_{i,(r-1)|\hat{G}| + \rho} \gets \|\hat{\calX}(i, \mathrm{eq\_idx}_r, \rho)\|_F^2$ \Comment{(d) per-row per-irrep power, $N \times K|\hat{G}|$}
\For{$i = 1, \ldots, N$}
    \State $X_i \gets$ pad $\calX(i, :, :)$ to $p \times q \times n$
    \State $[\,\cdot\,, \calS_i, \,\cdot\,] \gets \mathrm{starG\_SVD}(X_i)$ \Comment{(e) singular tubes}
    \State $T_{i, k} \gets \|\calS_i(k, k, :)\|_F$ for $k = 1, \ldots, n_{\mathrm{svd}}$; sort $T_{i, :}$ descending
    \State $V_{ij} \gets \calX(i, j, 1)$ for $j$ with $\mathrm{inv\_mask}_j$ \Comment{(f) direct invariants}
    \State $\Sigma_i \gets \mathrm{svd}(\calX_{(1)}(i, :, :))$ \Comment{(g) spectral statistics}
    \State $S_i \gets [\sum \Sigma_i,\ \Sigma_{i,1},\ \Sigma_{i,1} / \Sigma_{i,\mathrm{end}},\ -\sum_k \tilde{\Sigma}_{i,k} \log \tilde{\Sigma}_{i,k}]$ where $\tilde{\Sigma}_i = \Sigma_i / \sum \Sigma_i$
\EndFor
\State $\Phi \gets [\bar{\calX}\ |\ \sigma\ |\ P^{\mathrm{col}}\ |\ P^{\mathrm{row}}\ |\ T\ |\ V\ |\ S]$
\State replace NaN/Inf with 0
\If{$\Theta$ not provided}
    \State $\mathrm{keep} \gets \{j : \mathrm{std}(\Phi_{:, j}) \geq 10^{-8}\}$
    \State $\mu \gets \mathrm{mean}(\Phi_{:, \mathrm{keep}})$;\ $s \gets \mathrm{std}(\Phi_{:, \mathrm{keep}})$;\ $s_j \gets \max(s_j, 1)$
    \State store $\Theta = (p, q, n_{\mathrm{svd}}, \mathrm{inv\_mask}, \mathrm{eq\_idx}, K, \mathrm{keep}, \mu, s)$
\EndIf
\State $\Phi \gets (\Phi_{:, \mathrm{keep}} - \mu) / s$
\State $\Phi \gets [\,\mathbf{1}_N\ |\ \Phi\,]$ \Comment{prepend unregularized intercept}
\State \Return $(\Phi, \Theta)$
\end{algorithmic}
\end{algorithm}

%% ========================================================================
\section{Equivariance Proofs}
%% ========================================================================

\begin{proposition}[Equivariance of $\starG$]\label{prop:equivariance}
For every finite group $G$, the $\starG$ product is left-equivariant in its first
argument:
\[(g \cdot \calA) \starG \calB = g \cdot (\calA \starG \calB).\]
The corresponding identity in the second argument,
$\calA \starG (g \cdot \calB) = g \cdot (\calA \starG \calB)$, holds when $G$ is
abelian. For general $G$ the second argument instead intertwines the left action
with right translation: writing $(\calB \triangleleft g)(h) = \calB(hg)$,
\[\calA \starG (\calB \triangleleft g) = (\calA \starG \calB) \triangleleft g.\]
\end{proposition}

\begin{proof}
By definition of the group action and the $\starG$ product,
\begin{align*}
\bigl((g\cdot\calA)\starG\calB\bigr)_{ij}(h)
   &= \sum_k \sum_{a\in G} \calA_{ik}(g^{-1}a)\,\calB_{kj}(a^{-1}h).
\end{align*}
Substituting $a'=g^{-1}a$ (so $a=ga'$ and $a^{-1}=a'^{-1}g^{-1}$),
\begin{align*}
   &= \sum_k \sum_{a'\in G} \calA_{ik}(a')\,\calB_{kj}(a'^{-1}g^{-1}h) \\
   &= (\calA\starG\calB)_{ij}(g^{-1}h)
    = \bigl(g\cdot(\calA\starG\calB)\bigr)_{ij}(h),
\end{align*}
which proves left-equivariance in the first argument for every $G$. For the second
argument the right-translation identity is immediate,
\begin{align*}
\bigl(\calA\starG(\calB\triangleleft g)\bigr)_{ij}(h)
   &= \sum_k \sum_{a\in G} \calA_{ik}(a)\,\calB_{kj}(a^{-1}hg)
    = (\calA\starG\calB)_{ij}(hg)
    = \bigl((\calA\starG\calB)\triangleleft g\bigr)_{ij}(h).
\end{align*}
Replacing right translation by the left action would require
$\calB_{kj}(a^{-1}g^{-1}h) = \calB_{kj}(g^{-1}a^{-1}h)$ for all $a$, i.e.\
$a^{-1}g^{-1} = g^{-1}a^{-1}$, which holds for all $a$ precisely when $g$ is
central; the left-action form of the second identity therefore holds for abelian
$G$, matching the \texttt{CommGroup} hypothesis under which it is discharged in the
Lean~4 development.
\end{proof}

\begin{corollary}[Invariance of Features]
The following are invariant under $g \cdot X$:
(i)~$\|\mathbf{s}_i\|_F$ (singular tube norms);
(ii)~$\|\hat{X}(:,:,\rho)\|_F^2$ (per-irrep Fourier power);
(iii)~$\bar{X}_j = \frac{1}{n}\sum_g X(j,g)$ (DC component).
\end{corollary}

\begin{proof}
(i) The group action permutes frontal slices; the SVD is computed per-irrep where the action multiplies each block by a unitary, preserving singular values. (ii) The group action shifts $g \to g'g$ inside the sum defining $\hat{X}(:,:,\rho)$, multiplying the block by $\rho(g')$, which is unitary, leaving the Frobenius norm unchanged. (iii) $\frac{1}{n}\sum_g X(j,g'g) = \frac{1}{n}\sum_h X(j,h) = \bar{X}_j$.
\end{proof}

%% ========================================================================
\section{Canonicity of the \texorpdfstring{$\starG$}{starG} Algebra}
\label{sec:necessity}
%% ========================================================================

The $\starG$ algebra is not an arbitrary member of the transform-based
($\star_M$/t-product) family: it is the unique such algebra compatible with the
group action. We make this precise below and draw the consequence that the
algebra's entire matrix-mimetic structure, and its low-rank optimum in particular,
is determined by $G$ rather than chosen.

Fix the left regular action of $G$ on tubes, $(g\cdot x)(h)=x(g^{-1}h)$,
equivalently $g\cdot\delta_a=\delta_{ga}$, where $\delta_g\in\bbC[G]$ is the
indicator of $g$. A \emph{transform-based product} is given by an invertible
$M\in\bbC^{n\times n}$ and a partition of the $n$ coordinates into blocks of
sizes $d_1^2,\dots,d_r^2$ (so $\sum_s d_s^2=n$), which identifies the transform
domain with $\mathcal B=\bigoplus_{s=1}^r M_{d_s}(\bbC)$; one sets
$a\star_M b=M^{-1}(Ma\cdot Mb)$, where $\cdot$ is block-diagonal matrix
multiplication. This is the $\star_M$/t-product template
\cite{kernfeld2015tensor,kilmer2021tensor}, with the abelian case $d_s\equiv1$.
Because $M$ is linear and blockmult is bilinear, associative and unital,
$\star_M$ is automatically an associative unital algebra; its identity is
$M^{-1}\!\big(\bigoplus_s I_{d_s}\big)$.

\begin{theorem}[Necessity and uniqueness of $F_G$]\label{thm:necessity}
Among all transform-based products $\star_M$ on $\bbC[G]$ that
\begin{enumerate}
\item[(i)] are associative with a two-sided identity, normalized so that the
identity is the group-identity tube $\delta_e$ (equivalently $M\delta_e=\bigoplus_s I_{d_s}$), and
\item[(ii)] are left-equivariant, $(g\cdot\calA)\star_M\calB=g\cdot(\calA\star_M\calB)$
for all $g\in G$ and all $\calA,\calB$,
\end{enumerate}
the product is exactly group convolution, $\star_M=\starG$; the block sizes are
forced to be the irreducible dimensions with their multiplicities,
$\{d_1,\dots,d_r\}=\{d_\rho\}_{\rho\in\hat G}$ as multisets; and $M=\theta\circ F_G$
for a unique $\theta\in\mathrm{Aut}\big(\bigoplus_\rho M_{d_\rho}(\bbC)\big)$,
i.e.\ up to (a) a per-irrep change of basis
$\hat\calA(\rho)\mapsto P_\rho\,\hat\calA(\rho)\,P_\rho^{-1}$, $P_\rho\in\mathrm{GL}_{d_\rho}(\bbC)$,
and (b) a permutation of blocks of equal dimension. These are precisely the
choices already made in constructing $F_G$ (irrep representatives and their
ordering). Conversely $\starG$ with $M=F_G$ satisfies (i) and (ii).
\end{theorem}

\begin{proof}
\emph{Part I: $\star_M=\starG$.} Let $\iota$ be the two-sided identity; by (i)
$\iota=\delta_e$, and $g\cdot\delta_e=\delta_g$. Setting $\calA=\delta_e$ in (ii),
\begin{equation}\label{eq:transl}
\delta_g\star_M b=(g\cdot\delta_e)\star_M b=g\cdot(\delta_e\star_M b)=g\cdot b,
\qquad \forall\,g\in G,\ b\in\bbC[G],
\end{equation}
so left translation is left $\star_M$-multiplication by $\delta_g$. For any
$x=\sum_g x(g)\delta_g$, bilinearity and \eqref{eq:transl} give
\[
x\star_M b=\sum_g x(g)\,(\delta_g\star_M b)=\sum_g x(g)\,(g\cdot b)
=\Big(\textstyle\sum_g x(g)\,g\cdot\Big)b=x\starG b,
\]
since $\big(\sum_g x(g)\,g\cdot b\big)(h)=\sum_g x(g)\,b(g^{-1}h)=(x\starG b)(h)$.
Hence $\star_M=\starG$.

\emph{Part II: block sizes and transform.} By Part I,
$M(\calA\starG\calB)=M\calA\cdot M\calB$, so $M$ is an isomorphism of
$\bbC$-algebras from $(\bbC[G],\starG)$ onto $\mathcal B=\bigoplus_s M_{d_s}(\bbC)$.
By Maschke's theorem $\bbC[G]$ is semisimple, and by the Artin--Wedderburn
theorem $\bbC[G]\cong\bigoplus_{\rho\in\hat G}M_{d_\rho}(\bbC)$; the multiset of
block dimensions of a semisimple algebra is an isomorphism invariant (uniqueness
of the Wedderburn decomposition~\cite{serre1977linear}), so $\{d_s\}=\{d_\rho\}$.
The Peter--Weyl map $W\colon\calA\mapsto\bigoplus_\rho\hat\calA(\rho)$, whose
coordinate matrix is $F_G$, is also an algebra isomorphism onto
$\bigoplus_\rho M_{d_\rho}(\bbC)$ (the convolution theorem,
Theorem~\ref{thm:pw_spectral}). Hence
$\theta:=M\circ W^{-1}\in\mathrm{Aut}\big(\bigoplus_\rho M_{d_\rho}(\bbC)\big)$.
An automorphism of a direct sum of matrix algebras permutes the simple summands,
mapping each onto one of equal dimension (as $M_d(\bbC)\cong M_{d'}(\bbC)$ iff
$d=d'$), and restricts on each summand to an inner automorphism
$X\mapsto P_\rho X P_\rho^{-1}$ by the Skolem--Noether theorem. This is exactly
the freedom in (a) and (b). The converse is Proposition~\ref{prop:equivariance}
together with Theorem~\ref{thm:pw_spectral}.
\end{proof}

\begin{remark}[Strength of Part I]
Part I is stronger than the theorem's hypotheses suggest: its derivation uses
neither associativity nor the transform-based template. Any \emph{bilinear}
product on $\bbC[G]$ whose two-sided identity is $\delta_e$ and which is
left-equivariant is already group convolution. Associativity and the template
enter only in Part II, to force the block sizes and the transform.
\end{remark}

\begin{remark}[On hypothesis (i)]
Associativity and the two-sided property of the identity are automatic for every
transform-based product (block-diagonal matrix multiplication is associative and
unital, and $M$ is a bijection); the sole content of (i) is the normalization that
the identity is the group-identity tube $\delta_e$, equivalently that $M$ is an
algebra isomorphism rather than one twisted by a unit. Without this normalization
the unit-twisted products $a\mapsto a\starG u^{-1}\starG b$ ($u$ any unit) are
associative, left-equivariant, and transform-based but are not convolution.
\end{remark}

\begin{remark}[Equivalent view via Schur]
Equation \eqref{eq:transl} says $M$ conjugates the regular representation into
$\bigoplus_s(A_s\otimes I_{d_s})$ with $A_s(g)=(M\delta_g)_s$. Since each block is
a full matrix algebra, $\{A_s(g)\}_g$ spans $M_{d_s}(\bbC)$, so $A_s$ is
irreducible (Burnside). Matching $\bigoplus_s d_s\cdot A_s$ with the regular
representation $\bigoplus_\rho d_\rho\cdot\rho$ and using uniqueness of the
isotypic decomposition (Schur's lemma) recovers $\{d_s\}=\{d_\rho\}$ and the
stated ambiguity.
\end{remark}

\begin{corollary}[Intrinsic Eckart--Young optimum]\label{cor:intrinsic}
Once the Parseval normalization of Theorem~\ref{thm:ey_full} is fixed, the residual
freedom of Theorem~\ref{thm:necessity} is restricted from the full per-irrep
$\mathrm{GL}_{d_\rho}(\bbC)$ to per-irrep unitary $P_\rho\in U(d_\rho)$ together with
permutations of equal-dimension blocks (a non-unitary change of basis would alter
the block singular values, so the invariance below is a post-normalization
statement). Under this freedom every
$\starG$-SVD singular value, each singular-tube norm $\|\mathbf s_i\|_F$, the
$\starG$-rank, and the optimal error $\|\calA-\calA_k\|_F$ are unchanged. Hence
the optimal rank-$k$ approximation is an invariant of $(G,\calA)$, independent of
the choice of irrep basis and block order (up to the usual degenerate
singular-value freedom).
\end{corollary}

\begin{proof}
A per-irrep unitary change of basis acts on each Fourier block by
$\hat\calA(:,:,\rho)\mapsto(I\otimes P_\rho)\hat\calA(:,:,\rho)(I\otimes P_\rho)^H$,
a unitary equivalence preserving singular values and Frobenius norm; block
permutations only relabel summands. The singular tubes are assembled from the
block singular values (Theorem~\ref{thm:ey_full}, Step~3), so all listed
quantities are invariant, and truncation commutes with the unitary equivalence.
\end{proof}

\begin{remark}[Eckart--Young optimality selects $F_G$]
Eckart--Young optimality is not special to $F_G$: Mor~\cite{mor2025eckart} shows
it holds for an entire family of transforms ($M=DQ$ with $Q$ unitary and $D$ a
nonzero real diagonal scaling), so best low-rank optimality alone cannot single out
any one transform. The operative content is therefore the selection principle of
Theorem~\ref{thm:necessity}: equivariance uniquely determines the group-Fourier
transform, singling out $F_G$ from that family and making the $\starG$ truncated
SVD the canonical, intrinsic optimum (Corollary~\ref{cor:intrinsic}) rather than
one arbitrary member of an Eckart--Young family.
\end{remark}

\paragraph{Relationship to prior work.}
Dunbar and Newman~\cite{dunbar2025tensor} give a representation-theoretic
\emph{interpretation} of the $\star_M$ product: through Schur's lemma and a
symmetry-adapted basis, a group representation furnishes a matrix $M$ that
block-diagonalizes the product into isotypic components, so the group-Fourier
transform is a natural and sufficient choice. Theorem~\ref{thm:necessity}
establishes the converse and sharper direction: requiring the product to be
equivariant \emph{forces} $M$ to be the non-abelian group-Fourier transform $F_G$,
unique up to a per-irrep change of basis and a permutation of equal-dimension
blocks. They answer which $M$ respects the group, by construction; we answer which
$M$ is compelled by equivariance, by necessity and uniqueness. The abelian case
recovers Jaming's characterization of the Fourier transform as the only linear map
carrying convolution to a pointwise product~\cite{jaming2009characterization}. The
block-diagonalization itself is classical (Peter--Weyl, Artin--Wedderburn); what is
new is the necessity direction and its uniqueness clause.

\paragraph{Numerical corroboration.}
We verify Theorem~\ref{thm:necessity} and Corollary~\ref{cor:intrinsic} numerically
for a panel of groups spanning the abelian and non-abelian cases (cyclic $\bbZ_4$
and $\bbZ_6$; symmetric $S_3$, dimensions $1,1,2$; quaternion $Q_8$, dimensions
$1,1,1,1,2$; octahedral $O$, dimensions $1,1,2,3,3$), in complex double precision
(\texttt{necessity\_fourier.py}). Across $200$ seeded trials per group, every
identity the theorem asserts holds to $\le 10^{-13}$: $F_G$ block-diagonalizes
convolution, the equivariant product is convolution (Part~I), and a per-irrep basis
change or an equal-dimension block permutation preserves it; by
Corollary~\ref{cor:intrinsic}, a per-irrep \emph{unitary} basis change leaves the
$\starG$-SVD singular values unchanged. Each hypothesis is confirmed load-bearing by
an $O(1)$ discriminator: the all-one-dimensional partition cannot reproduce
non-commutative convolution (commutator norm $1.41$), a cross-\emph{dimension} block
swap leaves the automorphism group and breaks both the product and equivariance
(gap $\ge 15$), the unit-twisted product with identity $\neq \delta_e$ differs from
convolution (gap $\ge 5$), and a \emph{non-unitary} per-irrep change moves the
singular values (change $O(1)$, exercised where an irrep has dimension $\ge 2$).

%% ========================================================================
\section{Extension to Space Groups: the Band-Limited \texorpdfstring{$\starSig$}{starSigma} Algebra}
\label{sec:spacegroups}
%% ========================================================================

The finite and band-limited $\starG$ calculus extends to the full symmetry of a
crystal, its space group, after one classical structural step, the Mackey
little-group construction, which fibers the calculus over the Brillouin zone. The
result is a reduction rather than new analysis: on a symmetric $k$-grid the
space-group calculus \emph{is} a finite $\starG$ algebra, so every finite-group
guarantee of this paper transfers unchanged. We treat the unitary space group; the
representation theory below (Mackey induction, the Bloch crossed-product picture,
projective little-group representations) is
classical~\cite{mackey1958unitary,bradley1972,serre1977linear} and has been
developed concurrently for the full space-group setting, including magnetic groups,
by Bekka and Brouder~\cite{bekka2025crystals}; the contribution here is the
reduction and the inherited optimality and machine verification, not the
representation theory. A neural architecture for all space groups, using
symmetry-adapted Fourier bases, was proposed concurrently~\cite{zhang2025singlearchitecture};
it enforces symmetry as an architectural \emph{invariance}, whereas the reduction
here yields a machine-verified, optimality-backed calculus whose selection rules are
\emph{enforceable} on any predictor's output (Corollary~\ref{cor:selection}) and
whose low-rank truncation is provably optimal.

Let $\Sigma$ be a $d$-dimensional space group, an extension
$1\to T\to\Sigma\xrightarrow{\pi}P\to1$ with $T\cong\bbZ^d$ the lattice
translations (abelian, normal, of finite index $|P|$) and $P$ the crystallographic
point group. We develop the symmorphic (split) case $\Sigma=T\rtimes P$; the
non-symmorphic case is Remark~\ref{rem:nonsymm}. Because $T$ is discrete, Pontryagin
duality makes its dual $\Omega:=\hat T$ a \emph{compact} torus, the Brillouin zone;
this compactness is what brings $\Sigma$ within reach of the finite theory. The
point group acts on $\Omega$ by $(p\cdot k)(t)=k(p^{-1}t)$; write
$P_k=\{p\in P:p\cdot k=k\}$ for the little co-group and $\Sigma_k=T\rtimes P_k$ for
the little group.

\begin{theorem}[Mackey little-group construction]\label{thm:mackey}
Every irreducible unitary representation of the symmorphic $\Sigma=T\rtimes P$ is
equivalent to exactly one
$\pi_{[k],\sigma}=\operatorname{Ind}_{\Sigma_k}^{\Sigma}(k\otimes\sigma)$, of
dimension $[P:P_k]\,d_\sigma$, where $[k]$ ranges over $P$-orbits in $\Omega$ and
$\sigma\in\widehat{P_k}$; distinct pairs give inequivalent irreps and these exhaust
$\hat\Sigma$.
\end{theorem}

\begin{proof}
Mackey's theorem for a regular semidirect product with abelian normal
subgroup~\cite{mackey1958unitary}, with the crystallographic specialization
in~\cite{bradley1972} and the finite instance in~\cite[\S8.2]{serre1977linear};
the action is regular because $P$ is finite, so every orbit is finite (closed) and
$\Omega/P$ is standard Borel. As a finite extension of an abelian group, $\Sigma$
is type~I~\cite{thoma1968}, so the direct-integral theory below applies.
\end{proof}

Decompose $a\in\bbC[\Sigma]$ as $a=\sum_{p\in P}a_p\otimes p$ with $a_p\in\bbC[T]$,
and let $\hat a_p(k)=\sum_{t\in T}a_p(t)e^{-ik\cdot t}$ be the lattice (Bloch)
transform.

\begin{theorem}[Bloch = crossed product]\label{thm:bloch}
The Bloch transform is a $*$-isomorphism (at the $C^*$-completion)
$\bbC[\Sigma]\cong C(\Omega)\rtimes P$ under which the $\starSig$ product is the
twisted convolution over $P$,
\begin{equation}\label{eq:twistedconv}
\widehat{(a\starSig b)}_r(k)=\sum_{p\in P}\hat a_p(k)\,
\hat b_{p^{-1}r}\!\big(p^{-1}\!\cdot k\big).
\end{equation}
Over a fundamental domain of $\Omega$ the $\starSig$ product is fiberwise
block-diagonal matrix multiplication, the fiber at $k$ carrying the little-group
irrep blocks $\bigoplus_{\sigma\in\widehat{P_k}}$.
\end{theorem}

\begin{proof}
Equation~\eqref{eq:twistedconv} is the semidirect-product multiplication
$(a_p\otimes p)(b_q\otimes q)=a_p\,(p\cdot b_q)\otimes pq$ transported through the
Bloch transform, using the character shift $\widehat{p\cdot c}(k)=\hat c(p^{-1}\!\cdot
k)$; this is the multiplication of the crossed product $C(\Omega)\rtimes P$, and
$\bbC[T\rtimes P]\cong C(\hat T)\rtimes P$ is the standard identification for an
abelian normal subgroup. Fiberwise diagonalization over each orbit is Mackey
imprimitivity, reproducing Theorem~\ref{thm:mackey}.
\end{proof}

Passing to a finite crystal, Born--von~Karman periodic boundary conditions on a
uniform $N\times\cdots\times N$ supercell replace $T$ by the finite
$T_N=(\bbZ/N)^d$ and $\Omega$ by the $P$-invariant uniform grid
$\Omega_N=\widehat{T_N}$ ($P$-invariant because $P\subset\mathrm{GL}(d,\bbZ)$
preserves $N\bbZ^d$ for every uniform $N$), yielding the \emph{finite} group
$\Sigma_N:=T_N\rtimes P$ of order $N^d|P|$.

\begin{theorem}[Reduction to finite $\starG$]\label{thm:sigmareduction}
Born--von~Karman boundary conditions replace the space group $\Sigma$ by the genuine
finite group $\Sigma_N=T_N\rtimes P$, and the induced product on $N$-periodic
signals (band-limited to $\Omega_N$) is the finite twisted convolution
$\star_{\Sigma_N}$; equivalently, in the Bloch representation the product
\eqref{eq:twistedconv} restricted to the $P$-invariant grid $\Omega_N$ is closed and
its algebra equals $C(\Omega_N)\rtimes P=\bbC[\Sigma_N]$. Because $\Omega_N$ is a
Plancherel-measure-zero finite subset of the continuous $\Omega$, this is a group
replacement, not a subalgebra of the infinite algebra. Since $\Sigma_N$ is an
ordinary finite group, the $\starG$-SVD (Theorem~\ref{thm:svd_exist}),
Eckart--Young optimality (Theorem~\ref{thm:ey_full}), canonicity
(Theorem~\ref{thm:necessity}), and the Lean-verified reductions apply verbatim to
$\bbC[\Sigma_N]$, with each irrep read as a little-group irrep $\pi_{[k],\sigma}$
($k\in\Omega_N$) of size $[P:P_k]\,d_\sigma$ and Plancherel weight
$[P:P_k]\,d_\sigma/(N^d|P|)$.
\end{theorem}

\begin{proof}
$\Omega_N=\widehat{T_N}$ and $P$-invariance make the twisted convolution
\eqref{eq:twistedconv} preserve the grid (vanishing of $\hat a_p$ off $\Omega_N$
forces the product off $\Omega_N$ to vanish), so the restricted algebra is
$C(\Omega_N)\rtimes P=\bbC[\Sigma_N]$; the block statement is
Theorem~\ref{thm:mackey} for the finite group $\Sigma_N$. Note $\Sigma_N$ is a
semidirect product, described by Theorem~\ref{thm:bloch}, not a direct product, so
the product-group ring isomorphism (Theorem~\ref{thm:prod_full}) is available for
composing $\Sigma$ with an internal symmetry but does not itself describe
$\Sigma_N$.
\end{proof}

\begin{corollary}[Selection rules across the Brillouin zone]\label{cor:selection}
A crystal tensor field transforming in a $P$-representation $\tau$ decomposes at
each $k$ into the little-group channels $\widehat{P_k}$; the allowed $\starSig$
channels are the irreps in $\tau{\downarrow}P_k$ (twisted by $k$), branching by the
subduction (compatibility) relations as $k$ moves to a lower-symmetry stratum. The
orthogonal projector onto the allowed channels at each $k$,
$\Pi(k)=\sum_{\sigma\ \mathrm{allowed}}\tfrac{d_\sigma}{|P_k|}\sum_{p\in P_k}
\overline{\chi_\sigma(p)}\,\rho(p)$, is a closed-form $\starSig$ idempotent whose
application sends every symmetry-forbidden channel to zero. Its $k=\Gamma$ instance
(little group $P$) is the point-group projector applied to cubic crystals in the
main text.
\end{corollary}

\begin{proof}
The $\Pi(k)$ are sums of mutually orthogonal isotypic
projectors~\cite[\S2.6]{serre1977linear}, hence idempotent; $P$ fixes
$\Gamma=0\in\Omega$, so $P_\Gamma=P$.
\end{proof}

\begin{proposition}[Enforcement is a guaranteed error reduction]\label{prop:enforce_guarantee}
Let $P$ be the orthogonal projector onto the symmetry-allowed subspace $S$ (the
point-group Reynolds average, Corollary~\ref{cor:selection}), and let the true
tensor satisfy $C^\star\in S$. Then for any prediction $C$,
\[
\|C-C^\star\|^2 = \|PC-C^\star\|^2 + \|(I-P)C\|^2,
\]
so $\|PC-C^\star\|\le\|C-C^\star\|$, with the error reduction equal to the
symmetry-forbidden energy $\|(I-P)C\|^2$. Enforcement therefore never increases the
error and removes exactly the forbidden channel the paper drives to machine zero;
the benefit is largest precisely when the predictor violates symmetry most, that is,
in the low-data or extrapolation regime.
\end{proposition}

\begin{proof}
$P$ is an orthogonal projection and $C^\star=PC^\star$, so $PC-C^\star=P(C-C^\star)$
and $(I-P)C=(I-P)(C-C^\star)$ are orthogonal complements of $C-C^\star$; Pythagoras
gives the identity. Empirically (\texttt{enforcement\_screening.py}), on the de Jong
tensors, which are symmetry-consistent to a residual forbidden fraction $2\times
10^{-4}$, projection reduces the reconstruction error of $98$--$100\%$ of
individual predictions across a training-fraction sweep, and composing with the
nearest-PSD step recovers every otherwise-viable material a raw Born-stability screen
rejects.
\end{proof}

\begin{corollary}[Direct-integral Eckart--Young]\label{cor:directintegral}
For a Hilbert--Schmidt tensor field over the infinite crystal there is a measurable
field of fiberwise SVDs, and for any measurable rank \emph{profile} $r(\cdot)$ the
fiberwise truncation minimizes the Plancherel-weighted Frobenius error over fields
of that profile, the error being the Brillouin-zone integral of the fiberwise
singular-value tails.
\end{corollary}

\begin{proof}
Fiberwise this is the matrix Eckart--Young theorem; measurability of a jointly chosen
SVD field and integrability of the tail are classical measurable direct-integral
facts~\cite{folland2016course,franco2024measurability}, cited not reproved. The only
$\starSig$-specific step is the fiberization, the $N\to\infty$ limit of
Theorem~\ref{thm:sigmareduction}.
\end{proof}

\begin{remark}[Rank profile versus global rank]
Corollary~\ref{cor:directintegral} is stated for a rank \emph{profile}. Converting a
single global tubal-rank budget into a per-fiber profile is ill-posed when the
singular-value field has level sets of positive Plancherel measure, so the clean
global-rank optimality of Theorem~\ref{thm:ey_full} is asserted only in the
band-limited Theorem~\ref{thm:sigmareduction}; the continuous-zone statement is
per-profile.
\end{remark}

\begin{remark}[Values, not topology]
These statements concern singular \emph{values} and irrep content, gauge-invariant at
each $k$. Any band-topological content (Chern number, Wannier localizability) resides
in the winding of the singular \emph{vectors} across $\Omega$, to which the value
spectrum is blind: two fields with identical singular-value fields can carry different
Chern number. No topological diagnostic is claimed. Concretely, a $\starSig$
certificate can guarantee the irrep content and the optimal rank of a retained
subspace but cannot certify that the subspace admits an exponentially localized
symmetric Wannier basis: that existence question is precisely the vanishing of the
topological obstruction carried by the vector field, so any Wannier-localizability
guarantee would require singular-vector (bundle) data beyond the value spectrum, and
predictions passed through the enforcement layer inherit selection-rule exactness but
no topological guarantee.
\end{remark}

\begin{remark}[Non-symmorphic groups]\label{rem:nonsymm}
When $\Sigma$ does not split, the little-group blocks become projective
representations of $P_k$ with the $k$-dependent factor system
$\omega_k(R,R')=\exp[-i(R^{-1}k-k)\cdot\tau_{R'}]$ set by the fractional translations
$\tau_R$~\cite{bradley1972}. The twisted group algebra $\bbC_{\omega_k}[P_k]$ is still
finite-dimensional and semisimple over $\bbC$ (twisted Maschke), so
Theorems~\ref{thm:bloch}--\ref{thm:sigmareduction} and
Corollary~\ref{cor:selection} hold with ``irrep'' read as ``projective irrep'': after
Born--von~Karman, $\Sigma_N$ is a genuine finite group whether or not it splits, so
the finite-group theorems and their Lean formalization apply to $\bbC[\Sigma_N]$
directly, the twisted algebra merely naming the little-co-group blocks. The cocycle is
trivial at interior $k$ and nontrivial only on parts of the zone boundary, the
algebraic origin of non-symmorphic band stickings; the same twisted-algebra mechanism,
with a different cocycle, covers spin double groups (spin--orbit coupling without time
reversal). The twist carries no computational penalty. The twisted algebra
$\bbC_{\omega_k}[P_k]$ has the same dimension $|P_k|\le|P|\le 48$ as the untwisted
one, and its projective irreps still satisfy $\sum_\sigma d_\sigma^2=|P_k|$; a
nontrivial cocycle only reallocates that fixed total into fewer, larger blocks (the
band stickings), never into a larger algebra. The $\starSig$-SVD therefore factors,
exactly as in the symmorphic case, into $N^d$ independent little-group fibers whose
block sizes are bounded by the point-group order, so runtime is linear in the
Born--von~Karman volume $N^d$ (embarrassingly parallel over $k$), memory is the
Plancherel-weighted block storage, which equals the size of the signal itself, and no
object of dimension $|\Sigma_N|$ is ever formed.
\end{remark}

\begin{remark}[Antiunitary time reversal, a flagged extension]
Time-reversal symmetry, whether in magnetic (Shubnikov) groups or in ordinary
``grey'' time-reversal-symmetric crystals, adjoins an \emph{antiunitary} operator, so
representations become Wigner corepresentations and the fiber algebras become matrix
algebras over $\bbR$, $\bbC$, or $\bbH$ (the Dyson threefold way; Herring's rules and
Kramers degeneracy). A Frobenius-norm SVD and its Eckart--Young optimum still exist
fiberwise, but the reduction to the complex-linear finite $\starG$ machinery of
Theorem~\ref{thm:sigmareduction}, and its Lean formalization, are not verbatim in the
real or quaternionic case; we record this as a flagged extension.
\end{remark}

%% ========================================================================
\section{Beyond Discrete Translations: the Wigner--Mackey Compact-Fiber Reduction}
\label{sec:wignermackey}
%% ========================================================================

The space-group reduction of Section~\ref{sec:spacegroups} kept the base compact by
discretizing the translations: with $T\cong\bbZ^d$ discrete, its dual $\Omega=\hat T$
is a compact torus, and Born--von~Karman turns the whole calculus finite. We now
remove that discreteness and treat \emph{continuous} translations $N\cong\bbR^n$, the
setting of the relativistic (Poincar\'e) and Euclidean motion groups. The dual
$\widehat N\cong\bbR^n$ (momentum space) is then non-compact, so there is no finite
grid and no global finite $\starG$ algebra. The point of this section is that the
finite/compact calculus nonetheless reaches the physically relevant content: the
non-compactness lives entirely in the \emph{base} (the orbit of momenta), while the
\emph{fiber} (the little-group content) stays compact and finite-dimensional, and the
non-compact covariance is discharged symbolically rather than numerically. The
machinery is the same Mackey construction as for crystals, with $N\rtimes H$ in place
of $T\rtimes P$.

Let $\Sigma=N\rtimes H$ be a regular semidirect product with $N\cong\bbR^n$ abelian
and $H$ acting by automorphisms, hence on $\widehat N$ by
$(h\cdot\chi)(x)=\chi(h^{-1}x)$. For $\chi\in\widehat N$ write $H_\chi=\{h\in
H:h\cdot\chi=\chi\}$ for the little group and $O=H\cdot\chi\cong H/H_\chi$ for its
orbit. Mackey's theorem for a regular semidirect product with abelian normal
subgroup~\cite{mackey1958unitary} gives the irreducibles exactly as in
Theorem~\ref{thm:mackey}, now over a continuous orbit:
$\pi_{[\chi],\sigma}=\operatorname{Ind}_{N\rtimes
H_\chi}^{\Sigma}(\chi\otimes\sigma)$ with $\sigma\in\widehat{H_\chi}$, realized on the
sections $L^2(O,\mu_O)\otimes V_\sigma$ of a homogeneous vector bundle over the orbit
with fiber the little-group irrep $V_\sigma$.

We first record that band-limiting on a compact group is not only exact on
band-limited signals (the band-limited reduction of the main text) but quantitatively
controlled on smooth ones.

\begin{proposition}[Band-limited approximation bound]\label{prop:bandlimit_error}
Let $K$ be a compact group with Casimir eigenvalue $\lambda_\rho\ge0$ on
$\rho\in\widehat K$, and let $f\in L^2(K)$ have finite Sobolev norm
$\|f\|_{H^s}^2=\sum_{\rho}(1+\lambda_\rho)^s\,d_\rho\,\|\hat f(\rho)\|_{\mathrm{HS}}^2$
for some $s>0$. Let $f_L$ be the truncation of $f$ to $\{\rho:\lambda_\rho\le L\}$.
Then
\[
\|f-f_L\|_{L^2}^2\ \le\ (1+L)^{-s}\,\|f\|_{H^s}^2 ,
\]
so the finite $\star_K$ product and its $\star_K$-SVD approximate the continuous
ones with error $O(L^{-s/2})$ (geometric in $L$ when $f$ is analytic). Consequently
the exact band-limited statement extends to a controlled approximation for smooth
signals, with the fiber truncation below incurring exactly this error.
\end{proposition}

\begin{proof}
$\|f-f_L\|^2=\sum_{\lambda_\rho>L}d_\rho\|\hat f(\rho)\|^2\le(1+L)^{-s}
\sum_{\lambda_\rho>L}(1+\lambda_\rho)^s d_\rho\|\hat f(\rho)\|^2\le(1+L)^{-s}
\|f\|_{H^s}^2$ by Parseval. Submultiplicativity of the $\starG$ product and Weyl's
singular-value perturbation inequality propagate the bound to the product and its
SVD.
\end{proof}

\begin{remark}[Truncation error is approximation error, not symmetry error]
\label{rem:truncation_exactness}
Band-limiting does not move the calculus to the approximate side of the
exact/approximate dichotomy of the main text's compounding section. The
band-limited subspace is invariant under the group action and closed under the
$\star_K$ product (the band-limited reduction), and the truncation
$f \mapsto f_L$ is itself an equivariant projection; a band-limited operator
therefore commutes with $K$ exactly, its iterates are exactly equivariant at
every depth, and no energy ever enters a symmetry-forbidden channel. What
Proposition~\ref{prop:bandlimit_error} bounds is the \emph{resolution} error
against an un-band-limited target; under iteration that error accumulates the
way any discretization error does (one application of the proof's
submultiplicativity argument per step), but it accumulates entirely within the
symmetry-allowed channels. The two axes are orthogonal: $L$ buys resolution,
never equivariance. In practice $L$ is selected from the data's spectral
tail, which is observable: the empirical Fourier energy above $L$ estimates
$\|f-f_L\|^2$ directly, so one takes the smallest $L$ whose tail energy is
below the target tolerance, against a cost that grows polynomially, as
$\sum_{\lambda_\rho\le L} d_\rho^2$ block entries (for $\mathrm{SO}(3)$,
$\sum_{l\le L}(2l+1)^2 = O(L^3)$), while the tail falls as $O(L^{-s})$ for
$H^s$ signals and geometrically for analytic ones. The lattice-gauge
demonstration below makes the trade concrete: the ground-state energy
converges geometrically in the band limit, a wider band being needed only at
stronger coupling.
\end{remark}

We instantiate this band-limited compact reduction on $\mathrm{SU}(2)$, the gauge
group of a lattice-gauge toy, entirely in the finite block domain
(\texttt{su2\_lattice\_gauge\_demo.py} in the repository). Band-limiting to $j\le2$,
including the half-integer blocks that carry the gauge charge, the electric-energy
operator is exactly the $\star_{\mathrm{SU}(2)}$ Casimir (block-diagonal, eigenvalue
$j(j+1)$ to $9\times10^{-16}$); the $\star_{\mathrm{SU}(2)}$-SVD attains the per-block
Eckart--Young error, the discarded singular-tube mass, to $4\times10^{-16}$; and the
gauge-singlet isotypic projector annihilates every non-singlet block exactly. For the
single-plaquette Kogut--Susskind Hamiltonian in its class-function sector, the
electric Casimir plus the magnetic Wilson-loop term (multiplication by the
fundamental character, which selects $\Delta j=\pm\tfrac12$ exactly), the
ground-state energy converges geometrically in the band limit: successive corrections
fall by a factor of about $0.04$ at magnetic coupling $\lambda=2$ and about $0.17$ at
$\lambda=6$, the stronger-coupling case needing a wider band, exactly the analytic
regime of Proposition~\ref{prop:bandlimit_error}.

\begin{theorem}[Wigner--Mackey compact-fiber reduction]\label{thm:wignermackey}
Fix an orbit $O\subset\widehat N$ whose little group $H_\chi$ is \emph{compact}, and
band-limit the fiber to a finite set $S\subset\widehat{H_\chi}$. Then:
\begin{enumerate}
\item at each $\chi'\in O$ the $\starSig$ product restricted to the fiber content is
the finite $\star_{H_\chi}$ blockwise matrix multiplication
$\bigoplus_{\sigma\in S}\mathrm{End}(V_\sigma)$, so the $\starG$-SVD
(Theorem~\ref{thm:svd_exist}), Eckart--Young optimality
(Theorem~\ref{thm:ey_full}), and canonicity (Theorem~\ref{thm:necessity}) apply
\emph{verbatim} to the fiber algebra;
\item the fiberwise singular values are constant along the orbit: for
$\chi'=h\cdot\chi$ the fiber operator is conjugated by the unitary $\sigma(h')$
($h'\in H_\chi$ from the Mackey cocycle), so its singular values, and hence the
$\starG$-SVD certificate, are $H$-invariant, a frame-independent scalar field.
\end{enumerate}
The numerical content of the $\starSig$-SVD over a compact-little-group orbit is
therefore the finite $\star_{H_\chi}$-SVD computed once at the base point; the
non-compact directions (the boosts moving $\chi$ along $O$) act by unitary
conjugation and add no numerical content. This is Theorem~\ref{thm:sigmareduction}
with the roles exchanged: a compact fiber over a non-compact base, in place of a
compact base with finite fibers.
\end{theorem}

\begin{proof}
Part~1 is Theorem~\ref{thm:bloch} read fiberwise: the crossed-product multiplication
$C(O)\rtimes H$ is, at each point of the orbit, the little-group convolution algebra
$\bbC[H_\chi]$, whose semisimple block form is
$\bigoplus_{\sigma\in\widehat{H_\chi}}\mathrm{End}(V_\sigma)$ by Artin--Wedderburn;
band-limiting to $S$ keeps finitely many blocks, on which the finite-group theorems
cited are exactly those of the finite $\starG$ calculus. Part~2 is the imprimitivity
conjugation of Theorem~\ref{thm:mackey}: transport along the orbit is implemented by
unitaries, which preserve singular values; it is the continuous-orbit form of the
gauge-invariance-of-values remark of Section~\ref{sec:spacegroups}.
\end{proof}

We validate part~2 numerically on the Poincar\'e massive orbit
(\texttt{wigner\_mackey\_orbit\_invariance.py} in the repository). A spin fiber,
band-limited to $j\le2$, is transported along $200$ points of the mass shell
$p^2=m^2$ by composed Lorentz boosts whose raw transport has spectral norm up to
$22$ and induces Wigner rotations of up to $0.85$~rad. The induced fiber action
nonetheless lands in the compact little group, the Wigner rotation is orthogonal
to $5\times10^{-14}$ and its spin representation unitary to $4\times10^{-16}$, so
the $\starG$-SVD singular-value field and the Eckart--Young rank-$k$ truncation
error are constant along the orbit to $7\times10^{-16}$. The identical transport
applied to a boost-covariant four-vector observable, moved by the raw non-unitary
boost, changes its singular values by a median factor of $5.8$: the numerical
signature of the residual of Remark~\ref{rem:noncompact_residual}, where
boost-covariant content is not orbit-invariant and needs the direct-integral
$\starSig$. The demonstration is thus a genuine test of the theorem's
\emph{scope}, separating the compact-fiber-reducible content from the residual,
not merely of the fact that unitaries preserve singular values.

\begin{corollary}[Linearization discharges covariance symbolically]\label{cor:linearization}
Let $F$ be a $\Sigma$-covariant map ($F(g\cdot x)=g\cdot F(x)$) and $b$ a background
with stabilizer $\Sigma_b\supseteq H_\chi$. The differential $DF_b$ intertwines the
isotropy representation of $\Sigma_b$ on the tangent space, so it (and, for an
invariant energy, its Hessian) is $H_\chi$-equivariant and is block-diagonalized
exactly by $\star_{H_\chi}$; when $H_\chi$ is compact the blocks are
finite-dimensional and the $\star_{H_\chi}$-SVD gives the optimal symmetry-preserving
low-rank model of the fluctuation operator. For a background $b'=g\cdot b$ on the same
orbit, $DF_{b'}=\rho(g)\,DF_b\,\rho(g)^{-1}$, so its spectrum and $\starG$-SVD
certificate are identical. Full $\Sigma$-covariance is thus realized by a symbolic
conjugation computed once: the non-compact (boost) covariance is discharged in the
exact symbolic layer, while the numerical $\starG$-SVD acts only on the compact
little-group fiber.
\end{corollary}

\begin{proof}
$g\cdot b=b$ for $g\in\Sigma_b$ gives $F(g\cdot x)=g\cdot F(x)\Rightarrow DF_b\,\rho(g)
=\rho(g)\,DF_b$, i.e.\ $H_\chi$-equivariance; block-diagonalization and per-block
optimality are Theorems~\ref{thm:svd_exist}--\ref{thm:ey_full}. The orbit relation is
covariance evaluated at $b'=g\cdot b$, and conjugation preserves singular values.
\end{proof}

\begin{remark}[The physical spectrum lies on compact-little-group orbits]
For the Poincar\'e group $\Sigma=\bbR^{1,3}\rtimes\mathrm{SO}(3,1)^{\uparrow}$
(or, to include half-integer spin, its universal cover
$\bbR^{1,3}\rtimes\mathrm{SL}(2,\bbC)$, for which every statement below holds
verbatim with $H_\chi=\mathrm{SU}(2)$ and the double cover of
$\mathrm{SO}(2)$ as the compact little groups) the
orbits in $\widehat N$ are the mass shells. The physical one-particle representations
sit on orbits with \emph{compact} little group: massive ($p^2=m^2>0$,
$H_\chi=\mathrm{SO}(3)$, spin) and massless with finite helicity ($p^2=0$, the
physical representations of $H_\chi=\mathrm{ISO}(2)$ factoring through the compact
$\mathrm{SO}(2)$, helicity). Theorem~\ref{thm:wignermackey} covers exactly these. The
non-compact-little-group orbits, Wigner's continuous-spin representations
(infinite-dimensional $\mathrm{ISO}(2)$) and the tachyonic $p^2<0$
($H_\chi=\mathrm{SO}(2,1)$), have never been observed in nature and fall to
the residual below. The same
organization, finite little-group (spin) labels over a continuous spectral fiber, is
the shape of the conformal bootstrap.
\end{remark}

\begin{remark}[The honest residual, and what is excluded]\label{rem:noncompact_residual}
Two regimes lie beyond the compact-fiber reduction. First, when the observable itself
transforms under the non-compact factor (a boost-covariant quantity with no fixed
frame), or when the whole non-compact block is needed numerically (conformal blocks of
$\mathrm{SO}(d,2)$ in the bootstrap; full-covariance nonlinear analysis with no
background to linearize about), no single little group suffices and one needs the
direct-integral $\starSig$ over the non-compact orbit, with Eckart--Young becoming a
generalized-$s$-number (Hilbert--Schmidt) statement. This is the non-compact-base
analogue of Corollary~\ref{cor:directintegral}, and unlike that corollary it is
genuine new analysis rather than a measurable-selection formality; we flag it as the
principal open extension. Second, genuinely infinite-dimensional symmetry groups,
$\mathrm{Diff}(M)$, the BMS group, or a gauge group taken as an
infinite-dimensional Lie group, lie outside the tensor-algebra frame entirely (no
local compactness, no Haar or Plancherel measure, as discussed in the main-text scope
analysis) and are kept symbolic or excluded. The reduction here is thus a bounded,
well-shaped extension, compact little group over a continuous fiber, not a claim over
arbitrary non-compact groups.
\end{remark}

%% ========================================================================
\section{Extended Data}
%% ========================================================================

\begin{figure}[!ht]
\centering
\includegraphics[width=0.7\textwidth]{fig_ablation.pdf}
\caption{\textbf{Extended Data Figure~1: Ablation of symmetry components.} Removing group structure from the product group experiment causes systematic degradation.}
\label{fig:ext_ablation}
\end{figure}

\begin{figure}[!ht]
\centering
\includegraphics[width=0.85\textwidth]{fig_scatter.pdf}
\caption{\textbf{Extended Data Figure~2: Predicted vs.\ true (synthetic).} (a)~$\starG$-SVD: perfect diagonal. (b)~Standard MLP: scattered.}
\label{fig:ext_scatter}
\end{figure}

\begin{table}[h]
\centering
\caption{Per-method hyperparameter settings used in all experiments. Hidden
widths $[64, 32]$, ReLU activations, He initialization, and an unregularized
bias per layer are common to all neural baselines. ``Native'' = original
training set of size $n$; ``$|G|$-aug'' = original training set augmented by
applying every $g \in G$ to each input, yielding $|G| \cdot n$ samples.}
\label{tab:hyperparams_full}
\setlength{\tabcolsep}{4pt}
\begin{tabular}{@{}p{2.8cm}p{4.0cm}p{2.6cm}p{1.6cm}p{1.1cm}p{1.4cm}p{1.3cm}@{}}
\toprule
\textbf{Method} & \textbf{Input} & \textbf{Architecture} &
\textbf{Train set} & \textbf{lr} & \textbf{Epochs} & \textbf{Batch} \\
\midrule
$\starG$-SVD + Ridge & $\starG$-features & Linear, ridge $\lambda \in \{10^{-3},\ldots,10^3\}$ (val.) & native & --  & --  & full \\
Standard MLP        & raw $\calX(:, e)$ ($z$-norm)             & $[n_f, 64, 32, 1]$  & native     & 0.003 & 300\textsuperscript{a} & 256\textsuperscript{b} \\
Invariant MLP       & $[\mathrm{mean}, \mathrm{std}, \min, \max]_g$ $\calX$ ($z$-norm) & $[4 n_f, 64, 32, 1]$ & native & 0.003 & 300 & 256 \\
Augmented MLP       & raw $\calX(:, e)$ ($z$-norm on aug.\ set) & $[n_f, 64, 32, 1]$  & $|G|$-aug & 0.003\textsuperscript{c} & 80--300\textsuperscript{d} & 256 \\
Neural $\starG$     & $\starG$-features                        & $\starG$ layers $[\cdot, 64, 32, 1]$ & native & 0.003 & 300 & 32 \\
\bottomrule
\end{tabular}
\medskip

\noindent\footnotesize
\textsuperscript{a}~80 epochs in the synthetic $\bbZ_{12}$ experiment;
300 elsewhere.\quad
\textsuperscript{b}~32 in the synthetic experiment.\quad
\textsuperscript{c}~0.005 in the synthetic experiment.\quad
\textsuperscript{d}~80 in the synthetic experiment, 200 in the
product-group experiment, 300 elsewhere; the smaller epoch budget for
heavily augmented training reflects the $|G|$-fold larger gradient budget per
epoch. Optimizer is Adam ($\beta_1 = 0.9, \beta_2 = 0.999, \varepsilon =
10^{-8}$) with early-stopping patience 20 on validation MSE for every neural
baseline.
\end{table}

\begin{table}[h]
\centering
\caption{Computational cost (wall-clock, single seed, 1{,}000 molecules).}
\begin{tabular}{@{}lccc@{}}
\toprule
\textbf{Method} & \textbf{Feature (s)} & \textbf{Train (s)} & \textbf{Total (s)} \\
\midrule
$\starG$-SVD + Ridge & 0.6 & $<$0.1 & 0.7 \\
Neural $\starG$ & 0.6 & 0.5 & 1.1 \\
Standard MLP & -- & 0.4 & 0.4 \\
Invariant MLP & -- & 0.5 & 0.5 \\
Augmented MLP & -- & 4.0 & 4.0 \\
\bottomrule
\end{tabular}
\end{table}

%% ========================================================================
\section{Wigner--Eckart Discovery: Extended Data}
%% ========================================================================

\begin{figure}[!ht]
\centering
\includegraphics[width=\textwidth]{irrep_bars.pdf}
\caption{\textbf{Extended Data Figure~4: Per-irrep predictive power.} Grouped bar chart showing $R^2$ from each irrep's features alone, for all 9 quantum properties. The qualitative pattern shift between scalar properties (A$_1$-dominated) and dipole vector components (T$_1$-dominated) is the data-driven signature of the Wigner--Eckart selection rules.}
\end{figure}

\begin{figure}[!ht]
\centering
\includegraphics[width=\textwidth]{irrep_heatmap.pdf}
\caption{\textbf{Extended Data Figure~5: Irrep decomposition heatmap.} $R^2$ for each (property, irrep) pair, sorted by tensor rank. The block structure separating rank-0 from rank-1 properties is visible as a qualitative change in the A$_1$ and T$_1$ columns across the horizontal separator.}
\end{figure}

The octahedral group $O$ was constructed programmatically from its 24 rotation matrices (6 face, 8 vertex, 6 edge rotations plus identity). The multiplication table was verified to satisfy the group axioms. The five irreducible representations were constructed as: A$_1$ (trivial), A$_2$ (determinant), T$_1$ (the rotation matrices themselves, 3D), and E + T$_2$ (from the rank-2 symmetric traceless tensor representation, decomposed into the 2D and 3D invariant subspaces). All representations were verified to be closed under the group multiplication.

\begin{algorithm}[h]
\caption{Per-irrep Fourier decomposition for Wigner--Eckart analysis}
\label{alg:per_irrep}
\begin{algorithmic}[1]
\Require feature batch $\calX \in \bbR^{N \times n_f \times |G|}$, irreps $\hat{G} = \{\rho_1, \ldots, \rho_M\}$ with dimensions $d_\rho$ and matrices $\rho(g)$, target vector $y \in \bbR^N$, ridge grid $\Lambda$
\Ensure per-irrep predictive scores $\{R^2_\rho\}_{\rho \in \hat{G}}$
\For{$\rho \in \hat{G}$}
    \For{$i = 1, \ldots, N$ and $j = 1, \ldots, n_f$ \textbf{in parallel}}
        \State $\hat{X}^\rho_{ij} \gets \sqrt{d_\rho / |G|} \sum_{g} \calX(i, j, g) \, \rho(g) \in \bbC^{d_\rho \times d_\rho}$
        \State $P^\rho_{ij} \gets \|\hat{X}^\rho_{ij}\|_F^2$ \Comment{$G$-invariant power}
    \EndFor
    \State assemble $\Phi^\rho \in \bbR^{N \times n_f}$ from $\{P^\rho_{ij}\}$
    \State split $\Phi^\rho, y$ into train/val/test (70/15/15); standardize $\Phi^\rho$
    \State $\lambda^\star \gets \arg\min_{\lambda \in \Lambda} \mathrm{MSE}_{\mathrm{val}}(\Phi^\rho, y; \lambda)$
    \State $w_\rho \gets (\Phi^{\rho \top}_{\mathrm{train}} \Phi^\rho_{\mathrm{train}} + \lambda^\star I)^{-1} \Phi^{\rho \top}_{\mathrm{train}} y_{\mathrm{train}}$
    \State $R^2_\rho \gets 1 - \mathrm{SS}_{\mathrm{res}}(\Phi^\rho_{\mathrm{test}} w_\rho, y_{\mathrm{test}}) / \mathrm{SS}_{\mathrm{tot}}(y_{\mathrm{test}})$
\EndFor
\State \Return $\{R^2_\rho\}$, $\{T_1 / A_1\text{ ratio}\}$, irrep heatmap data
\end{algorithmic}
\end{algorithm}

%% ========================================================================
\section{Baseline Implementation Algorithms}
\label{sec:baseline_alg}
%% ========================================================================

This section gives the complete per-method specifications and explicit
pseudocode for the neural baselines used in the main paper.

\paragraph{Standard MLP.} Input is the un-rotated raw feature vector
$X_{\mathrm{in}} = \calX(:, e) \in \bbR^{n_f}$ (frontal slice at the
identity), $z$-score normalized using training-set statistics. The model is
$[n_f \to 64 \to 32 \to 1]$. No symmetry information is used during training.
Trained for up to 300 epochs (synthetic and QM9 / product-group)
at learning rate $0.003$, batch size 32 (synthetic) or 256 (QM9). This is the
``no-symmetry'' control.

\paragraph{Invariant MLP.} Input is the concatenation of four hand-crafted
group-invariant pooling statistics applied along the group dimension:
$X_{\mathrm{in}} = [\,\mathrm{mean}_g \calX,\ \mathrm{std}_g \calX,\
\min_g \calX,\ \max_g \calX\,] \in \bbR^{4 n_f}$, $z$-score normalized. The
model is $[4 n_f \to 64 \to 32 \to 1]$. By construction the input is exactly
$G$-invariant so the model is invariant by composition; this is the ``manual
invariance'' control. Same optimizer schedule as the Standard MLP.

\paragraph{Augmented MLP.} Input is the un-rotated raw feature vector
$X_{\mathrm{in}} = \calX(:, e) \in \bbR^{n_f}$, but the training set is
expanded by including $X_{\mathrm{aug}} = \calX(:, g)$ for every $g \in G$,
with the same target label, yielding a $|G|$-fold augmented training set.
$z$-score statistics are computed on the augmented set. The model is
$[n_f \to 64 \to 32 \to 1]$. Test-time prediction uses the un-rotated
slice. This is the standard data-augmentation strategy and serves as a
proxy for invariance learned from data rather than enforced algebraically;
it is the closest non-equivariant baseline to an ENN. Trained for 80--300
epochs at learning rate $0.003$--$0.005$, batch size 32 (synthetic) or 256
(QM9 / product-group); the lower epoch count for the synthetic experiment
reflects the $|G|$-fold larger effective batch budget. A precise pseudocode
specification is given in Algorithm~\ref{alg:augmented_mlp}.

\paragraph{Neural $\starG$.} A symmetry-aware feed-forward network whose linear
layers are $\starG$ products with weight tensors $W^{(\ell)} \in \bbR^{n_{\ell+1}
\times n_\ell \times |G|}$ rather than ordinary matrix products. Forward
pass: $\calA^{(\ell+1)} = \mathrm{ReLU}\bigl(W^{(\ell)} \starG \calA^{(\ell)}
+ \mathbf{b}^{(\ell)}\bigr)$ for hidden layers, with a linear output and an
invariant pooling $y = \mathrm{mean}_{g, j} \calA^{(L)}(j, g)$. The hidden
widths are $[64, 32]$, matching the MLPs. Input is the invariant
$\starG$-feature vector $X_{\mathrm{in}}$ of Section~\ref{sec:feat_alg}
(Materials and Methods of the main text). Trained for 300 epochs at learning
rate $0.003$, batch size 256, with the same Adam settings and early
stopping. Equivariance is exact by construction (the per-layer rotation
variance is $\sim 10^{-28}$, at floating-point noise) so this baseline
isolates the cost of replacing a closed-form ridge regressor with a non-linear
trainable model on top of the same algebraic representation. Pseudocode is
given in Algorithm~\ref{alg:neural_starg}.

The Standard MLP and Invariant MLP follow the standard MLP training loop
that wraps Algorithm~\ref{alg:augmented_mlp} once augmentation is removed.

\begin{algorithm}[h]
\caption{Augmented MLP training}
\label{alg:augmented_mlp}
\begin{algorithmic}[1]
\Require training tensor $\calX^{\mathrm{tr}} \in \bbR^{n \times n_f \times |G|}$, target $y^{\mathrm{tr}} \in \bbR^n$, group $G$, validation $(\calX^{\mathrm{va}}, y^{\mathrm{va}})$, hidden widths $h = [64, 32]$, learning rate $\eta$, max epochs $E$, batch size $B$, patience $P$
\Ensure trained weights $W = \{W^{(1)}, W^{(2)}, W^{(3)}\}$, biases $b$
\State $\tilde{X} \gets \mathrm{reshape}(\mathrm{permute}(\calX^{\mathrm{tr}}, [1, 3, 2]), [\,n |G|, n_f\,])$ \Comment{stack all $|G|$ orbit copies}
\State $\tilde{y} \gets \mathrm{repmat}(y^{\mathrm{tr}}, |G|, 1)$ \Comment{labels are $G$-invariant; replicate}
\State $(\mu, s) \gets (\mathrm{mean}(\tilde{X}), \mathrm{std}(\tilde{X}) + 10^{-8})$ \Comment{$z$-norm on augmented set}
\State $\tilde{X} \gets (\tilde{X} - \mu) / s$
\State initialize $W^{(\ell)} \sim \mathcal{N}(0, 2 / \mathrm{fan\_in}_\ell)$ (He init); $b^{(\ell)} \gets 0$
\State Adam state: $m^{(\ell)}, v^{(\ell)} \gets 0$, step counter $t \gets 0$
\State best-$W \gets W$; $\mathrm{wait} \gets 0$; $\mathrm{best\_val} \gets +\infty$
\For{$\mathrm{epoch} = 1, \ldots, E$}
    \State shuffle $\tilde{X}$
    \For{each minibatch of size $B$}
        \State $t \gets t + 1$
        \State forward: $A^{(0)} \gets X_{\mathrm{batch}}$;\ $A^{(\ell)} \gets \mathrm{ReLU}(W^{(\ell)} A^{(\ell-1)} + b^{(\ell)})$ for $\ell < L$;\ $A^{(L)} \gets W^{(L)} A^{(L-1)} + b^{(L)}$
        \State $\mathcal{L} \gets \tfrac{1}{B} \|A^{(L)} - y_{\mathrm{batch}}\|^2$
        \State backprop $\nabla_W \mathcal{L}, \nabla_b \mathcal{L}$
        \State Adam update: $m, v$ exponential moving averages with $\beta_1 = 0.9, \beta_2 = 0.999, \varepsilon = 10^{-8}$
        \State $W^{(\ell)} \gets W^{(\ell)} - \eta \cdot \hat{m}^{(\ell)} / (\sqrt{\hat{v}^{(\ell)}} + \varepsilon)$
    \EndFor
    \State $\mathcal{L}_{\mathrm{val}} \gets$ MSE on $((\calX^{\mathrm{va}}(:, :, e) - \mu) / s, y^{\mathrm{va}})$
    \If{$\mathcal{L}_{\mathrm{val}} < \mathrm{best\_val}$}
        \State $\mathrm{best\_val} \gets \mathcal{L}_{\mathrm{val}}$;\ best-$W \gets W$;\ $\mathrm{wait} \gets 0$
    \Else
        \State $\mathrm{wait} \gets \mathrm{wait} + 1$;\ \textbf{if} $\mathrm{wait} \geq P$ \textbf{break}
    \EndIf
\EndFor
\State \Return best-$W$, best-$b$
\end{algorithmic}
\end{algorithm}

\begin{algorithm}[h]
\caption{Neural $\starG$ forward pass and gradient}
\label{alg:neural_starg}
\begin{algorithmic}[1]
\Require batch $\calX \in \bbR^{N \times n_f \times |G|}$, $\starG$-weights $\{W^{(\ell)}\}_{\ell=1}^{L}$ with $W^{(\ell)} \in \bbR^{n_{\ell+1} \times n_\ell \times |G|}$, biases $\{b^{(\ell)}\} \in \bbR^{n_{\ell+1}}$ (one scalar per output channel, broadcast constant along the group axis), group $G$
\Ensure scalar predictions $\hat{y} \in \bbR^N$
\State $\calA^{(0)} \gets \calX$
\For{$\ell = 1, \ldots, L$}
    \State $\calZ^{(\ell)} \gets W^{(\ell)} \starG \calA^{(\ell-1)} + b^{(\ell)}$ \Comment{Algorithm~\ref{alg:starg_product}}
    \If{$\ell < L$}
        \State $\calA^{(\ell)} \gets \mathrm{ReLU}(\calZ^{(\ell)})$
    \Else
        \State $\calA^{(\ell)} \gets \calZ^{(\ell)}$ \Comment{linear output}
    \EndIf
\EndFor
\State $\hat{y}_i \gets \tfrac{1}{n_L |G|} \sum_{j, g} \calA^{(L)}(i, j, g)$ \Comment{$G$-invariant pooling}
\State \Return $\hat{y}$
\end{algorithmic}
\end{algorithm}

The network is trained for 300 epochs (Adam, $\eta = 0.003$, batch size 32,
early-stopping patience 20), backpropagating through the $\starG$ product
(Algorithm~\ref{alg:starg_product}) layer by layer. Equivariance is preserved
exactly, to floating-point precision, for \emph{any} trained weights: the $\starG$
product factors through the per-irrep block multiplication, and the bias is
constrained to the trivial isotypic component (constant along the group axis), so
it acts only on the invariant channel and commutes with the group action. A free
per-group-element bias, by contrast, would break equivariance as soon as training
moved it off its zero initialization; the group-constant (trivial-isotypic) bias is
what keeps the trained network exactly equivariant.

\subsection{ENN baselines: SchNet, e3nn, MACE}
\label{sec:enn_baselines}

The three ENN baselines used in Section~2.9 are not reimplemented from
scratch: each is a published reference implementation, executed on the
same train/val/test splits and the same seeds as the $\starG$ and MLP
baselines. The configuration we used is recorded explicitly so that the
comparison is reproducible.

\begin{itemize}
\item \textbf{SchNet}~\cite{schutt2017schnet}.
  Reference implementation: \texttt{schnetpack} v2.0.4 (pinned in the
  repository's \texttt{requirements.txt}).
  Configuration: $128$ atom-basis features, $6$ interaction blocks,
  $20$-Gaussian radial basis, cosine cutoff at $5.0$\,\AA, MSE loss with
  L1 monitoring, Adam $\eta = 5{\times}10^{-4}$, batch $64$, max $200$
  epochs, early-stop patience $20$. We run the standard
  \texttt{spk.datasets.QM9} loader with \texttt{remove\_uncharacterized
  =True} to match PyG's $130{,}831$-molecule subset, the
  \texttt{ASENeighborList(cutoff=5.0)} transform, and a per-target
  $z$-norm via \texttt{RemoveOffsets}/\texttt{AddOffsets}.
  Scalar targets only ($\mu$, $\alpha$, gap, ZPVE).

\item \textbf{e3nn-based SE(3)-equivariant model}~\cite{thomas2018tensor,geiger2022e3nn}.
  A compact equivariant message-passing network built directly from
  \texttt{e3nn} v0.5.4 primitives. Three layers, hidden irreps
  \texttt{32x0e+16x1o+8x2e}, edge spherical harmonics
  \texttt{1x0e+1x1o+1x2e}, RBF $16$ Gaussians on the $0\!-\!5$\,\AA{}
  cutoff, gated equivariant non-linearities, sum-pool over atoms,
  $\texttt{FullyConnectedTensorProduct}$ head with output irreps matched
  to target rank (\texttt{1x0e} for scalars, \texttt{1x1o} for
  $\boldsymbol{\mu}$ vector, \texttt{1x2e+1x0e} for $\alpha$ tensor).
  Adam $\eta = 5{\times}10^{-4}$, batch $32$, $200$ epochs, patience
  $20$. Used for tensor-rank-matched comparison rather than as the SOTA
  ENN target.

\item \textbf{MACE}~\cite{batatia2022mace}.
  Reference implementation: \texttt{mace-torch} v0.3.15 (pinned in
  \texttt{requirements.txt}). Configuration:
  \texttt{ScaleShiftMACE} with $r_{\max} = 5.0$\,\AA, $8$ Bessel radial
  features, $5$-th order polynomial cutoff, $\ell_{\max}=3$,
  correlation $3$, two interaction blocks
  (\texttt{RealAgnosticInteractionBlock} first;
  \texttt{RealAgnosticResidualInteractionBlock} second), hidden irreps
  \texttt{128x0e+128x1o}, MLP irreps \texttt{16x0e}, $5$ elements
  (H/C/N/O/F), per-target shift/scale $(\bar y, \mathrm{std}\,y)$.
  Optimizer: Adam (\texttt{amsgrad}), $\eta=10^{-3}$, batch $32$,
  \texttt{ReduceLROnPlateau} (factor $0.5$, patience $15$), max $200$
  epochs, early-stop patience $25$ on validation MSE. Total parameter
  count: $945{,}168$. Scalar targets only.
\end{itemize}

\noindent
\emph{Implementation notes.} Three non-trivial integration adjustments
were required that may be useful to anyone reproducing the comparison.
(i)~In \texttt{schnetpack}'s \texttt{ModelOutput}, every metric must be
an \texttt{nn.Module} (lambda functions are silently rejected by
\texttt{nn.ModuleDict}); we use only \texttt{L1Loss} as the metric and
recompute RMSE/$R^2$ ourselves from saved test predictions.
(ii)~In \texttt{mace-torch} ${\geq}\,0.3.10$, \texttt{interaction\_cls}
and \texttt{interaction\_cls\_first} are mandatory, and
\texttt{hidden\_irreps}/\texttt{MLP\_irreps} must be \texttt{o3.Irreps}
objects, not strings; per-molecule \texttt{AtomicData} construction
requires a single shared \texttt{z\_table} for the full element set
(H/C/N/O/F) so that the per-graph \texttt{node\_attrs} have uniform
width when batched. (iii)~Loading \texttt{e3nn 0.4.4}'s pickled
\texttt{constants.pt} (transitively imported by \texttt{mace-torch})
fails under PyTorch ${\geq}\,2.6$'s default
\texttt{weights\_only=True}; we set
\texttt{TORCH\_FORCE\_NO\_WEIGHTS\_ONLY\_LOAD=1} for the comparison
runs only.

%% ========================================================================
\section{Angular Resolution: Octahedral versus Icosahedral}
\label{sec:ico_resolution}

The main text notes that the quantitative accuracy ceiling of a finite point
group is set by its angular resolution, and that larger polyhedral groups lift
that ceiling.  We demonstrate this directly in a controlled setting that does
not depend on rare molecular symmetry (\texttt{icosahedral\_resolution.py}).
We draw random functions band-limited to angular degree $\ell \le \ell_{\max}$
on the rotation group, observe each function only on the finite orbit of a
generic base direction under a point group $G$ (24 directions for the chiral
octahedral group $O$, 60 for the chiral icosahedral group $I$), fit a
harmonic model on that orbit, and measure the test $R^2$ on a dense set of
$4{,}000$ held-out rotations.  A test $R^2$ below one is aliasing: the finite
orbit cannot resolve the continuous signal.

\begin{table}[h]
\centering
\caption{Angular resolution by group: test $R^2$ on $4{,}000$ held-out
rotations for random signals band-limited to degree $\ell_{\max}$ (mean over
300 signals; seed 0).  The chiral octahedral group ($24$ frames) resolves
signals through $\ell=3$ but aliases $\ell=4$; the chiral icosahedral group
($60$ frames) resolves $\ell=4$ exactly.  The $\ell=4$ gap is stable across
seeds (per-seed octahedral $R^2$ = 0.908 / 0.943 / 0.896 for seeds 0/1/2,
each a mean over 300 signals).}
\label{tab:ico_resolution}
\begin{tabular}{@{}ccc@{}}
\toprule
\textbf{Signal degree $\ell_{\max}$} & \textbf{Octahedral $O$ (24 frames)} & \textbf{Icosahedral $I$ (60 frames)} \\
\midrule
1 & $1.000$ & $1.000$ \\
2 & $1.000$ & $1.000$ \\
3 & $1.000$ & $1.000$ \\
4 & $0.908$ & $1.000$ \\
\bottomrule
\end{tabular}
\end{table}

The pattern matches representation theory.  The octahedral group's largest
irreducible representation has dimension $3$ ($T_1$, $T_2$), so a generic
$24$-point orbit cannot linearly separate the nine degree-$4$ harmonics from
lower degrees and aliases them; the icosahedral group has a dimension-$5$
irreducible representation ($H$), and its $60$-point orbit resolves degree $4$
completely.  This confirms, quantitatively, that moving from $O$ to $I$ lifts
the angular-resolution ceiling, and it exercises the machine-validated
icosahedral $\starG$ algebra (\texttt{starg\_torch/icosahedral.py}: five
irreps $A, T_1, T_2, G, H$ with $\sum_\rho d_\rho^2 = 60$, group-Fourier
matrix unitary to $2 \times 10^{-15}$, $\starG$-SVD reconstruction to
$8 \times 10^{-6}$).

\section{Error Metrics in Physical Units and Literature Anchoring}
\label{sec:mae_anchor}

The main text reports $R^2$ throughout because the contribution of the
$\starG$ rows is representational (closed form, exact equivariance, per-irrep
attribution) rather than benchmark accuracy, and $R^2$ makes the per-irrep
attribution comparable across targets with different units.  For
cross-referencing with the QM9 literature, which conventionally reports mean
absolute error (MAE) in physical units, Table~\ref{tab:mae_anchor} restates
the full-QM9 results of the main text as MAE, and places our self-trained
SchNet baseline next to the published SchNet
figures~\cite{schutt2018schnet}.  Two facts are visible.  First, our
self-trained baselines are healthy: our SchNet is comparable to the published
SchNet, improving on the reported MAE for gap and $\alpha$ and lying within a
small factor for $\mu$ and ZPVE, so the ENN baselines in this paper are
trained to community standard.  Second, the $\starG$ absolute errors sit
far above the message-passing state of the art, exactly as the pooled-$R^2$
concession in the main text states; the $\starG$ value proposition is the
closed-form equivariant representation at roughly a hundred parameters, not
benchmark accuracy.

\begin{table}[h]
\centering
\caption{Full QM9 test MAE in physical units (mean $\pm$ std; same runs as the
main-text $R^2$ table: 5 seeds for the $\starG$ and SchNet rows, 3 for MACE).
Published SchNet values from
\cite{schutt2018schnet}.}
\label{tab:mae_anchor}
\begin{tabular}{@{}lcccc@{}}
\toprule
\textbf{Method} & \textbf{gap (eV)} & \textbf{$\alpha$ ($a_0^3$)} & \textbf{$\mu$ (D)} & \textbf{ZPVE (eV)} \\
\midrule
$\starG$-SVD + Ridge & $0.746 \pm 0.001$ & $1.746 \pm 0.045$ & $0.782 \pm 0.006$ & $0.029 \pm 0.000$ \\
$\starG$ neural      & $0.671 \pm 0.004$ & $1.333 \pm 0.059$ & $0.690 \pm 0.013$ & $0.027 \pm 0.002$ \\
SchNet (this work)   & $0.055 \pm 0.002$ & $0.127 \pm 0.069$ & $0.060 \pm 0.041$ & $0.0019 \pm 0.0002$ \\
SchNet (published)~\cite{schutt2018schnet} & $0.063$ & $0.235$ & $0.033$ & $0.0017$ \\
MACE (this work)     & $0.116 \pm 0.005$ & $0.317 \pm 0.004$ & $0.060 \pm 0.003$ & $0.0097 \pm 0.0004$ \\
\bottomrule
\end{tabular}
\end{table}

\section{Molecular Accuracy Benchmarks on QM9}
\label{sec:qm9_accuracy}

The main text (Scope: pooled predictive accuracy) concedes pooled predictive
accuracy on QM9 and confines the details to this section: the full-QM9
head-to-head against graph equivariant networks and matched-input MLP
baselines, the within-isomer audit that attributes most of the graph-network
margin to an input-representation gap, and the data-scarce sample-size sweep.
None of these numbers carries a capability claim of the paper; they delimit
its scope.

\paragraph{Full-QM9 head-to-head and matched-input comparison.}
Table~\ref{tab:fullqm9} reports test $R^2$ on the four scalar QM9 targets
(130{,}831 characterized molecules; 91{,}581 / 19{,}624 / 19{,}626
train / validation / test split; 5 seeds, 3 for MACE). The five molecule-level
methods (both $\starG$ variants and the three MLP baselines) consume the
\emph{identical} $48$-row angular feature tensor, carrying angular moments,
heavy-atom rows, atom-pair Coulomb couplings, and distance-distribution rows
under cyclic Z$_{12}$, and every MLP is trained with a modern recipe (target
standardization, validation-tuned weight decay, learning-rate schedule;
Methods of the main text), so that comparison is genuinely matched-input.
On it, plain standard and invariant MLPs exceed both $\starG$ variants on
gap, $\alpha$, and $\mu$ and tie on ZPVE; on shared input a well-trained
unconstrained network is the stronger pooled predictor. SchNet and MACE
consume the full atomic graph, a different-input comparison that sets the
pooled ceiling. Parameter counts compare only the trained predictor; the
$\starG$ features are a fixed, untrained preprocessing step.

\begin{table}[h]
\centering
\caption{Full QM9, four scalar targets, mean test $R^2 \pm$ std
(130{,}831 molecules; 91{,}581 / 19{,}624 / 19{,}626 train / val
/ test split; $\starG$, MLP, and SchNet rows 5 seeds, MACE row 3 seeds).
The five molecule-level methods consume the identical $48$-row angular
feature tensor (matched input); SchNet and MACE consume the full atomic
graph and set the pooled-$R^2$ ceiling, motivating the within-isomer audit in
Table~\ref{tab:isomer_audit}. The augmented MLP collapses regardless
of capacity (Section~\ref{sec:aug_capacity}).
$^{\dagger}$SchNet uses the PyTorch Geometric reference implementation
(\texttt{torch\_geometric.nn.models.SchNet}) with $128$ hidden channels,
$6$ interactions, $50$-Gaussian RBF, cutoff $10$\,\AA, $455{,}809$
trainable parameters. Full configurations for all three ENN baselines are in
Section~\ref{sec:enn_baselines}.}
\label{tab:fullqm9}
\begin{tabular}{@{}lrcccc@{}}
\toprule
\textbf{Method} & \textbf{Params} & \textbf{gap} & \textbf{alpha} & \textbf{mu} & \textbf{ZPVE} \\
\midrule
$\starG$-SVD + Ridge & $144$    & $0.481 \!\pm\! 0.002$ & $0.907 \!\pm\! 0.003$ & $0.472 \!\pm\! 0.006$ & $0.998 \!\pm\! 0.000$ \\
$\starG$ neural       & $62{,}625$ & $0.571 \!\pm\! 0.006$ & $0.945 \!\pm\! 0.005$ & $0.563 \!\pm\! 0.016$ & $0.998 \!\pm\! 0.002$ \\
MLP standard         & 5{,}249  & $0.687 \!\pm\! 0.008$ & $0.973 \!\pm\! 0.001$ & $0.672 \!\pm\! 0.007$ & $0.999 \!\pm\! 0.000$ \\
MLP invariant        & 14{,}465  & $0.686 \!\pm\! 0.004$ & $0.972 \!\pm\! 0.001$ & $0.642 \!\pm\! 0.009$ & $0.999 \!\pm\! 0.000$ \\
MLP augmented        & 5{,}249  & $0.002 \!\pm\! 0.002$ & $0.002 \!\pm\! 0.001$ & $0.001 \!\pm\! 0.001$ & $0.003 \!\pm\! 0.003$ \\
SchNet$^{\dagger}$   & $455{,}809$  & $0.996 \!\pm\! 0.000$ & $0.999 \!\pm\! 0.001$ & $0.995 \!\pm\! 0.005$ & $1.000 \!\pm\! 0.000$ \\
MACE                 & $945{,}168$ & $0.985 \!\pm\! 0.001$ & $0.997 \!\pm\! 0.000$ & $0.995 \!\pm\! 0.001$ & $1.000 \!\pm\! 0.000$ \\
\bottomrule
\end{tabular}
\end{table}

\paragraph{Within-isomer audit.}
On QM9 the HOMO--LUMO gap, the polarizability $\alpha$, and the dipole
magnitude $\mu$ are strongly size-correlated, so a model that captures gross
size and composition already explains the bulk of the pooled variance. The
within-isomer $R^2$ (restricted to molecular formulas with $\geq 5$
constitutional isomers, sample-weighted across formulas, same test
predictions) removes that signal (Table~\ref{tab:isomer_audit}). Every method
on the molecule-level summary caps near within-isomer $R^2 \approx 0.35$
(the nonlinear MLPs), with the closed-form linear $\starG$-SVD at $0.01$,
whereas SchNet ($0.991$) and MACE ($0.963$) consume the full atomic graph
carrying the bond topology that distinguishes constitutional isomers. Most of
the pooled-$R^2$ advantage of the graph networks is therefore an
input-representation gap rather than a purely algorithmic one; the per-atom
$\starG$ descriptor of the main text (Symmetry in the representation, not the
architecture) reaches within-isomer $R^2 = 0.785$ with a symmetry-blind
learner, confirming that the plateau is a property of the molecule-level
featurization.

\begin{table}[h]
\centering
\caption{Pooled vs.\ within-isomer test $R^2$ for the HOMO--LUMO gap
target. Within-isomer $R^2$ is computed over molecular formulas with
$\geq 5$ constitutional isomers and reported as the sample-weighted
mean across formulas. The collapse from pooled to within-isomer
quantifies how much of the headline $R^2$ is size-prediction signal.}
\label{tab:isomer_audit}
\begin{tabular}{@{}lcc@{}}
\toprule
\textbf{Method} & \textbf{Pooled $R^2$} & \textbf{Within-isomer $R^2$} \\
\midrule
$\starG$-SVD + Ridge & $0.481$  & $\phantom{-}0.011$ \\
$\starG$ neural       & $0.571$  & $\phantom{-}0.101$ \\
MLP standard         & $0.687$  & $\phantom{-}0.350$ \\
MLP invariant        & $0.686$  & $\phantom{-}0.345$ \\
MLP augmented        & $0.002$  & $-2.119$ \\
SchNet               & $0.996$ & $\phantom{-}0.991$ \\
MACE                 & $0.985$ & $\phantom{-}0.963$ \\
\bottomrule
\end{tabular}
\medskip

\footnotesize
\noindent\emph{Within-isomer $R^2$ averaged over $\sim$230 molecular
formulas (per seed, $\geq 5$ constitutional isomers each); covered
test molecules range from $\sim$10{,}500 (PyG-split SchNet) to
$\sim$19{,}250 ($\starG$ / MLP / MACE). The sample-weighted mean
across formulas is robust to the per-method test-set partitioning,
so the numerical comparison is unaffected. }
\end{table}

\paragraph{Data-scarce sample-size sweep.}
With $\bbZ_{12}$ rotational featurization and the same fair recipe for every
method, a sweep from $100$ to $1{,}000$ training molecules (5 seeds) shows a
regime crossover (Table~\ref{tab:qm9}, Figure~\ref{fig:learning_curves}): at
$100$ molecules the closed-form $\starG$-SVD ridge regressor leads every
trained baseline ($R^2 = 0.415$ at $64$ parameters, versus $0.199$ for the
standard MLP), and by $1{,}000$ molecules the unconstrained MLPs overtake
($0.558$ versus $0.493$); the crossover falls between $316$ and $1{,}000$
training molecules (Figure~\ref{fig:learning_curves}). The advantage is real
but regime-specific, largest where data is scarcest. The augmented MLP never becomes competitive at any
sample size, consistent with its full-scale collapse
(Section~\ref{sec:aug_capacity}).

\begin{table}[h]
\centering
\caption{QM9 HOMO--LUMO gap ($\bbZ_{12}$), test $R^2$ (mean $\pm$ std over 5
seeds) at the smallest and largest sample sizes of the sweep, with the same
fair training recipe for all methods. $\starG$-SVD leads at $n=100$; the
unconstrained MLPs overtake by $n=1{,}000$.}
\label{tab:qm9}
\begin{tabular}{@{}lccc@{}}
\toprule
\textbf{Method} & \textbf{Params} & \textbf{$R^2$ ($n{=}100$)} & \textbf{$R^2$ ($n{=}1{,}000$)} \\
\midrule
$\starG$-SVD + Ridge & \textbf{64}    & $\mathbf{0.415 \pm 0.187}$ & $0.493 \pm 0.095$ \\
$\starG$ neural       & $36{,}513$ & $0.074 \pm 0.196$          & $0.479 \pm 0.108$ \\
MLP standard         & $5{,}249$  & $0.199 \pm 0.176$          & $\mathbf{0.558 \pm 0.110}$ \\
MLP invariant        & $5{,}761$  & $-0.041 \pm 0.367$         & $0.505 \pm 0.094$ \\
MLP augmented        & $5{,}249$  & $-0.726 \pm 1.209$         & $0.025 \pm 0.027$ \\
\bottomrule
\end{tabular}
\end{table}

\begin{figure}[!ht]
\centering
\includegraphics[width=0.72\textwidth]{fig3d_learning_curves.pdf}
\caption{\textbf{Data-scarce learning curves on the QM9 HOMO--LUMO gap
($\bbZ_{12}$), all methods trained with the same fair recipe (5 seeds; bands
$\pm 1$ s.d.).} The closed-form $\starG$-SVD ridge regressor ($64$ parameters)
is the most sample-efficient method in the very-low-data regime, leading at
$100$--$316$ molecules, where the unconstrained MLPs overfit their thousands of
parameters. The standard and invariant MLPs catch up and overtake $\starG$ by
$1{,}000$ molecules; the augmented MLP never becomes competitive. The
$\starG$ advantage is real but regime-specific, concentrated where data is
scarcest.}
\label{fig:learning_curves}
\end{figure}

\begin{figure}[!ht]
\centering
\includegraphics[width=\textwidth]{pareto_gap.pdf}
\caption{\textbf{Parameter-efficiency vs predictive power on QM9 HOMO--LUMO
gap.}
(\textbf{a}) Pooled test $R^2$ vs trainable parameters (error bars $\pm$std).
MACE occupies the upper-right ($R^2 = 0.985$ at $945{,}168$ parameters);
$\starG$-SVD~+~Ridge occupies the far left at $144$ parameters
($R^2 = 0.481$), roughly $6{,}600\times$ fewer parameters than MACE, and the
matched-input MLPs sit above it ($R^2 \approx 0.69$ at a few thousand
parameters). The augmented MLP sits at the bottom ($R^2 \approx 0.00$),
illustrating the structural collapse of orbit-augmented learning at
full QM9 scale.
(\textbf{b}) Within-isomer test $R^2$ on the same axes (formulas with
$\geq 5$ constitutional isomers, sample-weighted mean). Every
molecule-level-summary method is capped well below the graph networks:
$\starG$-SVD at $\approx 0.01$, $\starG$ neural at $\approx 0.10$, and the
nonlinear MLPs at $\approx 0.35$, versus SchNet ($\approx 0.99$) and MACE
($\approx 0.96$). The large offset to the ENNs is the input-information gap
(bond topology in the full graph); the spread within the molecule-summary
methods reflects model nonlinearity.}
\label{fig:pareto}
\end{figure}

\section{Augmented-MLP Capacity Sweep}
\label{sec:aug_capacity}

The main text reports that the Augmented MLP collapses at full QM9 scale
($R^2 \approx 0.002$ on gap, $\alpha$, $\mu$) even on the matched $48$-row
feature tensor and with the modern training recipe.  To confirm this is not a
capacity limitation, we sweep the hidden width from the default $[64,32]$ up
to $[1024,512]$, a $110\times$ increase in parameters, holding the fair recipe
fixed (three seeds; \texttt{submit\_aug\_capacity.bsub}).

\begin{table}[h]
\centering
\caption{Augmented MLP at full QM9 on the matched $48$-row features: test
$R^2$ versus capacity (mean over 3 seeds). More capacity does not recover the
collapse; the standard MLP on the identical features reaches $0.687$ (gap) and
$0.973$ ($\alpha$).}
\label{tab:aug_capacity}
\begin{tabular}{@{}lrcc@{}}
\toprule
\textbf{Hidden widths} & \textbf{Params} & \textbf{gap $R^2$} & \textbf{$\alpha$ $R^2$} \\
\midrule
$[64, 32]$     & $5{,}249$   & $0.002$ & $0.002$ \\
$[256, 128]$   & $45{,}569$  & $0.009$ & $0.021$ \\
$[512, 256]$   & $156{,}673$ & $0.016$ & $0.020$ \\
$[1024, 512]$  & $575{,}489$ & $0.011$ & $0.031$ \\
\bottomrule
\end{tabular}
\end{table}

Across a $110\times$ parameter range the augmented model never exceeds
$R^2 = 0.031$, two orders of magnitude below the standard MLP on identical
input.  Orbit augmentation presents the network with $|G|$ distinct rotated
feature vectors mapped to one label; absent an architectural invariance, the
model regresses toward the label mean rather than learning the invariant, and
capacity does not change this.  This is the concrete failure that algebraic
equivariance avoids by construction.

\section{Probing Trained ENN Internals per Irrep}
\label{sec:enn_probe}

The main text states that standard ENN workflows do not \emph{expose} a
per-irrep predictive decomposition, not that the information is absent from
their learned representations.  To make that distinction concrete we probe
the trained e3nn baseline directly.  The e3nn hidden representation is
\texttt{32x0e + 16x1o + 8x2e}; for each molecule we form a rotation- and
permutation-invariant per-multiplet power (squared multiplet norm summed over
atoms), mirroring the per-irrep Fourier power that $\starG$ uses, then fit
RidgeCV per $\ell$-block and report test $R^2$
(\texttt{enn\_irrep\_probe.py}).

\begin{table}[h]
\centering
\caption{Per-$\ell$ probe of the trained e3nn baseline (seed 0): test $R^2$
of RidgeCV on invariant block powers of the hidden representation. The pattern is
reproduced across three seeds (gap and ZPVE): the scalar block stays dominant, with
$\ell=1,2$ recoverable but weaker.}
\label{tab:enn_probe}
\begin{tabular}{@{}lcccc@{}}
\toprule
\textbf{Target} & \textbf{$\ell=0$ (scalar)} & \textbf{$\ell=1$ (vector)} & \textbf{$\ell=2$ (rank-2)} & \textbf{all $\ell$} \\
\midrule
gap   & 0.944 & 0.625 & 0.413 & 0.948 \\
alpha & 0.974 & 0.772 & 0.751 & 0.983 \\
ZPVE  & 0.990 & 0.904 & 0.758 & 0.996 \\
\bottomrule
\end{tabular}
\end{table}

The probe shows that a trained ENN's internals do carry $\ell$-resolved
structure (the scalar block dominates for scalar targets, as expected), and
that it can be recovered by a bespoke post-hoc analysis of a trained model.
The contrast with $\starG$ is therefore precise, and it is threefold. The
probe presupposes the trained network and inherits its contingencies: the
recovered per-$\ell$ attribution is a property of one converged model (it
varies with seed and training budget) and carries no optimality or exactness
certificate. It is read-only: a diagnostic of what the model happened to
learn, with no mechanism for restricting or enforcing a channel. And it is
bespoke to one architecture's internal irrep typing. The $\starG$
decomposition is canonical for the group (the canonicity theorem), costs one
closed-form transform with no training, certifies its selection rules
exactly, and doubles as an input-side control: channels can be restricted,
enforced, or ablated \emph{before} any learner is trained.

\section{End-to-End Workflow}
\label{sec:workflow}
%% ========================================================================

The complete pipeline used to produce every numerical result in this paper
is summarized in Algorithm~\ref{alg:pipeline}. The pipeline is identical
across the synthetic, QM9, product-group, symmetry-discovery, and
Wigner--Eckart experiments; only the group $G$, the featurization
$\phi: \mathrm{molecule} \mapsto \calX$, and the regression head differ.

\begin{algorithm}[h]
\caption{End-to-end $\starG$-SVD~+~ridge pipeline}
\label{alg:pipeline}
\begin{algorithmic}[1]
\Require dataset $\{(\mathrm{mol}_i, y_i)\}_{i=1}^{N}$, group $G$, featurizer $\phi$, ridge grid $\Lambda$
\Ensure trained predictor $\hat{f}$, test scores
\State precompute $\calT_G$, $F_G$, irrep dimensions $\{d_\rho\}$
\For{$i = 1, \ldots, N$}
    \State $\calX_i \gets \phi(\mathrm{mol}_i; G)$ \Comment{tensorial featurization}
\EndFor
\State assemble $\calX \in \bbR^{N \times n_f \times |G|}$
\State split $\calX, y$ into train/val/test (70/15/15)
\State $(\Phi_{\mathrm{tr}}, \Theta) \gets$ Algorithm~\ref{alg:features}$(\calX_{\mathrm{tr}}, G)$
\State $\Phi_{\mathrm{va}}, \Phi_{\mathrm{te}} \gets$ Algorithm~\ref{alg:features}$(\cdot, G; \Theta)$
\State $\lambda^\star \gets \arg\min_{\lambda \in \Lambda} \mathrm{MSE}_{\mathrm{val}}(\Phi_{\mathrm{va}}, y_{\mathrm{va}}; \lambda)$
\State $w \gets (\Phi_{\mathrm{tr}}^\top \Phi_{\mathrm{tr}} + \lambda^\star I)^{-1} \Phi_{\mathrm{tr}}^\top y_{\mathrm{tr}}$
\State $R^2_{\mathrm{te}} \gets 1 - \mathrm{SS_{res}}(\Phi_{\mathrm{te}} w, y_{\mathrm{te}}) / \mathrm{SS_{tot}}(y_{\mathrm{te}})$
\State $\nu \gets \mathrm{Var}_{g \in G} \hat{y}(g \cdot \calX_{\mathrm{te}})$ \Comment{rotation-variance audit}
\State \Return $\hat{f}: \calX \mapsto \mathrm{Algorithm~\ref{alg:features}}(\calX; \Theta) \cdot w$, $R^2_{\mathrm{te}}$, $\nu$
\end{algorithmic}
\end{algorithm}

%% ========================================================================
\section{Formal Verification in Lean~4}
\label{sec:lean}
%% ========================================================================

The core algebraic results in this paper have been machine-verified in the
Lean~4 proof assistant~\cite{demoura2021lean4} using the Mathlib
library~\cite{mathlib2020}.  The formalization comprises roughly $2{,}200$
lines of Lean~4 across twelve modules, with zero unresolved proof obligations
(\texttt{sorry}).  It certifies the $\starG$ algebra laws, equivariance and
the invariance identities, the $\starG$-SVD and its Eckart--Young optimality,
the further matrix-mimetic operations ($\starG$-QR, symmetric
eigendecomposition, and least-squares), the product-group factorization, the
octahedral selection rules, and the concrete witness systems; the axiom
dependencies of every theorem are
exhibited explicitly, so the certificate is exactly as strong as the declared
axiom budget it rests on: six axioms in the real development, three
complex-local analogues confined to the complex module's cyclic witnesses,
and, for the \texttt{native\_decide} selection-rule computations only, Lean's
compiled-evaluation axiom \texttt{Lean.ofReduceBool}.

\subsection{Architecture}

The formalization is organized into twelve modules mirroring the paper's
logical structure (line counts from the repository at submission):

\begin{center}
\resizebox{\textwidth}{!}{
\begin{tabular}{@{}llrl@{}}
\toprule
\textbf{Module} & \textbf{Paper section} & \textbf{Lines} & \textbf{Content} \\
\midrule
\texttt{Basic.lean}         & SI~\S\S1--4   & 58  & Convolution tensor, $\starG$ product, transpose \\
\texttt{Algebra.lean}       & SI~\S4        & 92 & Associativity, distributivity, identity laws \\
\texttt{Equivariance.lean}  & SI~\S8        & 139 & Equivariance, Frobenius/Fourier-power invariance \\
\texttt{ProductGroup.lean}  & Theorem~2     & 146 & Product-group factorization, Kronecker irreps \\
\texttt{EckartYoung.lean}   & Theorem~1     & 183 & Matrix Eckart--Young, proved from two bridge lemmas \\
\texttt{SVD.lean}           & Theorem~1     & 228 & Irrep systems, $\starG$-SVD, Eckart--Young optimality \\
\texttt{MatrixMimetic.lean} & Methods       & 260 & $\starG$-QR, spectral, least-squares; transpose bridge lemma \\
\texttt{WignerEckart.lean}  & \S2.2         & 66  & Octahedral characters, derived selection rules \\
\texttt{ConcreteWitness.lean} & SI~\S\ref{sec:lean} & 163 & Axiom-free $C_2$ irrep system; optimality suite instantiated \\
\texttt{OctahedralWitness.lean} & \S2.3     & 396 & Four of five octahedral irreps, axiom-free, with irreducibility \\
\texttt{ComplexCore.lean}   & SI~\S\ref{sec:lean} & 390 & Complex-typed layer; cyclic (DFT) witnesses; convolution theorem \\
\texttt{Audit.lean}         & (all)         & 72  & \texttt{\#print axioms} certificate for every theorem \\
\bottomrule
\end{tabular}
}
\end{center}

\subsection{Axiom Budget}

Six standard results from finite-group harmonic analysis and matrix
linear algebra are axiomatized rather than derived here; most have no Mathlib
counterpart in the form we use. Axiomatizing them is a deliberate trust
boundary rather than a gap: each is a named, citable classical fact, stated in
a faithful form that cannot be satisfied trivially (\texttt{matrix\_qr}, for
instance, requires a genuinely upper-triangular factor, so the degenerate
$Q = I$, $R = M$ reading is excluded), and confining the trusted statements to
these classical facts keeps every $\starG$-specific reduction inside the
machine-checked kernel; discharging the Ky~Fan lemma and QR from Mathlib
primitives is mechanical but substantial formalization work, orthogonal to the
claims this paper certifies. Three former axioms are now proved outright: the
group-Fourier convolution theorem (\texttt{fourier\_multiplicative}), directly from
the homomorphism property by reindexing the group sum; the real symmetric spectral
theorem, from Mathlib's \texttt{Matrix.IsHermitian.spectral\_theorem} (giving the
$\starG$ symmetric eigendecomposition); and the least-squares minimizer, from the
normal equations over $\bbR$ (giving $\starG$ least-squares). Discharging these
dropped the custom-axiom count from nine to six. Every other
statement is derived from first principles.  The classical matrix
Eckart--Young theorem itself is \emph{not} among the axioms: it is proved in
\texttt{EckartYoung.lean} (\texttt{matrix\_eckart\_young}) from the two bridge
lemmas listed. The matrix-mimetic bridge lemma
(\texttt{fourierBlock\_starTranspose}: the Fourier transform sends the $\starG$
transpose to the transpose of each irrep block), on which the QR, spectral, and
least-squares $\starG$-orthogonality reductions rest, is likewise proved, not
assumed, and depends only on Lean's core axioms.

\begin{center}
\resizebox{\textwidth}{!}{
\begin{tabular}{@{}lll@{}}
\toprule
\textbf{Axiom} & \textbf{Content} & \textbf{Reference} \\
\midrule
\texttt{parseval\_group} & Plancherel identity, finite groups & \cite{serre1977linear}, \S2.4 \\
\texttt{fourier\_surjective} & Fourier inversion: every block field is attained & \cite{peterweil1927} \\
\texttt{fourier\_injective} & Fourier injectivity (used by $\starG$-QR/spectral) & \cite{peterweil1927} \\
\texttt{rank\_le\_proj\_capture} & rank-$\le k$ matrix is captured by a rank-$\le k$ projection & standard \\
\texttt{ky\_fan\_with\_witness} & Ky Fan (1949): optimal rank-$k$ projection for symmetric $M$ & standard \\
\texttt{matrix\_qr} & Matrix QR factorization ($Q$ orthonormal, $R$ triangular) & standard \\
\bottomrule
\end{tabular}
}
\end{center}

The complex-typed module \texttt{ComplexCore.lean}, which exists to build
concrete cyclic-group (DFT) witness systems over $\bbC$ that the real
development cannot express, declares three \emph{complex-local} analogues of
the Fourier axioms (\texttt{parseval\_groupC},
\texttt{matrix\_best\_rank\_k\_approxC}, \texttt{fourier\_surjectiveC}). They
are used only by that module's complex Eckart--Young corollaries; none of the
real-development theorems above depends on them, as the audit certificate
shows. Counting them, the repository contains nine \texttt{axiom}
declarations in total: six real, three complex-local.

In addition, the octahedral selection-rule theorems are decided by exhaustive
native computation (\texttt{native\_decide}) and therefore depend on Lean's
compiled-evaluation axiom \texttt{Lean.ofReduceBool}: they trust Lean's
compiler as an oracle for finite rational arithmetic over the $24$ group
elements.  The hypotheses of the optimality theorem are carried by an
\texttt{IrrepSystem} structure requiring completeness
($\sum_\rho d_\rho^2 = |G|$), irreducibility of each member (no proper
invariant subspace), \emph{absolute} irreducibility of each member (Schur
condition: every matrix commuting with the representation is scalar), and
pairwise inequivalence (no invertible intertwiner).  The Schur field is what
makes the Fourier axioms sound over $\mathbb{R}$: without it, a single
real-irreducible representation of complex type can satisfy
$\sum_\rho d_\rho^2 = |G|$ (for example, the two-dimensional rotation
representation of $C_4$, with $d^2 = 4 = |C_4|$) while spanning only part of
the group algebra, which would falsify Parseval and Fourier injectivity.
With absolute irreducibility, completeness, and inequivalence, the
Artin--Wedderburn decomposition forces the system to exhaust
$\mathbb{R}[G]$, so the axioms hold for every admissible instantiation.

\subsection{Key Proof Techniques}

\paragraph{Associativity of $\starG$ (Proposition~4.2(i)).}
Rather than fragile nested sum-exchange calls (\texttt{Finset.sum\_comm}),
we define an explicit \texttt{Equiv} on the 4-tuple product type
$\mathrm{Fin}\,p \times (G \times (\mathrm{Fin}\,m \times G))$ that
simultaneously permutes components and applies the bijection
$b \mapsto a^{-1}b$.  A single call to \texttt{Fintype.sum\_equiv} then
completes the proof, with the group-element arithmetic handled by the
\texttt{group} tactic.

\paragraph{Kronecker product of irreps (Theorem~2(iii)).}
The tensor-product representation $\rho_1 \otimes \rho_2$ is defined
entry-wise via \texttt{finProdFinEquiv.symm}, mapping
$\mathrm{Fin}(d_1 d_2)$ indices to pairs $\mathrm{Fin}\,d_1 \times
\mathrm{Fin}\,d_2$.  A \texttt{sum\_split} helper converts sums over
$\mathrm{Fin}(d_1 d_2)$ into double sums, after which the
\texttt{is\_hom} and \texttt{unitary} proofs factor naturally into
products of single sums via $\rho_i.\text{is\_hom}$ and
$\rho_i.\text{unitary}$.

\paragraph{Fourier power invariance (Corollary~7.2(ii)).}
The \texttt{fourierBlock\_leftAction} lemma shows that the group action
multiplies each Fourier block by $(I_\ell \otimes \rho(g))$.  An
\texttt{orthogonal\_preserves\_sum\_sq} lemma proves
$\sum_s (\sum_{s'} R_{s,s'} v_{s'})^2 = \sum_s v_s^2$ when
$R^\top R = I$, by expanding squares, exchanging sums, and applying
orthogonality.

\paragraph{Eckart--Young for $\starG$ (Theorem~1).}
The optimal rank-$k$ approximation is defined in the Fourier domain via
\texttt{fourier\_surjective}: its Fourier block at each irrep $\rho$ is, by
construction, the best rank-$k$ matrix approximation of
$\hat{\calA}(:,:,\rho)$, supplied by \texttt{matrix\_eckart\_young}.  That
matrix-level theorem is \emph{proved} in \texttt{EckartYoung.lean}: a
Pythagorean split over orthogonal projections reduces best rank-$\le k$
approximation to maximizing the captured trace over rank-$\le k$ projections
(\texttt{rank\_le\_proj\_capture}), and the Ky Fan lemma
(\texttt{ky\_fan\_with\_witness}) identifies the top-$k$ eigenprojection as
the maximizer.  The global bound follows by applying \texttt{parseval\_group}
to decompose the Frobenius error into per-irrep terms, multiplying each
per-irrep inequality by the positive Plancherel weight $d_\rho / |G|$, and
summing via \texttt{Finset.sum\_le\_sum}.

\paragraph{Wigner--Eckart selection rules (\S2.2).}
The octahedral rotation group is realised concretely as $O \cong S_4$
(\texttt{Equiv.Perm (Fin 4)}) through its action on the four body diagonals
of the cube.  The five irreducible characters are constructed rather than
hardcoded: $A_1$ trivial, $A_2$ the sign character, $T_2$ the standard
representation ($\chi(\sigma) = \mathrm{fix}(\sigma) - 1$), and $T_1 = T_2
\otimes \mathrm{sign}$ come from explicit representations, while $E$ is
recovered as a class function from the regular character and certified
irreducible by $\langle \chi_E, \chi_E \rangle = 1$, its dimension
$\chi_E(1) = 2$ entering the verified count \texttt{dim\_sum} below.
\texttt{char\_orthonormal} proves
$\langle \chi_\rho, \chi_\sigma \rangle = \delta_{\rho\sigma}$, certifying
that the five are pairwise-inequivalent irreducibles, and \texttt{dim\_sum}
checks $\sum_\rho d_\rho^2 = 24$.  Every selection rule is then
\emph{derived} as a character inner product, never asserted by definition:
a scalar ($A_1$) operator connects an irrep only to itself
(\texttt{scalar\_selection}); $T_1 \otimes T_1$ contains $A_1$ exactly once
(\texttt{vector\_couples\_scalar}); and the symmetric square decomposes as
$\mathrm{Sym}^2(T_1) = A_1 \oplus E \oplus T_2$ with zero $T_1$ multiplicity
(\texttt{polarizability\_no\_T1}, \texttt{polarizability\_decomposition}).
The finite rational arithmetic over the $24$ group elements is discharged by
\texttt{native\_decide}; the build fails if any multiplicity is wrong (for
example, if $T_1$ and $T_2$ are interchanged).

\subsection{Verification Status}

The formalization achieves:
\begin{itemize}
\item Zero \texttt{sorry} (unresolved proof obligations) across all nine
  modules.
\item Six declared axioms in the real development, each a standard textbook
  result of finite-group
  harmonic analysis or matrix linear algebra with no Mathlib counterpart in the
  form we use (three former axioms are now proved outright: the group-Fourier
  convolution theorem, the real symmetric spectral theorem from Mathlib's
  \texttt{Matrix.IsHermitian.spectral\_theorem}, and the least-squares minimizer
  from the normal equations); the complex witness module adds the three
  complex-local analogues tabulated above, nine \texttt{axiom} declarations in
  the repository in total.  Their
  exact usage, as certified by
  \texttt{\#print axioms} in \texttt{StarG/Audit.lean}: the Eckart--Young
  chain (\texttt{eckart\_young\_starG}, \texttt{eckart\_young\_error};
  Theorem~\ref{thm:ey_full}) uses \texttt{parseval\_group},
  \texttt{fourier\_surjective}, \texttt{rank\_le\_proj\_capture}, and
  \texttt{ky\_fan\_with\_witness}; the product-group Fourier factorization
  (Theorem~\ref{thm:prod_full}) now rests on the core axioms, its
  transform-multiplicativity being a theorem; the
  $\starG$-QR theorem uses \texttt{matrix\_qr}, while the $\starG$ symmetric
  eigendecomposition and $\starG$ least-squares minimizer carry no matrix axiom,
  being proved from Mathlib's \texttt{Matrix.IsHermitian.spectral\_theorem} and
  from the normal equations over $\bbR$ respectively; their
  $\starG$-orthogonality reductions rest on \texttt{fourier\_injective}
  together with the axiom-free transpose bridge lemma
  (\texttt{fourierBlock\_starTranspose}).  The octahedral selection-rule theorems
  additionally depend on \texttt{Lean.ofReduceBool} through
  \texttt{native\_decide}, as disclosed above.
\item Coverage of Theorems~\ref{thm:ey_full} and~\ref{thm:prod_full}, the
  algebraic identities of Proposition~\ref{prop:T_props} and the
  $\starG$-product properties, equivariance and Frobenius/Fourier-power
  invariance, the Kronecker product construction for product-group irreps,
  and the Wigner--Eckart selection rules from \S2.2.  Not formalized: the
  band-limited compact-group reduction (Theorem~4 of the main text) and its
  ENN corollary, whose mathematical content is the classical Peter--Weyl
  theorem.
\item All algebraic theorems (associativity, identity, distributivity,
  transpose, equivariance, Frobenius and per-irrep Fourier-power invariance)
  depend solely on Lean's three core axioms
  (\texttt{propext}, \texttt{Classical.choice}, \texttt{Quot.sound}); they
  introduce no project-level axioms.  (The formal \texttt{StarGSVD} structure
  records the unitarity of the factors but does not yet impose
  f-diagonality on $\calS$, so the accompanying existence lemma is a
  packaging statement rather than a verified factorization theorem; the
  optimality theorem above does not rely on it.)
\item A concrete, axiom-free \texttt{IrrepSystem} for the cyclic group $C_2$ (its
  trivial and sign representations) is constructed, and the full optimality suite
  is instantiated at it (\texttt{eckart\_young\_c2}, \texttt{starG\_qr\_c2},
  \texttt{starG\_spectral\_c2}, \texttt{starG\_lstsq\_c2}), each adding no custom
  axiom beyond its parent theorem. The theorems are therefore demonstrably
  non-vacuous rather than conditional over a never-inhabited hypothesis. The
  cyclic groups $\bbZ_n$ with $n\ge3$, whose complex irreducibles the real
  structure cannot express, are witnessed separately by the additive complex
  port described below; a concrete octahedral witness is the one remaining step.
\end{itemize}

\paragraph{Scope of the formalization: real versus complex.}
The formalization is stated over $\mathbb{R}$.  This faithfully covers groups
all of whose irreducible representations are of real type, such as the
octahedral group used in our Wigner--Eckart analysis (where
$\sum_\rho d_\rho^2 = 24$ holds over $\mathbb{R}$).  For cyclic groups the
complex DFT used in our experiments pairs conjugate one-dimensional complex
irreps into two-dimensional real blocks, so $\sum_\rho d_\rho^2 = |G|$ fails
over $\mathbb{R}$ and the cyclic groups $\bbZ_n$ with $n\ge3$ are not witnessed by
a real \texttt{IrrepSystem}.  The one cyclic group whose irreducible
representations are real, $C_2$, is witnessed concretely (above); for the octahedral
group $O\cong S_4$, four of its five irreducibles are formalized as genuine
irreducibles (the two one-dimensional representations and the three-dimensional
standard representation and its sign twist, all realized with rational matrix
entries and proven absolutely irreducible), so only the two-dimensional
representation, which requires irrational basis entries, and the assembled system
remain before the optimality suite instantiates at the flagship group.  To cover
the cyclic instantiations we carried out an additive complex port (module
\texttt{ComplexCore}): a complex-typed $\star_G$ layer whose unitarity is stated
with the conjugate transpose and whose Schur field is ordinary (over $\mathbb{C}$,
absolute irreducibility coincides with irreducibility).  In it we construct a
concrete complex \texttt{IrrepSystem} for every cyclic group directly from the DFT
characters $\rho_k(j)=\omega^{jk}$, instantiated at $\bbZ_3$, and prove the
group-Fourier convolution theorem and the resulting per-character
block-diagonalization; these results are axiom-free beyond Lean's core.  A complex
per-block Eckart--Young corollary is also instantiated, resting on three
complex-typed analogues of the classical Plancherel, matrix Eckart--Young, and
Fourier-surjectivity facts (declared as complex-local axioms, disjoint from and
leaving unchanged the six real axioms above).  This inhabits the cyclic
$\bbZ_n$ hypotheses that $\mathbb{R}$ cannot; the mathematics of the paper is
unaffected, and the remaining caveat concerns only the octahedral instantiation of
the real suite.

\noindent
To our knowledge, this is the first machine-checked derivation of an
Eckart--Young-type optimality theorem for symmetry-preserving tensor
approximation, conditional on the classical linear-algebra facts declared as
axioms above.

\end{document}
